\documentclass{article}
\usepackage[utf8]{inputenc}
% La plantilla original cargaba babel dos veces, con opciones distintas, lo que
% aborta la compilacion por choque de opciones. Una sola carga con las opciones
% unidas.
\usepackage[spanish, es-lcroman]{babel}
\usepackage{color}
\usepackage{ifthen}
\usepackage{xcolor}
\usepackage{titling}
\usepackage{fancyhdr}
\usepackage{geometry}
\usepackage{titlesec}
\usepackage{remreset}
\usepackage{titletoc}
\usepackage{graphicx}
\usepackage{tikz}
\usepackage{tikz}
\usetikzlibrary{arrows.meta,positioning}
\usetikzlibrary{arrows.meta,shapes.geometric}
% Gráficas de los bancos de prueba (carpeta graficas), tomadas del anexo CRISP-DM con la
% leyenda adaptada a la memoria.
\usepackage{pgfplots}
\usepackage{pgfplotstable}
\pgfplotsset{compat=1.18}
\pgfkeys{/pgf/number format/.cd, use comma, 1000 sep={.}}
\newcommand{\graficasdir}{graficas}
\usepackage{enumitem}
\usepackage{float}
\usepackage{setspace}
\usepackage{amsmath}
\usepackage{booktabs}
\usepackage{pifont}
\usepackage{colortbl}
\newcommand{\cumple}{\textcolor{green!55!black}{\ding{51}}}
\newcommand{\nocumple}{\textcolor{red!70!black}{\ding{55}}}
\newcommand{\parcial}{\textcolor{orange!85!black}{\ding{108}}}
\usepackage{longtable}
\usepackage{tabularx}
\usepackage{array}
\usepackage{multirow}
\usepackage{listings}
% Seudocodigo en espanol para las secciones de algoritmos (formato del curso de algoritmos).
\usepackage[spanish,ruled,linesnumbered]{algorithm2e}
\SetAlgorithmName{Algoritmo}{algoritmo}{Lista de algoritmos}
\SetKwInput{KwEntrada}{Entrada}
\SetKwInput{KwSalida}{Salida}
% natbib y no biblatex, para compilar con bibtex sin depender de biber.
\usepackage[numbers,sort&compress]{natbib}
\usepackage[hidelinks]{hyperref}

\graphicspath{{img/}{img/arquitectura/}{img/contexto/}{img/diseno/}{img/desarrollo/}{img/resultados/}}

\setstretch{1.15}

\pagenumbering{roman}
\pagestyle{fancy}

\fancyhead[R]{Memoria Trabajo de Grado}
\fancyhead[LO]{\ifthenelse{\isodd{\value{page}}}{\textbf{Pontificia Universidad Javeriana}}{\textbf{Ingeniería de Sistemas} \\ \codigoproyecto}}
\fancyfoot[R]{Página \thepage }
\fancyfoot[C]{}
\renewcommand{\footrulewidth}{1pt}

\definecolor{azul_oscuro}{HTML}{003770}

\title{Tukit: Sistema de Recomendación Multiobjetivo para la Construcción Asistida de Kits Escolares}
\author{
\begin{tabular}{c}
    Maria Paula Rodríguez Ruiz \\
    Juan Enrique Rozo Tarache \\
    Gabriel Aníbal Riaño Beltrán \\
    Daniel Felipe Castro Moreno
\end{tabular}
}

\newcommand{\codigoproyecto}{ISIS2610-G19}
\newcommand{\anioproyecto}{2026}
\newcommand{\director}{Ing. Julio Ernesto Carreño Vargas}

\geometry{textwidth=15.5cm,left=3cm}

\titlecontents{section}[0em]{\bfseries\large}{\thecontentslabel. }{}{\titlerule*[1pc]{.}\contentspage}

\titleformat{\section}
  {\centering\normalfont\Large}
  {\thesection}{1em}{}

\titleformat{\subsection}
  {\bfseries\large}{\hspace{2em}\thesubsection}{1em}{}

\titleformat{\subsubsection}
  {\normalfont\normalsize\bfseries}
  {\hspace{2em}\thesubsubsection}{1em}{}

\makeatletter
\@removefromreset{subsection}{section}
\makeatother

\renewcommand\thesection{\Roman{section}}
\renewcommand\thesubsection{\arabic{subsection}}
% Las referencias a subsecciones llevan el capítulo («VII.6»); los títulos y el índice
% conservan la numeración de la plantilla.
\makeatletter
\renewcommand\p@subsection{\thesection.}
\renewcommand\p@subsubsection{\thesection.}
\makeatother

\begin{document}

    \renewcommand{\contentsname}{TABLA DE CONTENIDO}
    \renewcommand{\listfigurename}{LISTA DE FIGURAS}
    \renewcommand{\listtablename}{LISTA DE TABLAS}
    \def\tablename{Tabla}
    
    \begin{titlepage}
    \begin{center}
        \vspace*{1cm}

        \Huge
        \textbf{\codigoproyecto}

        \vspace{0.5cm}
        \LARGE
        
        \Huge
        \thetitle
        
        \vspace{0.5cm}
        \LARGE
        
        \vspace{3.5cm}
        
        \theauthor
        
        \vfill

        \large
        PONTIFICIA UNIVERSIDAD JAVERIANA \\
        \vspace{0.5cm}
        FACULTAD DE INGENIERÍA \\
        \vspace{0.5cm}
        PROGRAMA DE INGENIERÍA DE SISTEMAS \\
        \vspace{0.5cm}
        BOGOTÁ, D.C. \\
        \vspace{0.5cm}
        \anioproyecto
        
    \end{center}
\end{titlepage}



    \setcounter{page}{1}

    \begin{center}

        \large
        \codigoproyecto \\

        \vspace{0.5cm}
        \LARGE
        
        \large
        \thetitle
        
        \vspace{0.5cm}
        \LARGE


        \vspace{0.5cm}

        \large
        \textbf{Autores:}\\
        \vspace{0.2cm}
        \theauthor \\
        
        \vspace{0.5cm}
        \LARGE


        \large
        \vspace{0.5cm}
         MEMORIA DE PROYECTO DE GRADO REALIZADO PARA CUMPLIR UNO DE LOS REQUISITOS PARA EL TÍTULO EN INGENIERÍA DE SISTEMAS \\

        \vspace{0.5cm}
        \LARGE


        \vspace{0.5cm}

        \large
        \textbf{Director}\\
        \vspace{0.2cm}
        \director\\

        \vspace{0.5cm}
        \LARGE


        \vspace{0.5cm}

        \large
        \textbf{Jurados del Trabajo Final de Grado}\\
        \vspace{0.2cm}
        {[Por confirmar]} \\
        {[Por confirmar]} \\

        \vspace{0.5cm}
        \LARGE


        \vspace{1cm}

         \large
        PONTIFICIA UNIVERSIDAD JAVERIANA \\
        \vspace{0.5cm}
        FACULTAD DE INGENIERÍA \\
        \vspace{0.5cm}
        PROGRAMA DE INGENIERÍA DE SISTEMAS \\
        \vspace{0.5cm}
        BOGOTÁ, D.C. \\
        \vspace{0.5cm}
        \anioproyecto
        
    \end{center}
    \begin{center}

        \large
        \textbf{PONTIFICIA UNIVERSIDAD JAVERIANA} \\
        \textbf{FACULTAD DE INGENIERÍA} \\
        \textbf{PROGRAMA DE INGENIERÍA DE SISTEMAS} \\

        
        \vspace{0.5cm}
        \LARGE


        \vspace{8cm}

        \large
        \textbf{Rector de la Pontificia Universidad Javeriana}\\
        \vspace{0.5cm}
        P. Luis Fernando Múnera Congote, S.J.\\
        \vspace{0.5cm}
        
        \textbf{Decano de la Facultad de Ingeniería}\\
        \vspace{0.5cm}
        Ing. Diego Alejandro Patiño Guevara, Ph.D.\\
        \vspace{0.5cm}
        
        \textbf{Director del Departamento de Ingeniería de Sistemas}\\
        \vspace{0.5cm}
        Ing. César Julio Bustacara Medina, Ph.D. \\
        \vspace{0.5cm}

        \textbf{Director de Carrera de Ingeniería de Sistemas}\\
        \vspace{0.5cm}
        Ing. Andrea del Pilar Rueda Olarte \\
        
    \end{center}
    \noindent \textbf{Artículo 23 de la Resolución No. 1 de Junio de 1946} \\

\vspace{0.25cm}

\noindent \textit{“La Universidad no se hace responsable de los conceptos emitidos por sus alumnos en sus proyectos de grado. Sólo velará porque no se publique nada contrario al dogma y la moral católica y porque no contengan ataques o polémicas puramente personales. Antes bien, que se vean en ellos el anhelo de buscar la verdad y la Justicia”}
    % Sin agradecimientos, por decisión de los autores.
    % \begin{center}
    \textbf{AGRADECIMIENTOS}
\end{center}

\vspace{0.25cm}

%Write a message if you feel gratitude for someone who has supported the development of the project. Your family, your partner, your friends, your principal, teachers, etc

% Redacte un mensaje de agradecimiento dirigido a quienes hayan apoyado el desarrollo del proyecto. Puede incluir a su familia, pareja, ami-gos, directores, profesores, o cualquier otra persona que considere fundamental en este proceso.
    
    \tableofcontents
    \newpage
    \listoffigures
    \newpage
    \listoftables

    \begin{center}
    \textbf{ABSTRACT}
\end{center}

\vspace{0.25cm}

Colombian households often receive extensive and ambiguous school supply lists, forcing them to visit multiple stores to fulfill them—frequently without finding a single, comprehensive kit that meets their specific needs. Tukit processes these lists—whether provided as images, PDFs, or free-text—by comparing each item against a canonical catalog, utilizing both lexical and semantic similarity. It then assembles optimized kits using a multi-objective approach based on the NSGA-II algorithm, balancing factors such as cost, brand fidelity, product availability, and thematic preferences. The project followed a hybrid, encompassing methodology that integrates SCRUM, Design Thinking, RUP, and CRISP-DM, in alignment with IEEE standards 830, 1016, and 1058.
\begin{center}
    {\textbf{RESUMEN}}
\end{center}

\vspace{0.25cm}

Los hogares colombianos reciben listas escolares extensas y ambiguas, y deben recorrer
varios comercios para completarlas sin encontrar un kit que se ajuste a sus necesidades.
Tukit lee la lista en imagen, PDF o texto libre, compara cada renglón con un catálogo
canónico combinando el parecido léxico y el semántico. Posteriormente, arma kits con una optimización
multiobjetivo basada en NSGA-II que equilibra costo, fidelidad de marca, disponibilidad y
preferencias temáticas. El trabajo siguió una metodología híbrida contenedora que combina
SCRUM, Design Thinking, RUP y CRISP-DM con los estándares IEEE 830, 1016 y 1058.
    
    \pagenumbering{arabic}
\section{INTRODUCCIÓN}\label{cap:introduccion}

Al comienzo de cada año escolar, las familias colombianas reciben del colegio una lista de
útiles que deben conseguir completa y en pocos días. El gasto es alto, porque el 35\,\% de
las familias destina entre \$300.000 y \$500.000 por niño \citep{noauthor_resultados_2025},
y la presión de tiempo también, porque el 70\,\% no planea la compra con anticipación
\citep{noauthor_resultados_2024}. A eso se suma el esfuerzo de interpretar cada renglón,
buscarlo en varias tiendas y comparar precios y marcas, una carga mental que lleva a decidir
peor \citep{sweller_cognitive_2011}. Cuando la compra se retrasa o queda incompleta, quien
lo paga es el estudiante, que empieza el año sin los materiales que el colegio pidió para
sus clases.\par

En Bogotá hay más de 1,17 millones de estudiantes matriculados
\citep{noauthor_estadisticas_nodate}, y este proyecto estudió el problema de cerca en el
Colegio Agustiniano Ciudad Salitre. Una encuesta a sus acudientes, con 60 respuestas
válidas, mostró que el 65\,\% dedica cuatro horas o más a la compra, que el 68,3\,\% compra
en tres o más lugares y que la queja más frecuente es no encontrar el producto o que esté
agotado (43,3\,\%). Las entrevistas con directivos, docentes y acudientes confirmaron esa
fricción.

Las alternativas que existen hoy no resuelven ese recorrido completo. Panamericana y Éxito
publican catálogos por categoría \citep{panamericana_cotizador_2025,exito_listas_2025},
Primavera permite armar un carrito a mano sobre su propio catálogo
\citep{primavera_armatulista_2026} y MercadoLibre vende kits fijos que arma cada vendedor
\citep{mercadolibre_kits_2025}. En todos los casos es el acudiente quien debe traducir la
lista del colegio en productos. Plataformas como TeacherLists, en Estados Unidos, convierten la lista del salón en un carrito de compra, pero ninguna
lee listas en papel o en foto ni está disponible en Colombia.

Frente a este panorama, el presente proyecto describe a Tukit como una plataforma que le permite al acudiente conseguir la lista de útiles completa.
Funciona como un sistema de recomendación multiobjetivo con un modelo B2B2C, en el que el
colegio aporta la lista oficial y la familia hace la compra. El acudiente elige el colegio y
el grado de cada hijo, o sube su propia lista en imagen, PDF o texto libre, e indica su
presupuesto y sus preferencias. A cambio recibe hasta cuatro kits completos (Ahorrador,
Premium, Fan y A tu Medida) y puede cambiar
cualquier producto antes de pagar en línea. Para llegar a esos kits, el sistema recorre tres etapas.\par

\begin{itemize}
  \item \textbf{Percepción.} Un modelo multimodal de inteligencia artificial lee la lista y
        la convierte en renglones.
  \item \textbf{Normalización.} Cada renglón se identifica como un útil, un texto escolar,
        una licencia digital u otro tipo, y se separa la cantidad que pide el colegio del
        empaque en que se vende el producto. Luego, cada útil se busca en un catálogo de
        conceptos canónicos con una búsqueda híbrida. Un canal léxico corrige la ortografía
        y exige que lo pedido esté nombrado, y un canal semántico compara significados para
        reconocer sinónimos. Si el sistema no tiene certeza suficiente, el renglón pasa a
        revisión humana.
  \item \textbf{Optimización.} Un algoritmo calcula el frente de Pareto de los kits
        posibles, es decir, aquellos en los que no se puede bajar el costo sin ceder en la
        marca o en las preferencias. Si las combinaciones son pocas (hasta 4.096), las revisa
        todas; si son más, usa NSGA-II \citep{deb_fast_2002}. De ese frente sale el kit de
        cada perfil.
\end{itemize}

Esta estrategia toma de un esquema RAG la indexación del catálogo y la recuperación de los
productos, pero en lugar de que un modelo de lenguaje redacte la respuesta, los kits los
construye el optimizador y el equipo de Tukit los revisa antes de publicarlos
(\emph{human-in-the-loop}). Todo corre sobre una arquitectura de seis servicios con cinco
bases de datos, elegida entre tres alternativas con el método ATAM
\citep{bass_software_2021}.\par

El trabajo siguió una metodología híbrida contenedora organizada en cinco fases, una por cada
objetivo específico. SCRUM \citep{schwaber_scrum_2020} organizó la construcción en ciclos de
dos semanas, Design Thinking \citep{brown_design_2008} guió la propuesta de valor y el modelo
de negocio, RUP \citep{kruchten_rup_2003} aportó la gestión de riesgos y la trazabilidad, y
CRISP-DM \citep{chapman_crisp-dm_2000} fue el marco de los componentes que dependen de datos.
Los estándares IEEE 830 \citep{ieee_830_1998}, 1016 \citep{ieee_1016_2009} y 1058
\citep{ieee_1058_1998} definieron la forma de los requerimientos, el diseño y el plan del
proyecto. Aunque el piloto se hizo en un colegio de Bogotá, el modelo puede llevarse a otros
colegios y ciudades, porque sumar un colegio consiste en cargar y revisar su lista, sin
cambiar el sistema.\par
    \setcounter{subsection}{0}
\section{DESCRIPCIÓN GENERAL}

El presente capítulo expone una visión de conjunto del proyecto. En primer término se aborda
la problemática que lo origina, relacionada con el tiempo, el costo y el esfuerzo que supone
para las familias transformar la lista de útiles entregada por el colegio en una compra
completa, en ausencia de herramientas que asistan dicho proceso. Sobre esa base se plantea
la pregunta de investigación y se describe la solución propuesta, Tukit, un sistema de
recomendación multiobjetivo apoyado en inteligencia artificial que interpreta la lista, la
relaciona con un catálogo de productos y genera kits ajustados al presupuesto y a las
preferencias de cada hogar. Posteriormente se exponen la justificación de la solución, el
objetivo general y los objetivos específicos, y finalmente el marco metodológico, los
entregables y los estándares que orientaron su desarrollo, elementos que permiten
comprender el diseño y el alcance del sistema presentado en los capítulos siguientes.\par

\subsection{Oportunidad, Problema}

\subsubsection{Contexto del Problema}

El comercio electrónico en Colombia ha mantenido un crecimiento sostenido. En 2025, las
ventas en línea alcanzaron COP 105 billones, un 5,9\,\% más que el año anterior, y la
categoría de libros y papelería representó el 0,81\,\% de las transacciones
\citep{noauthor_informe_nodate-1}. No obstante, ese crecimiento no ha transformado la manera
en que se adquieren los útiles escolares, una tarea que se repite cada año y que, solo en
Bogotá, involucra a más de 1.171.390 estudiantes matriculados
\citep{noauthor_estadisticas_nodate}. De acuerdo con las encuestas de temporada escolar de
FENALCO, el 35\,\% de las familias destina entre \$300.000 y \$500.000 por niño
\citep{noauthor_resultados_2025}, y el 70\,\% no planea la compra con anticipación. Apenas
el 6\,\% la realiza con meses de margen, mientras que el resto compra a medida que surgen
las necesidades o pocas semanas o días antes del inicio de clases
\citep{noauthor_resultados_2024}. A esta presión de tiempo se suma la de los precios, pues
el 91\,\% de los padres percibe un alza, y el precio continúa siendo el principal criterio
de elección (33\,\%), seguido de la calidad (23\,\%) y, de forma creciente, del gusto de los niños (22\,\%) \citep{noauthor_resultados_2025}.\par

La dificultad implica tiempo, dinero, y un esfuerzo
mental considerable. La teoría de la carga cognitiva establece que, cuando la complejidad de
una tarea supera la capacidad de la memoria de trabajo, la calidad de las decisiones
disminuye \citep{sweller_cognitive_2011}. Esta situación se presenta en la compra escolar,
en la que el acudiente debe buscar cada artículo por separado en plataformas genéricas,
comparar manualmente precio, disponibilidad y especificación, y conformar un carrito que
rara vez cubre la lista completa. Iyengar y Lepper demostraron que el exceso de opciones
inhibe la compra, ya que al pasar de 6 a 24 alternativas la tasa de compra se redujo del
30\,\% al 3\,\%, aun cuando la exhibición amplia atrajo a más personas (60\,\% frente a
40\,\%) \citep{iyengar_when_2000}. La abundancia de alternativas genera una sensación de desorden que termina por paralizar la decisión, y en una lista escolar este efecto se repite
con cada renglón. En el comercio electrónico en general, cerca del 70\,\% de los carritos se
abandona antes del pago \citep{noauthor_50_nodate}.\par

Los actores que lideran este segmento en Colombia trasladaron su modelo físico al canal
digital sin modificarlo, de modo que publican un catálogo y delegan en el cliente la
búsqueda, la comparación y la validación. Panamericana ofrece cotizaciones por categorías
genéricas, sin relación con la lista de cada colegio \citep{panamericana_cotizador_2025};
Primavera permite armar manualmente un carrito sobre su propio catálogo
\citep{primavera_armatulista_2026}; MercadoLibre comercializa kits preconfigurados por
terceros \citep{mercadolibre_kits_2025}, y Éxito ofrece kits por nivel escolar sin
distinguir entre colegios ni optimizar varios criterios a la vez
\citep{exito_listas_2025}. En el ámbito internacional, TeacherLists convierte las listas
de los salones en carritos prellenados \citep{teacherlists_retail_2026}, pero no interpreta listas impresas o fotografiadas, no
reconoce las variaciones del vocabulario local (esferos o lapiceros, marcadores o plumones)
ni distingue entre presupuestos familiares. En consecuencia, se identifica la ausencia de un actor que
transforme el modelo reactivo de compra en uno proactivo y asistido, y esta brecha
constituye la oportunidad que motiva el presente proyecto.\par

\subsubsection{Formulación del Problema}


Los canales de venta de útiles escolares son abundantes e incluyen papelerías físicas,
grandes superficies y mercados electrónicos generales. Sin embargo, no existe un mecanismo
que traduzca de manera automática el texto libre y heterogéneo de una lista escolar en un
pedido correcto y entregable. Dicha traducción exige determinar qué significa cada renglón
de la lista dentro de un vocabulario que nunca ha sido estandarizado, y establecer qué
combinación de productos equilibra el costo, la marca que prefiere el acudiente y la
disponibilidad real.\par

En consecuencia, el problema se descompone en tres dimensiones que deben resolverse en orden
y de forma integrada. La primera es la percepción, que consiste en convertir una lista
recibida en PDF, en fotografía o como texto libre en una representación estructurada que
pueda ser procesada por un programa. La segunda es la normalización, que asocia cada
renglón, con frecuencia vago, incompleto o mal escrito, con el producto real del catálogo
que le corresponde, resolviendo las ambigüedades y las distintas formas de nombrar un mismo
artículo. La tercera es la optimización, que construye kits considerando simultáneamente el
costo, las preferencias de la familia y la calidad, y presenta un número reducido de
alternativas claramente diferenciadas que facilitan la decisión.\par

La fricción de este proceso, definida como el esfuerzo innecesario en la compra de útiles, se analiza desde el enfoque de completación de tareas de
Sismeiro y Bucklin \citep{sismeiro_modeling_2004}, que describe la compra en línea como una
secuencia de tareas que el usuario debe completar de forma sucesiva para que la venta se
concrete. Bajo esta perspectiva, cada renglón que el acudiente debe buscar, comparar y
validar por su cuenta constituye una tarea adicional y, por lo tanto, un punto en el que
puede abandonar la compra. 
Si bien el problema admite aproximaciones desde el mercadeo, la
logística o la educación, una solución tecnológica resulta especialmente pertinente. Los
agentes de recomendación reducen el esfuerzo de búsqueda y mejoran la calidad de la
decisión del consumidor \citep{haubl_consumer_2000}, y el precedente de TeacherLists
muestra que el paso automático de la lista al carrito es viable
\citep{teacherlists_retail_2026}. Adicionalmente, la madurez reciente de los modelos
multimodales \citep{shen_ocr_2026}, la búsqueda vectorial \citep{karpukhin_dense_2020} y la
optimización multiobjetivo \citep{deb_fast_2002} permite abordar las tres dimensiones dentro
de un mismo sistema. A partir de lo anterior, el problema de investigación se formula de la
siguiente manera.\par

\begin{quote}
\itshape ¿Cómo construir kits de útiles escolares personalizados y alineados con los
requerimientos específicos de cada institución educativa, considerando el perfil
socioeconómico y las preferencias de la familia, de manera que se reduzca la carga cognitiva
del proceso de compra?
\end{quote}

\subsubsection{Propuesta de Solución}

Tukit es un sistema de recomendación multiobjetivo que opera como una plataforma B2B2C
(\emph{business-to-business-to-consumer}), es decir, que llega a las familias a través de los
colegios. El colegio aporta la lista oficial, y el acudiente realiza y paga la compra. La
plataforma recibe dicha lista en imagen, PDF o texto libre y la convierte en kits de compra
completos, cada uno construido a partir del equilibrio entre objetivos que compiten entre
sí, como el costo, la calidad y las preferencias de cada familia. Esta conversión se realiza
mediante un \emph{pipeline} de tres etapas. La primera es la percepción de la lista; la segunda, la normalización de cada renglón contra
un catálogo de conceptos canónicos, en la que los casos que no alcanzan la certeza suficiente
se remiten a revisión humana; y la tercera, la optimización multiobjetivo del conjunto de
compra.\par

En la plataforma intervienen tres actores, cada uno con un canal propio. El colegio, por
medio de su administrador, entrega y publica la lista oficial de cada grado. El equipo de
Tukit opera desde una aplicación de administración, en la que revisa las listas antes de que
lleguen a las familias, atiende la cola de renglones que el sistema no resolvió con certeza
y valida o corrige cada uno de ellos. Por su parte, el acudiente accede mediante una
aplicación web progresiva, en la que selecciona el colegio y el grado de cada hijo, o carga
su propia lista en el formato que prefiera, e indica su presupuesto y sus preferencias.\par

Como resultado, el acudiente recibe tres propuestas de kit para la lista completa. El kit
Ahorrador corresponde a la opción de menor costo; el Premium, a la de mejores marcas; y el
Fan, a la que incluye los personajes solicitados para cada hijo. Cuando el acudiente fija un
tope de gasto, recibe además el kit A tu Medida, que ofrece la mejor marca posible dentro de
ese monto. Cada propuesta presenta los productos, las cantidades que deben comprarse según el
empaque y el valor total. Finalmente, el acudiente puede reemplazar productos de manera
individual, elegir el kit de su preferencia y realizar el pago en línea.\par

\subsubsection{Justificación de la Solución}

El desarrollo de un sistema de recomendación multiobjetivo que construya kits escolares a
partir de la lista oficial de cada colegio representa una solución pertinente y viable, en
la medida en que permite simplificar un proceso de compra que hoy recae por completo sobre
las familias. En la actualidad, la adquisición de útiles escolares exige al acudiente
interpretar cada renglón de la lista, buscarlo en distintas tiendas y comparar por su cuenta
precios, marcas y disponibilidad, tal como se describió en las secciones anteriores. Un
sistema que asuma esa tarea y entregue la compra ya organizada reduce el tiempo y el
esfuerzo que demanda el proceso, y contribuye a que el estudiante cuente con los materiales
que su colegio solicitó al inicio del año escolar.\par

Presentar pocas alternativas ya armadas facilita la decisión, puesto que el exceso de
opciones dificulta la elección y tiende a postergar la compra \citep{iyengar_when_2000},
mientras que un sistema que le propone una
selección reducida disminuye su esfuerzo y mejora la calidad de su decisión
\citep{haubl_consumer_2000}. Por esta razón, Tukit se concibe como un asistente que entrega kits completos entre los cuales el acudiente
solo debe elegir. Esta forma de operar beneficia a cada uno de los actores. El acudiente
evita recorrer varios establecimientos y puede ajustar la alternativa elegida antes de
pagar; el colegio se asegura de que las familias compren a partir de la versión oficial de
su lista y puede orientarlas sin exigir marcas ni proveedores; y el equipo de Tukit concentra su revisión en los
renglones que el sistema no resolvió con certeza, en lugar de revisar cada lista de forma
manual. Esa revisión garantiza, a su vez, que la calidad de los kits no dependa únicamente
de la automatización. Cuando el sistema no alcanza la certeza suficiente sobre un renglón,
este se remite a una persona antes de que la lista llegue a las familias, siguiendo las
pautas que recomiendan indicar con claridad lo que el sistema desconoce y facilitar la
corrección de sus errores \citep{amershi_guidelines_2019}. 

Desde el punto de vista técnico, la solución es viable porque las técnicas en que se apoya
cuentan con respaldo empírico. Los modelos multimodales alcanzan en la lectura de documentos
un desempeño comparable al de los métodos tradicionales \citep{shen_ocr_2026}, el algoritmo
NSGA-II constituye una referencia para equilibrar varios objetivos en conflicto
\citep{deb_fast_2002}, y las mediciones realizadas sobre las listas reales del corpus
confirman que el sistema responde en tiempos adecuados para su uso, como se detalla en el
Capítulo~\ref{cap:desarrolloSolucion}. La arquitectura de microservicios permite, además,
que cada componente crezca según su propia demanda \citep{newman_building_2021}, y la
incorporación de un nuevo colegio solo requiere cargar y revisar su lista oficial, lo que
hace posible replicar la solución en otras instituciones y ciudades. En lo económico, el
sistema se construye sobre herramientas de código abierto, como PostgreSQL, RabbitMQ, Redis
y la biblioteca pymoo \citep{blank_pymoo_2020}, que evitan costos de licenciamiento, y
consume los modelos de inteligencia artificial mediante servicios de pago por uso, de modo
que el gasto crece en proporción al número de listas procesadas y no exige una inversión
inicial elevada. Como cada lista oficial se procesa una sola vez y queda disponible para
todas las familias del mismo grado, su costo se distribuye entre muchos usuarios, y el
modelo B2B2C se vuelve más eficiente a medida que se incorporan nuevos colegios, lo que hace
que la solución sea sostenible en el tiempo.\par

La solución es, además, económicamente viable, porque su modelo de negocio genera un margen
positivo desde la primera venta y requiere un volumen moderado para cubrir sus costos. Tukit
obtiene sus ingresos de una comisión sobre cada kit vendido, que en el escenario base es del
15\,\% sobre un valor promedio de \$180.000 por kit, es decir, \$27.000 por venta. A ese
ingreso se le restan los costos variables, que corresponden principalmente a la comisión de
la pasarela de pagos y al uso de los servicios de inteligencia artificial, y suman cerca de
\$7.500 por kit, lo que deja una contribución de \$19.500 por cada venta. Como los costos
fijos del primer año se mantienen bajos, alrededor de \$15.000.000 en infraestructura y
herramientas, el punto de equilibrio se alcanza con aproximadamente 770 kits vendidos, una
cifra modesta frente a los 1.500 kits que se proyectan para ese año con un número reducido
de colegios aliados en Bogotá. Con una inversión inicial cercana a \$70.000.000 y una
proyección de 8.000 y 25.000 kits para el segundo y el tercer año, el valor presente neto
del proyecto a tres años, descontado al 20\,\% anual, es de cerca de \$235 millones. Estas
cifras corresponden a una proyección construida sobre supuestos del equipo y dependen, sobre
todo, de la velocidad con que se incorporen nuevos colegios.\par

\subsection{Descripción del Proyecto}

\subsubsection{Objetivo General}

Diseñar una arquitectura de sistema de recomendación multiobjetivo que asista la
construcción de kits escolares, midiendo su efecto sobre la carga cognitiva percibida, la
usabilidad del sistema y la satisfacción del usuario con padres del sector Salitre de
Bogotá durante 2026.


\subsubsection{Objetivos Específicos}
\label{sec:objetivos}

\begin{itemize}
    \item \textbf{OE1.} Caracterizar el dominio escolar a través del análisis del negocio
    y el levantamiento de requerimientos con los actores clave, identificando las
    restricciones, necesidades y criterios de éxito del sistema, y estableciendo las
    bases conceptuales y de datos necesarias para el diseño y desarrollo del sistema de
    recomendación.
    \item \textbf{OE2.} Diseñar la arquitectura del sistema de recomendación, definiendo
    los componentes principales, las decisiones de diseño y los atributos de calidad
    mediante la evaluación ATAM \citep{bass_software_2021}, y documentando la solución en
    el Documento de Diseño de Software (SDD) que sirva como guía de referencia para las
    fases de implementación posteriores.
    \item \textbf{OE3.} Desarrollar los módulos de percepción y normalización semántica
    del sistema, modelando la extracción de información desde fuentes heterogéneas (PDF,
    imágenes y texto libre) e implementando el mecanismo de búsqueda vectorial densa para
    la correspondencia entre ítems de la lista institucional y el catálogo de productos,
    incluyendo la traducción de preferencias en lenguaje natural a restricciones
    estructuradas (NL2C).
    \item \textbf{OE4.} Implementar la generación de recomendaciones personalizadas
    mediante un enfoque multiobjetivo que logre un equilibrio técnico entre el costo, la
    fidelidad de marca y la disponibilidad, permitiendo la construcción de configuraciones
    diferenciadas que satisfagan las diversas restricciones de los perfiles de usuario
    identificados, integradas en el sistema, con una interfaz web progresiva.
    \item \textbf{OE5.} Evaluar el sistema mediante pruebas funcionales, no funcionales,
    de carga y de usabilidad con usuarios reales, verificando el correcto funcionamiento
    de cada componente del sistema y midiendo la experiencia del usuario en términos de
    satisfacción y reducción de la carga cognitiva percibida
    \citep{sweller_cognitive_2011} frente al proceso tradicional
    de compra.
\end{itemize}

\subsection{Marco metodológico, entregables, estándares y justificación}

El proyecto siguió una metodología híbrida contenedora, llamada así porque reúne varios
marcos de trabajo y le asigna a cada uno una función distinta. SCRUM \citep{schwaber_scrum_2020} organizó la
construcción en ciclos cortos de dos semanas, en los que el equipo planeaba, construía y
revisaba un avance funcional del sistema, y un tablero Kanban registró el estado de cada
tarea. Design Thinking \citep{brown_design_2008} guió el entendimiento de las familias y de los colegios, y con
él la propuesta de valor y el modelo de negocio. RUP \citep{kruchten_rup_2003} aportó la gestión de riesgos y la
trazabilidad, es decir, la posibilidad de seguir cada requerimiento hasta el componente de
la arquitectura que lo atiende y la prueba que lo verifica. CRISP-DM \citep{chapman_crisp-dm_2000,wirth_crisp-dm_2000} fue el proceso de
referencia para todos los componentes que dependen de datos: la lectura de las listas, su
correspondencia con el catálogo, la traducción de preferencias en lenguaje natural y el armado de
los kits. Para ello se extendió más allá de su primera etapa, la comprensión del negocio, como
se detalla en la Tabla~\ref{tab:crispdm-oe}. Por último, las normas IEEE e ISO
definieron la estructura de los documentos de ingeniería.

El trabajo se dividió en cinco fases, una por cada objetivo específico, de modo que el
resultado de cada fase sirviera como evidencia del cumplimiento de su objetivo. La
Tabla~\ref{tab:fases} resume cada fase con su objetivo, la etapa de CRISP-DM que le
corresponde, los métodos que se usaron y lo que produjo. La comprensión del negocio, que
es la primera etapa de CRISP-DM, se desarrolla en el Capítulo~III, donde se presentan las
entrevistas, la encuesta a los acudientes, la propuesta de valor y el modelo de negocio.

\begin{table}[H]
\centering\footnotesize
\setlength{\tabcolsep}{4pt}
\renewcommand{\arraystretch}{1.1}
\caption{Fases del proyecto, su relación con CRISP-DM y sus métodos.}
\label{tab:fases}
\begin{tabularx}{\textwidth}{@{}p{2.3cm}p{0.8cm}p{2.4cm}Xp{2.9cm}@{}}
\toprule
Fase & Obj. & Etapa de CRISP-DM & Métodos y técnicas & Resultado \\
\midrule
Caracterización del problema & OE1 & Comprensión del negocio &
Entrevistas, mapas de empatía, encuesta a acudientes, lienzo del modelo de negocio
\citep{osterwalder_business_2010}, FODA y recolección de listas &
Requerimientos, modelo de negocio y catálogo de conceptos \\
Diseño de la arquitectura & OE2 & No aplica (ingeniería de software) &
Atributos de calidad \citep{richards_fundamentals_2025,iso_25010_2011}, árbol de
utilidad, ATAM sobre tres alternativas \citep{bass_software_2021} y vistas 4+1
\citep{kruchten_41_1995} &
Arquitectura y documento de diseño \\
Percepción y normalización & OE3 & Datos, preparación, modelado y evaluación &
Corpus de oro anotado por una persona; kappa de Cohen \citep{cohen_coefficient_1960} entre
las reglas y un modelo de lenguaje; comparación de motores de lectura y de estrategias de
búsqueda; calibración del umbral con revisión humana &
Módulos de percepción y normalización \\
Optimización e integración & OE4 & Modelado, evaluación y despliegue &
Formulación multiobjetivo, NSGA-II \citep{deb_fast_2002}, búsqueda exacta, utilidad
multiatributo (MAUT), traducción de preferencias (NL2C) e integración en sprints &
Optimizador, aplicación web y aplicación de administración \\
Validación & OE5 & Evaluación &
Pruebas funcionales, no funcionales y de carga; SUS \citep{brooke_sus_1996}, NPS
\citep{reichheld_one_2003} y verbalización concurrente \citep{lazar_research_2017} &
Informe de pruebas y evaluación con usuarios \\
\bottomrule
\end{tabularx}
\end{table}

\subsubsection{Extensión de CRISP-DM a los componentes que dependen de datos}
\label{sec:crispdm-extension}

CRISP-DM \citep{chapman_crisp-dm_2000,wirth_crisp-dm_2000} define seis etapas: comprensión del
negocio, comprensión de los datos, preparación de los datos, modelado, evaluación y despliegue.
La metodología se aplicó en la Fase~1 solo a su primera etapa. A partir de la Fase~3 se extendió a
todos los componentes cuyo comportamiento depende de datos, es decir, la percepción y la
normalización del OE3 y la optimización, la utilidad multiatributo (MAUT) y la traducción de
preferencias en lenguaje natural a restricciones (NL2C) del OE4. Para cada uno se recorrieron las
seis etapas, y la Tabla~\ref{tab:crispdm-oe} resume, por objetivo, la actividad de cada etapa, la
evidencia que produjo y la sección de esta memoria donde se presenta.

\begin{table}[H]
\centering\footnotesize
\setlength{\tabcolsep}{3pt}
\renewcommand{\arraystretch}{1.1}
\caption{Extensión de CRISP-DM por objetivo: etapa, actividad, evidencia y sección.}
\label{tab:crispdm-oe}
\begin{tabularx}{\textwidth}{@{}p{0.9cm}p{2.4cm}XXp{1.9cm}@{}}
\toprule
OE & Etapa & Actividad & Evidencia & Sección \\
\midrule
OE1 & Comprensión del negocio & Entrevistas, encuesta a acudientes (N=65 recibidas, n=60
analizables) y levantamiento de listas de 17 colegios & Dossier de negocio y SRS &
Cap.~III y IV \\
\midrule
OE3 & Comprensión de los datos & Caracterización del corpus de listas (3.608 renglones) y
armado del conjunto de oro CGS-10 (diez documentos, 1.293 ítems) & Corpus y conjunto de oro &
Sec.~\ref{sec:percepcion} \\
OE3 & Preparación & Limpieza del texto, corrección ortográfica con el léxico del catálogo, núcleo
en singular y separación de atributos & Código de normalización &
Sec.~\ref{sec:normalizacion} \\
OE3 & Modelado & Cinco enfoques de lectura; cuatro estrategias de recuperación; resolución
léxica con aporte del coseno & Bancos de los prototipos &
Sec.~\ref{sec:percepcion-dev} y \ref{sec:normalizacion} \\
OE3 & Evaluación & F1 de la lectura; AUC y precisión del umbral $\tau$ con 542 ítems; revisión
humana de los renglones dudosos & Tablas de resultados &
Sec.~\ref{sec:percepcion} y \ref{sec:recuperacion} \\
OE3 & Despliegue & Servicios de inteligencia y de catálogo; cola de revisión del administrador &
Sistema construido & Sec.~\ref{sec:pipeline} \\
\midrule
OE4 & Comprensión de los datos & Encuesta de preferencias (pregunta P21, 60 respuestas) y 174
listas del corpus & Dossier de negocio y banco de listas & Sec.~\ref{sec:optimizacion-dev} \\
OE4 & Preparación & Líneas comprables, candidatos por renglón y curaduría del catálogo &
Auditorías de los kits & Sec.~\ref{sec:optimizacion-dev} \\
OE4 & Modelado & NSGA-II con cuatro objetivos, frente exacto, MAUT y NL2C &
Código y banco de calibración & Sec.~\ref{sec:optimizacion-cal} y \ref{sec:nl2c} \\
OE4 & Evaluación & Comparación con el óptimo exacto, hipervolumen y cuatro auditorías de kits &
Tablas de resultados & Sec.~\ref{sec:optimizacion} \\
OE4 & Despliegue & Aplicación web, aplicación de administración y pago & Sistema
construido & Sec.~\ref{sec:web} \\
\bottomrule
\end{tabularx}
\end{table}

\subsubsection{Ciclos de calibración con revisión humana}
\label{sec:ciclos-calibracion}

Las etapas de CRISP-DM se recorrieron en varios ciclos, y en los componentes de percepción,
normalización y optimización cada ciclo incluyó una revisión humana. Un ciclo tuvo cuatro pasos.
Primero, se ejecutó el componente sobre una muestra. Segundo, una persona revisó los resultados:
el administrador de Tukit en la cola de renglones dudosos, o el equipo, al auditar los kits renglón
por renglón. Tercero, esa revisión produjo etiquetas o correcciones, es decir, qué producto era el
correcto y qué regla había fallado. Cuarto, el equipo ajustó los datos, las reglas o las constantes
del componente y repitió la evaluación con la misma métrica. Las cuatro auditorías de los kits de la
Tabla~\ref{tab:kits} son un ejemplo: en cada una, una persona clasificó cada renglón y, a partir de
esa clasificación, se ajustaron la comparación de atributos y la curaduría. Estos ciclos no son
aprendizaje por refuerzo, porque no hay una señal de recompensa ni una política que se actualice
sola. Los parámetros los cambia el equipo después de leer los resultados, y cada cambio queda
registrado con su evidencia en la Tabla~\ref{tab:constantes}. El diseño sigue las pautas de
interacción humano-IA de Amershi et al. \citep{amershi_guidelines_2019}, entre ellas mostrar
la incertidumbre del sistema y permitir la corrección por parte de la persona. El Capítulo~VI
describe cada ciclo junto a su componente. La última etapa, el despliegue, corresponde a la
integración en el sistema construido, y su operación con colegios y familias se presenta como
trabajo futuro en el Capítulo~VIII.

Los entregables del proyecto, el objetivo especifico (OE), los estándares que guiaron su elaboración y la razón de cada
uno se resumen en la Tabla~\ref{tab:entregables} (Anexo~\ref{anx:tablas}).

% [movida al Anexo: tab:entregables]

    \setcounter{subsection}{0}
\section{CONTEXTO DEL PROYECTO}

La compra de útiles escolares parece, a primera vista, una tarea sencilla de comercio. Sin
embargo, automatizarla exige resolver problemas que pertenecen a campos distintos, como la
lectura de documentos sin un formato común, la interpretación de un vocabulario que cada
colegio escribe a su manera, la elección entre una gran cantidad de combinaciones de
productos y la presentación de esas opciones de forma que no abrumen a quien decide. Este
capítulo reúne el contexto necesario para comprender cómo Tukit aborda cada uno de esos
retos. En primer lugar, se presentan los antecedentes conceptuales, que parten del proceso
general del sistema y describen las técnicas en que se apoya cada una de sus etapas, junto
con la arquitectura y las tecnologías que lo sostienen. Posteriormente, se analiza el
entorno en el que surge el proyecto, a partir de la caracterización del dominio escolar y
de la comparación con las soluciones que existen hoy en el mercado, con el fin de evidenciar
el vacío que la propuesta busca cubrir.\par

\subsection{Antecedentes}

En esta sección se presentan los conceptos necesarios para comprender el funcionamiento de
Tukit. Se parte del dominio, es decir, de la lista escolar y de los retos que plantea
convertirla en una compra, y a partir de ellos se introducen el tipo de sistema que se
construyó y las técnicas que resuelven cada uno de esos retos.\par

\subsubsection{La lista escolar y los retos de convertirla en una compra}

La lista de útiles es el documento con el que un colegio indica a las familias qué
materiales necesita cada estudiante para el año escolar. Aunque su propósito es el mismo en
todas las instituciones, cada colegio la elabora a su manera. Algunas listas se publican
como tablas con varias columnas, otras como textos con viñetas, y muchas llegan a las
familias como fotografías compartidas en grupos de mensajería. Además, el mismo artículo se
nombra de formas distintas, como ``cartulina en octavos de colores'' o ``cartulina octavos
surtida'', mientras que nombres casi iguales, como ``block oficio'' y ``block carta'',
corresponden a productos diferentes.\par

Convertir esa lista en una compra completa implica superar varios retos sucesivos. El
primero es leer el documento sin perder renglones, sin importar su formato ni la calidad de
la imagen. El segundo es entender qué pide cada renglón, a pesar de los errores de
ortografía, de los sinónimos y de la diferencia entre la cantidad que solicita el colegio y
el empaque en que se vende el producto. El tercero es encontrar en el catálogo el producto
que corresponde, sin ofrecer uno parecido pero equivocado, ya que la familia solo descubre
ese error cuando ya ha comprado. El cuarto es elegir, entre todas las combinaciones de
productos posibles, aquellas que mejor equilibran el costo, la calidad y las preferencias de
la familia. Finalmente, el quinto es presentar esas alternativas de manera que faciliten la
decisión en lugar de dificultarla. Los conceptos que se exponen a continuación responden a
estos retos.\par

\subsubsection{Sistemas de recomendación}

Un sistema de recomendación es un software que sugiere a una persona los elementos que con
mayor probabilidad le resultan útiles, a partir de información sobre sus necesidades, sus
preferencias o el comportamiento de otros usuarios \citep{adomavicius_toward_2005}. La
literatura distingue varios enfoques. Los sistemas basados en contenido recomiendan
elementos parecidos a los que la persona ya eligió, los de filtrado colaborativo recomiendan
lo que eligieron usuarios con gustos similares, y los basados en conocimiento razonan sobre
las necesidades explícitas del usuario y sobre las características de cada producto
\citep{burke_hybrid_2002}. Los sistemas híbridos combinan varios de estos enfoques para
aprovechar sus fortalezas.\par

En la compra escolar, el filtrado colaborativo resulta poco útil, porque una familia compra
útiles una vez al año y su lista depende del colegio y del grado, no de lo que compraron
otras familias. Por esta razón, Tukit se apoya en el enfoque basado en conocimiento, ya que
recomienda a partir de los requerimientos de la lista oficial y de las preferencias que
declara el acudiente. A diferencia de la mayoría de los sistemas, que recomiendan productos
individuales, Tukit recomienda conjuntos completos, un problema conocido como recomendación
de paquetes, en el que importa la coherencia entre los elementos que se ofrecen juntos
\citep{avny_brosh_bruce_2022}. Como cada paquete debe equilibrar varios criterios al mismo
tiempo, el sistema se clasifica como un sistema de recomendación multiobjetivo.\par

\subsubsection{Comprensión de documentos y modelos multimodales}

El primer reto, leer la lista, pertenece al campo de la comprensión de documentos, que
estudia cómo convertir una imagen o un archivo en información estructurada. El enfoque
tradicional es el reconocimiento óptico de caracteres (OCR), que identifica en la imagen las
zonas con texto y reconoce cada carácter \citep{smith_overview_2007}. Este enfoque funciona
bien con documentos impresos y nítidos, pero lee línea por línea sin comprender la
estructura de la página, de modo que no distingue si dos líneas pertenecen a la misma fila
de una tabla o si una viñeta continúa el renglón anterior.\par

Los modelos multimodales de lenguaje superan esa limitación, porque reciben al mismo tiempo
una imagen y una instrucción escrita y razonan sobre la disposición y el contenido de la
página en conjunto. Evaluaciones recientes muestran que, a partir de la imagen únicamente,
estos modelos alcanzan un desempeño comparable al de los flujos que combinan OCR y modelos
de lenguaje \citep{shen_ocr_2026}, aunque su precisión disminuye con el texto manuscrito
\citep{liu_ocrbench_2024}. Para reducir el riesgo de que el modelo produzca información que
no está en el documento, se le exige responder con un formato fijo que se verifica antes de
aceptarse. Este es el enfoque que utiliza Tukit para leer listas en fotografía, PDF o
texto.\par

\subsubsection{Representación vectorial y búsqueda híbrida}

El segundo y el tercer reto, entender cada renglón y encontrar su producto, corresponden a un
problema de búsqueda. La búsqueda léxica compara las palabras de dos textos y tolera
errores de ortografía mediante la distancia de edición, que cuenta cuántas letras deben
cambiarse para pasar de una palabra a otra. Su limitación es que no reconoce sinónimos, de
modo que ``estuche'' y ``cartuchera'' le resultan completamente distintos.\par

La búsqueda semántica resuelve ese problema mediante la representación vectorial del texto.
Un modelo de lenguaje convierte cada frase en un vector, conocido como \emph{embedding}, que
funciona como un conjunto de coordenadas en un espacio donde las frases con significado
parecido quedan cerca \citep{reimers_sentence-bert_2019}. La cercanía entre dos vectores se
mide con la similitud coseno, y a este tipo de búsqueda se le denomina recuperación densa
\citep{karpukhin_dense_2020}. Sin embargo, la búsqueda semántica presenta el riesgo
contrario al de la léxica, ya que puede considerar parecidos dos textos que tratan del mismo
tema sin pedir lo mismo.\par

Estas técnicas suelen integrarse en la generación aumentada por recuperación (RAG), en la
que un modelo de lenguaje redacta una respuesta a partir de los documentos recuperados
\citep{lewis_retrieval-augmented_2020}. En Tukit no es necesario redactar ninguna respuesta,
porque el resultado de la búsqueda es el concepto del catálogo y sus productos. Por ello, el
sistema toma de este esquema la indexación y la recuperación, y combina ambas búsquedas de
forma que la léxica determina qué conceptos son admisibles y la semántica desempata entre
ellos o reconoce sinónimos que no están registrados en el catálogo.\par

\subsubsection{Optimización multiobjetivo y algoritmo NSGA-II}

El cuarto reto consiste en elegir un producto para cada renglón de la lista. Si cada renglón
tiene tres productos posibles, una lista de \(n\) renglones admite \(3^{n}\) combinaciones,
una cantidad que crece tan rápido que una lista de cuarenta renglones supera los
\(10^{19}\) kits. Además, no existe un único kit que sea el mejor, porque el más barato no
es el de mejores marcas, y el que incluye los personajes que pidió el niño no es el más
económico. Los problemas en los que varios criterios compiten entre sí se estudian en la
optimización multiobjetivo \citep{ehrgott_multiobjective_2002}.\par

En este tipo de problemas se dice que una solución domina a otra cuando no es peor en ningún
criterio y es mejor en al menos uno. Las soluciones que ninguna otra domina forman el frente
de Pareto, es decir, el conjunto de alternativas en las que solo es posible mejorar un
criterio cediendo en otro. Cuando las combinaciones son pocas, el frente se puede obtener
revisándolas todas; cuando son muchas, se recurre a algoritmos que lo aproximan en un tiempo
razonable. El más utilizado para problemas con pocos objetivos es NSGA-II
\citep{deb_fast_2002}, un algoritmo genético que mantiene una población de soluciones, las
combina y las modifica en cada generación, y conserva las que pertenecen a los mejores
frentes y las que mantienen mayor variedad entre sí. La calidad del frente obtenido puede
medirse con el hipervolumen, que calcula el espacio de criterios que el frente domina
\citep{zitzler_multiobjective_1999}. Tukit utiliza este enfoque para construir sus kits y
toma del frente la alternativa que mejor responde a cada perfil.\par

\subsubsection{Revisión humana en sistemas de inteligencia artificial}

Ninguna de las técnicas anteriores es infalible, y en la compra escolar un error tiene un
costo directo para la familia. Por esta razón, los sistemas de inteligencia artificial que
toman decisiones con consecuencias para las personas suelen incorporar la revisión humana
en los casos de menor certeza, un esquema conocido como \emph{human-in-the-loop}. Las pautas
de interacción entre personas e inteligencia artificial recomiendan que el sistema indique
con claridad lo que no sabe y que facilite la corrección de sus errores
\citep{amershi_guidelines_2019}. En Tukit, los renglones que el sistema no resuelve con
certeza suficiente pasan a una cola de revisión del equipo antes de que la lista llegue a
las familias.\par

\subsubsection{Carga cognitiva y arquitectura de la decisión}

El último reto se relaciona con la forma en que las personas toman decisiones. La teoría de
la carga cognitiva establece que la memoria de trabajo tiene una capacidad limitada, y que
cuando una tarea exige procesar demasiada información a la vez, la calidad de las decisiones
disminuye \citep{sweller_cognitive_2011}. En la misma línea, se ha demostrado que un número
excesivo de opciones puede desalentar la elección en lugar de facilitarla
\citep{iyengar_when_2000}. La arquitectura de la decisión propone que la forma de presentar
las alternativas, es decir, cuántas son, en qué orden aparecen y cuál se sugiere, influye en
la elección sin limitar la libertad de la persona \citep{thaler_nudge_2008}. A partir de
estos principios, Tukit presenta pocos kits bien diferenciados, cada uno con una razón clara
para elegirlo, y permite que el acudiente modifique cualquier producto antes de pagar.\par





\subsection{Análisis del Contexto}
\label{sec:analisis-contexto}

\subsubsection{Soluciones existentes}

A nivel internacional, varias plataformas han buscado simplificar la compra de útiles
escolares a partir de la lista que publica cada institución. En Estados Unidos, TeacherLists
reúne las listas de cerca de dos millones de salones de clase y las pone a disposición de
cadenas minoristas, que pueden convertir cada lista en un carrito de compra prellenado
\citep{teacherlists_retail_2026}. Amazon School Lists ofrece un servicio similar dentro de
su propia tienda, en el que los docentes publican la lista y los padres agregan sus
artículos al carrito \citep{amazon_schoollists_2025}. Por su parte, 1st Day
School Supplies vende kits armados con los artículos exactos que solicita cada docente y los
entrega en el hogar o en el colegio \citep{firstday_faqs_2026}. Estas soluciones confirman
que pasar de la lista al carrito de forma automática es viable y que las familias valoran
esa comodidad. Sin embargo, todas dependen de que la lista se registre manualmente en su
plataforma con un formato estructurado, ofrecen una sola forma de completarla y no están
disponibles en Colombia. En el caso de 1st Day School Supplies, además, solo pueden
comprar las familias de colegios aliados, dentro de fechas límite de pedido y con las
marcas que define la empresa.\par

En el contexto colombiano, las alternativas disponibles tampoco parten de la lista oficial
de cada colegio. Panamericana y Éxito organizan su oferta escolar por categorías o por
nivel educativo \citep{panamericana_cotizador_2025,exito_listas_2025}, Primavera permite
que el acudiente arme manualmente su lista sobre el catálogo propio
\citep{primavera_armatulista_2026}, y MercadoLibre comercializa kits fijos que arma cada
vendedor \citep{mercadolibre_kits_2025}. En todos los casos, la familia debe traducir por su
cuenta cada renglón de la lista en un producto, comparar alternativas y verificar que la
compra quede completa, lo que mantiene la carga cognitiva del proceso tradicional
\citep{sweller_cognitive_2011}.\par

En el ámbito académico, la recomendación de paquetes estudia cómo sugerir conjuntos de
productos coherentes entre sí, por ejemplo mediante representaciones vectoriales de los
artículos que componen cada paquete \citep{avny_brosh_bruce_2022,du_enhancing_2025} o
mediante modelos de lenguaje que generan o evalúan paquetes completos
\citep{liu_llmcbr_2025,bayrak_evaluating_2024}. De forma paralela, la generación aumentada
por recuperación se ha aplicado al comercio electrónico para responder preguntas sobre
productos \citep{xu_itemrag_2026} o completar las búsquedas de los usuarios
\citep{sun_product-aware_2024}. Estos trabajos aportan técnicas valiosas, pero suponen un
historial de interacciones o una consulta escrita por el usuario, y no abordan el caso en
el que el punto de partida es un documento institucional sin formato fijo que debe
cumplirse completo.\par

\subsubsection{Alternativas de solución consideradas}

Ante este panorama, se analizaron distintas formas de resolver el problema, considerando su
viabilidad técnica y su efecto sobre la experiencia de las familias y la operación del
servicio. A continuación, se describen estas alternativas.\par

\begin{enumerate}
  \item \textbf{Listas digitalizadas manualmente.} Siguiendo el modelo de TeacherLists, el
        colegio o el equipo de Tukit transcribe cada lista en un formulario estructurado, y
        cada renglón se asocia de antemano con un producto.
        \begin{itemize}
          \item \textit{Fortaleza:} la correspondencia entre renglón y producto es exacta,
                porque la define una persona.
          \item \textit{Debilidad:} exige transcribir cada lista y cada cambio de lista, no
                permite que el acudiente cargue la suya y no escala a cientos de colegios.
        \end{itemize}

  \item \textbf{Kits preconfigurados por colegio y grado.} Siguiendo el modelo de 1st Day
        School Supplies, se arma un único kit fijo para cada lista.
        \begin{itemize}
          \item \textit{Fortaleza:} la compra se resuelve en un solo paso y la operación es
                sencilla.
          \item \textit{Debilidad:} ofrece una sola alternativa, sin considerar el
                presupuesto ni las preferencias de cada familia, y requiere armar y
                actualizar manualmente cada kit.
        \end{itemize}

  \item \textbf{Generación directa con un modelo de lenguaje.} Un modelo de lenguaje recibe
        la lista y redacta el kit completo, siguiendo el esquema de generación aumentada
        por recuperación \citep{lewis_retrieval-augmented_2020}.
        \begin{itemize}
          \item \textit{Fortaleza:} se adapta con facilidad a listas redactadas de formas
                muy distintas.
          \item \textit{Debilidad:} puede proponer productos, cantidades o precios que no
                están en el catálogo, sus resultados son difíciles de auditar y no garantiza
                el mejor equilibrio entre costo, calidad y preferencias.
        \end{itemize}

  \item \textbf{Reconocimiento óptico de caracteres con reglas y búsqueda exacta.} La lista
        se lee con OCR tradicional \citep{smith_overview_2007} y cada renglón se busca en el
        catálogo por coincidencia exacta de palabras.
        \begin{itemize}
          \item \textit{Fortaleza:} es económica y sus resultados son fáciles de explicar.
          \item \textit{Debilidad:} falla con fotografías inclinadas, con tablas de varias
                columnas y con las múltiples formas de nombrar un mismo artículo.
        \end{itemize}

\item \textbf{Análisis de maquetación con Docling.} Docling, desarrollado por IBM,
      combina el reconocimiento de caracteres con modelos que identifican la estructura
      del documento, como títulos, párrafos, tablas y columnas, y la convierte en un
      formato organizado \citep{auer_docling_2024}.
      \begin{itemize}
        \item \textit{Fortaleza:} recupera con fidelidad las tablas y las columnas de los
              documentos PDF con una estructura definida, es de código abierto y puede
              ejecutarse en la infraestructura propia.
        \item \textit{Debilidad:} su desempeño disminuye con las fotografías tomadas con
              el celular, que son frecuentes entre las listas que comparten las familias,
              y no interpreta el contenido. Por ejemplo, no reconoce que ``Cartuchera con
              esferos y borrador'' describe tres artículos que se compran por separado, de
              modo que la separación de renglones y su normalización quedarían por
              resolver.
      \end{itemize}
        \item \textbf{Plataformas comerciales de análisis documental.} Servicios como Amazon
      Textract \citep{aws_textract_2026} e IBM Watson Discovery \citep{ibm_discovery_2026}
      extraen el texto, las tablas y los campos de un documento en la nube, y permiten
      buscar información dentro de grandes colecciones de archivos.
      \begin{itemize}
        \item \textit{Fortaleza:} ofrecen alta precisión en la extracción de texto y
              tablas, y no requieren desarrollar ni mantener un modelo propio.
        \item \textit{Debilidad:} su costo crece con cada página procesada y los documentos
              se tratan en la infraestructura de un tercero. Además, solo resuelven la
              lectura, ya que no interpretan el vocabulario escolar ni asocian cada
              renglón con un producto del catálogo, de modo que la normalización y la
              construcción de los kits quedarían por resolver.
      \end{itemize}
\end{enumerate}



Frente a estas alternativas, la solución presentada en la Propuesta de Solución combina
lo mejor de cada una. Conserva la lista oficial del colegio como punto de partida, igual que
las listas digitalizadas, pero la lee de forma automática en cualquier formato mediante un
modelo multimodal \citep{shen_ocr_2026}. Aprovecha la flexibilidad de los modelos de
lenguaje, pero solo para leer la lista, mientras que la asociación con el catálogo se
resuelve con una búsqueda híbrida que no inventa productos \citep{karpukhin_dense_2020}. En
lugar de un único kit fijo, construye varias alternativas que equilibran el costo, la
calidad y las preferencias mediante optimización multiobjetivo \citep{deb_fast_2002}, y
mantiene la revisión humana en los casos dudosos, lo que permite auditar cada resultado.\par

\subsubsection{Comparación con la solución propuesta}

La Tabla~\ref{tab:comparacion-soluciones} resume cómo se comparan las soluciones existentes
con Tukit en los aspectos que más pesan para las familias colombianas.\par

% Tabla tab:comparacion-soluciones movida al Anexo A.3 (tablas del cuerpo)

    \setcounter{subsection}{0}
\section{ANÁLISIS DEL PROBLEMA}

% =====================================================================
\subsection{Requerimientos}

En esta sección se resumen los requerimientos del sistema y se detallan los más importantes.
Para ello, primero se presentan los actores que intervienen en el problema y el proceso que
se siguió para levantar sus necesidades, y luego se explica cómo se organizaron y
priorizaron los requerimientos resultantes.\par

\subsubsection{Actores del sistema}

La caracterización del dominio identificó cuatro actores con intereses distintos frente al
sistema. La Tabla~\ref{tab:actores} los presenta junto con su relación con Tukit y su nivel
de poder e interés, es decir, qué tanto pueden influir en el éxito del sistema y qué tanto
les afecta su resultado. El colegio ocupa la posición más estratégica, porque controla el
acceso a las familias mediante la lista oficial y, a la vez, tiene interés en que sus
acudientes compren con menos dificultad.\par

\begin{table}[H]
\centering\footnotesize
\setlength{\tabcolsep}{4pt}
\renewcommand{\arraystretch}{1.25}
\caption{Actores del sistema y su relación con Tukit.}
\label{tab:actores}
\begin{tabularx}{\textwidth}{@{}p{2.6cm}Xp{3.2cm}@{}}
\toprule
\textbf{Actor} & \textbf{Relación con el sistema} & \textbf{Poder e interés} \\
\midrule
Acudiente & Usuario final que recibe, ajusta y paga el kit. & Poder medio, interés alto \\
Colegio & Publica la lista oficial y es el canal de confianza hacia las familias. & Poder
alto, interés alto \\
Equipo de Tukit & Opera el catálogo, revisa las listas y valida los renglones dudosos. &
Poder alto, interés alto \\
Estudiante & Influye en la elección del kit, por ejemplo con los personajes que prefiere,
aunque no usa el sistema. & Poder bajo, interés alto \\
\bottomrule
\end{tabularx}
\end{table}

\subsubsection{Proceso de levantamiento}

Los requerimientos se obtuvieron principalmente a partir de una encuesta a acudientes del
Colegio Agustiniano Ciudad Salitre, aplicada de forma presencial en el Dia de la Familia en el
colegio con autorización del rector. Esta modalidad permitió reunir en pocas horas a familias de todos los niveles,
desde preescolar hasta la educación media, y conversar en paralelo con docentes y directivos
sobre la forma en que la institución elabora sus listas.\par

El cuestionario constaba de dieciocho preguntas, diecisiete cerradas y una abierta, y su
tiempo de respuesta tuvo una mediana de 4,6 minutos. Su diseño siguió un enfoque que profundiza en qué intenta lograr la persona cuando compra, más
allá del producto en sí, e incluyó preguntas sobre el tiempo, los desplazamientos y las
dificultades que le exige la compra. Las
preguntas abordaron el perfil del hogar, el tiempo y los lugares que demanda la compra, la
principal dificultad que enfrenta la familia, los criterios con que elige los productos, su
disposición a pagar un sobreprecio, la satisfacción del estudiante con sus útiles y la
confianza que le daría el respaldo del colegio. Se recibieron 65 respuestas, de las cuales
se analizaron 60, ya que las cinco restantes no autorizaron el tratamiento de sus datos,
conforme a la Ley 1581 de 2012 \citep{congreso_ley1581_2012}.\par

Los hogares encuestados pertenecen en su mayoría al estrato 4 (68,3\,\%), el 78,3\,\% tiene
un solo hijo en el colegio y el 78,3\,\% compra por aplicaciones móviles al menos algunas
veces al mes, lo que indica que la población está familiarizada con la compra en línea. A
partir de las respuestas, cada hallazgo se tradujo en uno o varios requerimientos del
sistema, como resume la Tabla~\ref{tab:encuesta-requerimientos}.\par

\begin{table}[H]
\centering\footnotesize
\setlength{\tabcolsep}{4pt}
\renewcommand{\arraystretch}{1.25}
\caption{Hallazgos de la encuesta a acudientes y requerimientos que se derivan de ellos
($n=60$).}
\label{tab:encuesta-requerimientos}
\begin{tabularx}{\textwidth}{@{}XX@{}}
\toprule
\textbf{Hallazgo} & \textbf{Requerimiento que se deriva} \\
\midrule
El 65\,\% dedica cuatro horas o más a la compra de los útiles de un hijo, y el 68,3\,\%
compra en tres o más lugares. &
Entregar la lista completa resuelta en un solo lugar, en lugar de producto por producto. \\
El 76,7\,\% recibe la lista con más de un mes de anticipación. &
El problema no es la falta de tiempo para comprar, sino el esfuerzo que exige hacerlo, de
modo que el sistema debe reducir ese esfuerzo. \\
La dificultad más mencionada (43,3\,\%) es no encontrar el producto o encontrarlo agotado, y
el 23,3\,\% menciona los libros y textos escolares. &
Considerar la disponibilidad al armar los kits e informar con claridad los renglones que no
se venden en la plataforma, como los textos que distribuyen las editoriales. \\
El precio es el primer criterio para el 38,3\,\%, las marcas exactas que pide el colegio
para el 31,7\,\%, la calidad para el 18,3\,\% y los personajes que prefiere el hijo para el
11,7\,\%. &
Armar los kits equilibrando costo, fidelidad de marca, calidad y personajes, y ofrecer un
kit para cada prioridad. \\
El 70\,\% pagaría como máximo un 5\,\% más por recibir el kit armado, y el 30\,\% aceptaría
un sobreprecio de 10\,\% o más. &
Recomendar por omisión el kit de menor costo, y sugerir uno de mayor calidad solo a quien
declare que acepta pagar más. \\
El 90\,\% confiaría más en la plataforma si el colegio la recomienda (58,3\,\% ``mucho
más''). &
Operar con un modelo B2B2C, en el que el colegio publica la lista oficial y respalda el
servicio. \\
El 58,3\,\% usaría el servicio con probabilidad alta (4 o 5 en una escala de 1 a 5). &
Validar la propuesta de valor con un piloto en el mismo colegio. \\
\bottomrule
\end{tabularx}
\end{table}

En las conversaciones con docentes y directivos se trataron la forma en que el colegio
elabora y publica sus listas y los inconvenientes que observa cuando las familias compran
productos distintos a los solicitados. De ellas surgieron los requerimientos de la
institución, como publicar la lista oficial de cada grado, garantizar que las familias
compren sobre esa versión y no exigir marcas ni proveedores específicos.\par

Como complemento, se analizaron las listas reales de diecisiete colegios del sector, lo que
permitió conocer la variedad de formatos y de formas de nombrar los útiles y definir los
requerimientos de lectura y normalización. Con toda esta información se elaboró la
especificación de requerimientos siguiendo el estándar IEEE 830 y su sucesor, ISO/IEC/IEEE
29148 \citep{ieee_830_1998,iso_29148_2018}.\par
\subsubsection{Organización y priorización}

Los requerimientos se agrupan en tres categorías. Los requerimientos funcionales describen lo
que el sistema debe hacer, como leer una lista o generar los kits, y se organizan según el
componente del sistema que los atiende. Los requerimientos no funcionales describen cómo
debe comportarse el sistema, por ejemplo en tiempo de respuesta, seguridad o facilidad de
uso, y se organizan según los atributos de calidad de la norma ISO/IEC 25010
\citep{iso_25010_2011}. Por último, los requerimientos propios del aprendizaje automático
describen las condiciones que deben cumplir los modelos de inteligencia artificial, como la
calidad mínima de la lectura de listas o el comportamiento ante casos de baja certeza.\par

Cada requerimiento se priorizó con el método MoSCoW, que lo clasifica como indispensable
(\emph{Must}), importante pero no crítico (\emph{Should}), deseable (\emph{Could}) o fuera
del alcance de esta versión (\emph{Won't}). La especificación reúne 98 requerimientos, de los
cuales 43 son funcionales, 39 no funcionales y 16 propios del aprendizaje automático. La
Tabla~\ref{tab:moscow} (Anexo~\ref{anx:tablas}) muestra su distribución por componente.\par

% [movida al Anexo: tab:moscow]

La proporción de requerimientos indispensables es alta, y responde a la naturaleza del
sistema. Una obligación legal no puede clasificarse como opcional, y tampoco pueden serlo los
requerimientos que definen qué ocurre cuando algo sale mal, por ejemplo cuando un pago falla
o cuando la pasarela de pagos envía dos veces la misma confirmación. Sin ellos, el sistema
podría cobrar dos veces o dejar un pedido a medias, lo que resulta inaceptable en una
plataforma que trata datos de menores y recibe dinero de las familias.\par

\subsubsection{Requerimientos más importantes}

A continuación se presentan los requerimientos más relevantes del sistema, organizados en no
funcionales y funcionales. Cada requerimiento se identifica con un código propio de esta
sección, seguido de su nombre y, en el caso de los no funcionales, de la característica de
calidad de la norma ISO/IEC 25010 a la que pertenece \citep{iso_25010_2011}.\par

\paragraph{Requerimientos no funcionales.} Estos requerimientos definen cómo debe comportarse
el sistema. Los más importantes se relacionan con la protección de los datos de las familias
y de los menores, con la seguridad de los pagos y con el tiempo de respuesta, ya que de ellos
depende la confianza del acudiente en la plataforma.\par

\begin{itemize}
  \item \textbf{\textit{RNF01 -- No almacenar datos de pago (Seguridad)}}
    \begin{itemize}
      \item \textit{Descripción:} el sistema no debe guardar números de tarjeta ni datos
            bancarios. La captura del medio de pago debe ocurrir por completo en la pasarela
            de pagos.
      \item \textit{Justificación:} protege la información financiera de las familias y
            reduce el riesgo de que una falla de seguridad en Tukit exponga datos de pago.
    \end{itemize}

  \item \textbf{\textit{RNF02 -- Protección de los datos de los menores (Seguridad)}}
    \begin{itemize}
      \item \textit{Descripción:} el sistema debe guardar solo los datos indispensables del
            estudiante, cifrarlos y eliminar su nombre de la lista antes de enviarla al
            modelo de lectura, que es un servicio externo.
      \item \textit{Justificación:} la Ley 1581 de 2012 y el Decreto 1377 de 2013 dan un
            tratamiento especial a los datos de niños y adolescentes
            \citep{congreso_ley1581_2012,presidencia_decreto1377_2013}.
    \end{itemize}

  \item \textbf{\textit{RNF03 -- Tiempo de generación de los kits (Eficiencia de
        desempeño)}}
    \begin{itemize}
      \item \textit{Descripción:} una vez normalizada la lista, el sistema debe generar los
            kits en 20 segundos o menos en el 95\,\% de los casos, y el proceso completo,
            desde la carga de la lista hasta la creación de los kits, no debe superar los
            cinco minutos.
      \item \textit{Justificación:} el acudiente espera este resultado frente a la pantalla,
            y una espera excesiva lleva a abandonar el proceso.
    \end{itemize}

  \item \textbf{\textit{RNF04 -- Acceso desde el celular sin instalación (Usabilidad)}}
    \begin{itemize}
      \item \textit{Descripción:} la aplicación del acudiente debe funcionar desde el
            navegador del celular como una aplicación web progresiva, sin necesidad de
            instalarla.
      \item \textit{Justificación:} la mayoría de los acudientes compra desde el celular, y
            evitar la instalación elimina una barrera para el primer uso.
    \end{itemize}



\item \textbf{\textit{RNF05 -- Capacidad en temporada escolar (Eficiencia de
      desempeño)}}
  \begin{itemize}
    \item \textit{Descripción:} el sistema debe atender a 50 usuarios simultáneos y del
          orden de 1.000 pedidos en el día de mayor demanda, sin que los tiempos de
          respuesta se deterioren de forma significativa y con una tasa de error inferior
          al 5\,\%.
    \item \textit{Justificación:} la demanda se concentra en pocas semanas del año. La
          meta equivale a diez veces los cerca de 95 pedidos diarios que se proyectan para
          la temporada del primer año, lo que deja un margen amplio frente a la demanda
          esperada.
  \end{itemize}

\item \textbf{\textit{RNF06 -- Crecimiento sin rediseño (Escalabilidad)}}
  \begin{itemize}
    \item \textit{Descripción:} la arquitectura debe permitir crecer hasta 1.000 usuarios
          simultáneos aumentando el número de copias de cada servicio, sin cambiar su
          diseño ni su modelo de datos.
    \item \textit{Justificación:} el sistema debe poder incorporar nuevos colegios y
          ciudades después del piloto sin reconstruir la plataforma.
  \end{itemize}

\item \textbf{\textit{RNF07 -- Disponibilidad en las horas de mayor uso (Fiabilidad)}}
  \begin{itemize}
    \item \textit{Descripción:} el sistema debe estar disponible al menos el 90\,\% del
          tiempo durante las franjas de mayor demanda, de 12:00 a 14:00 y de 18:00 a 22:00.
          El mantenimiento programado se realiza entre las 02:00 y las 03:00.
    \item \textit{Justificación:} en esas franjas se concentra la actividad de las familias,
          y una interrupción en temporada escolar puede llevar al acudiente a comprar por
          otro medio.
  \end{itemize}

\item \textbf{\textit{RNF08 -- Tolerancia a fallas (Fiabilidad)}}
  \begin{itemize}
    \item \textit{Descripción:} la falla de un servicio no debe afectar a los demás. Si el
          servicio externo de lectura deja de responder, el sistema debe seguir mostrando
          el catálogo y los kits ya generados, y ofrecer al acudiente la opción de ingresar
          su lista manualmente.
    \item \textit{Justificación:} la lectura depende de un proveedor externo que puede
          fallar, y esa falla no debe impedir que las familias consulten o compren lo que
          ya está listo.
  \end{itemize}
\end{itemize}

\paragraph{Requerimientos funcionales.} Estos requerimientos describen lo que el sistema
debe hacer, y se agrupan según el recorrido de una lista, desde su carga hasta la compra.
Se enfocan en los tres actores que usan la plataforma: el acudiente, el colegio y el equipo
de Tukit.\par

\begin{itemize}
  \item \textbf{\textit{RF01 -- Cargar y leer la lista}}
    \begin{itemize}
      \item \textit{Descripción:} el acudiente debe poder seleccionar la lista oficial de su
            colegio y grado, o cargar la suya como fotografía, PDF o texto, y el sistema
            debe convertirla en renglones con su texto y su cantidad.
      \item \textit{Justificación:} cada colegio publica su lista en un formato distinto, y
            sin esta capacidad el acudiente tendría que transcribirla a mano.
    \end{itemize}

  \item \textbf{\textit{RF02 -- Reconocer las distintas formas de nombrar un artículo}}
    \begin{itemize}
      \item \textit{Descripción:} el sistema debe asociar cada renglón con su concepto del
            catálogo, aunque esté escrito con sinónimos o errores de ortografía, y enviar a
            revisión los renglones que no logre resolver con certeza suficiente.
      \item \textit{Justificación:} las listas de diecisiete colegios nombran los útiles de
            1.451 formas que corresponden a 99 conceptos, y sin este mecanismo la mayoría de
            los renglones quedaría sin producto.
    \end{itemize}

  \item \textbf{\textit{RF03 -- Generar los kits}}
    \begin{itemize}
      \item \textit{Descripción:} el sistema debe generar hasta cuatro kits completos
            (Ahorrador, Premium, Fan y A tu Medida), con un máximo de tres opciones de
            producto por renglón.
      \item \textit{Justificación:} ofrecer pocas alternativas bien diferenciadas facilita la
            decisión del acudiente \citep{iyengar_when_2000}.
    \end{itemize}

  \item \textbf{\textit{RF04 -- Personalizar y pagar el kit}}
    \begin{itemize}
      \item \textit{Descripción:} el acudiente debe poder cambiar cualquier producto, elegir
            el kit que prefiera, pagarlo en línea y solicitar el retracto dentro del plazo
            legal.
      \item \textit{Justificación:} la recomendación orienta la decisión sin imponerla, y la
            Ley 1480 de 2011 obliga a garantizar el retracto en las compras a distancia
            \citep{congreso_ley1480_2011}.
    \end{itemize}

  \item \textbf{\textit{RF05 -- Publicar y revisar las listas oficiales}}
    \begin{itemize}
      \item \textit{Descripción:} el colegio debe poder publicar la lista oficial de cada
            grado, y el equipo de Tukit debe revisarla y validar sus renglones antes de que
            llegue a las familias.
      \item \textit{Justificación:} garantiza que las familias compren sobre la versión
            oficial de la lista y no sobre copias desactualizadas.
    \end{itemize}
\end{itemize}
% =====================================================================
\subsection{Restricciones}
\label{sec:restricciones}

La restricción de mayor peso sobre el diseño es legal, no técnica. La Ley 1581 de 2012
\citep{congreso_ley1581_2012} establece el régimen general de protección de datos
personales en Colombia, y el Decreto 1377 de 2013 \citep{presidencia_decreto1377_2013} lo
desarrolla con un régimen especial para los datos de niños, niñas y adolescentes. Una lista
escolar es siempre la lista de un menor, de modo que la restricción alcanza todo el flujo y
no un módulo aislado. El sistema trata los datos para una sola finalidad, la prestación del
servicio, y atiende los derechos de los artículos 14 y 15 de la ley mediante un caso de
soporte dentro de los plazos legales.

La propuesta inicial era
anonimizar el nombre del menor antes de guardar la lista y antes de invocar el modelo
externo. El sistema construido protege ese dato de otra forma, por minimización: no lo recoge. No pide el
nombre legal; el apodo con que el acudiente identifica a cada hijo es opcional y se cifra en
reposo, y la percepción extrae solo los ítems y el grado (PV-03). Es mejor porque un dato
que no se recoge no puede llegar a los registros de operación, a los conjuntos de evaluación
ni al modelo de percepción, y no depende de que un reconocedor de entidades lo detecte.
Queda un caso: el nombre que figure en el documento cargado viaja con la página hacia la API
externa, y su depuración (PE-06) está pendiente.

La segunda restricción proviene del dominio de pagos. El sistema no almacena ningún dato del
medio de pago del acudiente: ni número de tarjeta, ni fecha de expiración, ni código de
seguridad. Solo guarda el identificador de la transacción y su estado, porque la captura del
medio de pago ocurre en el dominio de Wompi, que el sistema usa en su entorno de pruebas
detrás de un puerto de pasarela con un adaptador simulado (ADR 0010). Esta restricción
mantiene a Tukit fuera del alcance de cumplimiento de los datos de tarjeta. La confirmación
no depende de la respuesta sincrónica de la pasarela: llega por el aviso firmado de Wompi y
por la consulta de la transacción cuando el acudiente regresa, y ambos caminos pasan por un
mismo procedimiento idempotente, porque el aviso puede llegar duplicado o fuera de orden, o
perderse.

Las restricciones técnicas se derivan del ADR 0013, un monorrepositorio con despliegue
independiente por servicio. El sistema construido tiene seis desplegables con código, cinco
servicios de dominio y una puerta de enlace, empaquetados en contenedores estándar para que
el entorno de evaluación sea el mismo que el de desarrollo. Hoy se levantan con Docker
Compose sobre un único PostgreSQL 16, con una base lógica por servicio, RabbitMQ y Redis; la
orquestación productiva se decidirá con el despliegue de finales de noviembre de 2026. Cada
evento se escribe en una bandeja de salida en la misma transacción que el cambio de estado, y
RabbitMQ lo lleva a la cola de cada consumidor (ADR 0037). Sin RabbitMQ, los consumidores leen
las bandejas, y el sistema funciona completo con una sola base de datos, a costa de concentrar
en ella la coordinación.

% =====================================================================
\subsection{Especificación Funcional}
\label{sec:especificacion-funcional}

El SRS v6.1 define cuatro actores. El acudiente (padre, madre o tutor) es el comprador final:
carga la lista, declara sus preferencias, elige un kit y paga, y solo ve sus propias
solicitudes, pedidos y casos. El operador de Tukit revisa las listas y los kits: atiende la
cola de revisión y valida las líneas, pero no publica listas oficiales; también opera la
logística de los pedidos. El administrador de
Tukit tiene los permisos del operador y, además, publica las listas oficiales (CO-01), consulta
las solicitudes de kits y gestiona el catálogo, la logística y los casos de soporte desde el
panel de administración. El niño o la niña no opera ninguna pantalla: influye en la compra,
sobre todo en el kit Fan, pero es sujeto de datos protegidos y no usuario. La entrega física
tampoco tiene un actor propio: la hacen repartidores propios, y el despacho exige el nombre y
el celular de quien entrega (UM-03).

El flujo funcional encadena cuatro etapas. La lista escolar, fotografiada o entregada como
documento, pasa primero por la \emph{percepción}, que la convierte en texto estructurado
con un modelo multimodal, y luego por la \emph{normalización}, que asigna cada renglón de
texto libre a un concepto del catálogo mediante recuperación híbrida. Los renglones que la
normalización no cierra con certeza suficiente pasan a la \emph{resolución}. La
\emph{optimización}, por último, elige sobre el catálogo los productos que mejor equilibran
costo, fidelidad de marca, disponibilidad y personaje, con pesos derivados de lo que el
acudiente declara por cada hijo: cuánto quiere gastar y, si quiere, el personaje favorito
(PS-01). El resultado son los kits que el acudiente pidió, de los cuales lleva uno al pago:
el Ahorrador y el Premium siempre, el Fan si indicó un personaje y «A tu Medida» si eligió el
punto medio de gasto o declaró un tope.

La resolución concentra la dificultad de la normalización. Cuando la confianza de una
asignación cae por debajo del umbral de PE-04, el ítem queda marcado como ambiguo y espera
una decisión humana en la cola de revisión del panel (AD-01), donde cada corrección queda
registrada con su autor y su marca de tiempo (AD-02).

El
diseño preveía que una persona revisara solo los ítems por debajo del umbral. En el sistema
construido, el equipo revisa una vez cada lista nueva completa antes de publicarla, con el
documento original y la lectura del sistema lado a lado, y con un plazo de 4 h entre la
carga y la aprobación (PE-04). Es mejor porque detecta también los errores leídos con
confianza alta, que un filtro por ítem deja pasar, y porque la lista revisada sirve después a
todos los acudientes del mismo colegio y grado (NR-07). La validación de este flujo con
usuarios está planificada y pendiente (Sección~\ref{sec:oe5}).

La magnitud del problema la fijó el levantamiento de listas escolares de diecisiete
colegios: 3.608 ítems expresados en 1.451 formas distintas de nombrar los artículos, que el
catálogo reduce a 99 conceptos canónicos. Pocos conceptos concentran casi todo el volumen.
Según las métricas del catálogo, que registran 3.610 renglones, 3.600 de ellos con concepto
asignado, 25 conceptos cubren el 80,97\,\% de estos y 38 cubren el 90\,\%. Los totales de las
distintas fuentes difieren levemente sin una explicación registrada (anexo CRISP-DM). El
resto es una cola larga de variaciones superficiales de los mismos conceptos. Esa
distribución explica por qué ninguna solución de coincidencia exacta de cadenas resuelve el
problema y por qué la resolución es una parte estructural del flujo y no una salvedad.

    \setcounter{subsection}{0}
\section{DISEÑO DE LA SOLUCIÓN}

Este capítulo presenta el resultado de la Fase~2: la arquitectura de Tukit y las decisiones del
diseño detallado. Su
argumento central es que las fronteras entre servicios se justifican por las propiedades que
los datos no pueden compartir, no por la carga que el sistema va a soportar.

La exposición sigue la norma ISO/IEC/IEEE 42010 \citep{iso_42010_2022} en la separación
entre preocupaciones de los interesados, puntos de vista y vistas, y el modelo 4+1 de
Kruchten \citep{kruchten_41_1995} para distinguir la vista de desarrollo de la vista
física, distinción que resulta central en la Sección~\ref{sec:monorepo}. Las alternativas se
evaluaron con ATAM y las decisiones quedaron en registros de decisión de arquitectura.

% =====================================================================
\subsection{Impulsores arquitectónicos}
\label{sec:impulsores}

El diseño inicial justificaba la descomposición en microservicios por
\emph{escalabilidad}. Un cálculo de carga con las cifras del mercado objetivo (diecisiete
colegios en el corpus y una demanda concentrada en las semanas previas al inicio del año
lectivo) mostró que la infraestructura mínima excede la demanda proyectada en unos tres
órdenes de magnitud. Esa justificación no se sostenía. El rediseño condujo a tres criterios,
ninguno de rendimiento, que resume la Tabla~\ref{tab:impulsores}.

\begin{table}[htbp]
\centering
\caption{Impulsores arquitectónicos de la descomposición en servicios}
\label{tab:impulsores}
\small
\begin{tabular}{@{}p{3.1cm}p{5.0cm}p{5.6cm}@{}}
\toprule
\textbf{Criterio} & \textbf{Fundamento} & \textbf{Qué separa} \\
\midrule
Recuperabilidad asimétrica &
El objetivo de punto de recuperación difiere en órdenes de magnitud entre almacenes &
El catálogo se reconstruye desde sus archivos de datos, una vez exportada al repositorio la
curaduría del panel, y admite perder un día de cambios (CF-02a); un pedido no puede perder un
solo byte. Compartir almacén obliga a dimensionar el
respaldo por el más exigente y a exponer el crítico a la carga de escritura del volátil \\
\addlinespace
Frontera de cumplimiento &
Ley 1581 de 2012 \citep{congreso_ley1581_2012} y Decreto 1377 de 2013
\citep{presidencia_decreto1377_2013} &
Los datos de menores de edad se aíslan en un servicio con base, proceso y llaves propios.
La superficie que un auditor debe examinar se reduce a un solo componente \\
\addlinespace
Perfil de recurso &
Heterogeneidad de la demanda computacional &
La inferencia de modelos y la optimización multiobjetivo son intensivas en CPU y memoria;
la gestión transaccional requiere disponibilidad y consistencia fuerte. Colocarlas juntas
obliga a dimensionar todo por la más exigente \\
\bottomrule
\end{tabular}
\end{table}

Con estos criterios, el número de servicios que ya corría en los prototipos resultó
aproximadamente adecuado; lo que estaba mal era la ubicación de algunas fronteras.

% =====================================================================
\subsection{De catorce servicios documentados a seis desplegables}
\label{sec:consolidacion}

La arquitectura de los documentos académicos iniciales proponía trece microservicios de
dominio y una puerta de enlace, para separar al máximo las responsabilidades. Dos hechos la
cuestionaron. El primero fue de costo y operabilidad: catorce desplegables implican catorce
canales de integración y despliegue, catorce bases candidatas y una malla de red que un
equipo de cuatro personas no puede operar, frente a una demanda proyectada de unos treinta y
seis pedidos diarios en temporada. El segundo fue empírico: la base de grafos Neo4j,
prevista para la sustituibilidad, estaba apagada por configuración y el grafo de conceptos
cabía en memoria, por lo que el ADR 0006 la retiró y resolvió el grafo en proceso; su
evaluación como GraphRAG consta en el anexo CRISP-DM. El sistema construido sustituyó a su
vez ese grafo por una tabla curada y dirigida de sustituciones en el servicio de catálogo,
que tiene modelo, migración y pruebas. Ningún servicio la consulta en la versión evaluada, de modo que la sustitución se delimitó como línea de la evolución del catálogo (Sección~\ref{sec:trabajo-futuro}).

El ADR 0001 aplicó los tres criterios de la Tabla~\ref{tab:impulsores} y redujo los catorce
servicios a seis más la puerta de enlace. Los microservicios anteriores, desarrollados entre
el 2 y el 18 de septiembre de 2026, hicieron esa fusión el 10 de septiembre. Identidad y
privacidad quedaron en un servicio de cuentas; catálogo, pedidos y última milla, en uno de
comercio, y PQRS, analítica y notificaciones, en uno de plataforma. Aparte quedaron
instituciones, una fachada para el panel de administración y el servicio de búsqueda.

Esa reagrupación
dejaba un servicio de respaldo para el panel de administración y un servicio de
instituciones separado de las cuentas. Con la razón de cambio como criterio, los colegios
pasaron a identidad, el catálogo se separó de comercio (ADR 0011), el panel consume la misma
API que la aplicación del acudiente y la auditoría pasó a ser una biblioteca compartida
(ADR 0002). El servicio de búsqueda anterior se dividió entre catálogo e inteligencia,
porque sus objetivos de recuperación no permiten compartir almacén. El resultado es mejor
para un equipo de cuatro personas: hay menos procesos que operar, cada servicio tiene un
solo motivo para cambiar y ningún contrato se duplica. El sistema construido tiene seis
desplegables con código, cinco servicios de dominio (identidad, catálogo, inteligencia,
comercio y operación) y la puerta de enlace, sobre cinco bases.

El diseño de la
Fase 2 preveía un séptimo desplegable, ingesta, que extraía de forma masiva los productos y
precios de las tiendas. El sistema construido lo reemplazó por un catálogo curado que crece
con la demanda: un producto entra cuando una lista lo pide y el equipo confirma la
correspondencia. Es mejor porque concentra el esfuerzo en los productos que las listas
piden, reduce los candidatos de cada concepto de miles a decenas y evita publicar un
producto parecido pero equivocado. En el sistema construido reúne 1.050 productos (1.040
activos), 1.057 ofertas de ocho tiendas y 149 conceptos en archivos de datos versionados, y las herramientas que
consultan las tiendas solo completan o enriquecen productos. Lo que el equipo cura en el
panel queda en la base del catálogo hasta que se exporta a esos archivos (ADR 0028).

% =====================================================================
\subsection{Alternativas evaluadas}
\label{sec:alternativas}

Se compararon tres alternativas mediante el método ATAM \citep{bass_software_2021}, cuyos
resultados completos se consignan en el documento de evaluación ATAM. La
Tabla~\ref{tab:alternativas} sintetiza el contraste.

\begin{table}[htbp]
\centering
\caption{Alternativas arquitectónicas evaluadas}
\label{tab:alternativas}
\small
\begin{tabular}{@{}p{2.6cm}p{3.5cm}p{3.5cm}p{4.1cm}@{}}
\toprule
& \textbf{A. Monolito modular} & \textbf{B. Microservicios de grano fino} & \textbf{B$'$. Grano grueso (elegida)} \\
\midrule
Componentes & 1 desplegable, módulos internos & 13 servicios y puerta de enlace & 5 servicios de dominio y puerta de enlace, 5 bases \\
\addlinespace
Aislamiento de datos de menores & Por convención de código & Por servicio & Por servicio, base y llaves \\
\addlinespace
Punto de recuperación & Único para todo & Por servicio & Diferenciado por criticidad \\
\addlinespace
Complejidad operativa & Baja & Alta: coordinación distribuida en cada transacción & Media \\
\addlinespace
Costo de infraestructura & Mínimo & Alto & Moderado \\
\addlinespace
Puntaje ponderado (sobre 5,00) & 2,82 & 3,44 & 3,91 \\
\bottomrule
\end{tabular}
\end{table}

La alternativa~A se descarta por el primer criterio de la Tabla~\ref{tab:impulsores}: un
almacén único impone un mismo objetivo de recuperación a datos cuyo valor difiere en
órdenes de magnitud. La alternativa~B se descarta por complejidad operativa: fragmentar la
gestión de pedidos, inventario y pago en servicios separados convierte cada compra en una
transacción distribuida, y el equipo no dispone de los recursos de operación que eso exige
\citep{newman_building_2021}. La alternativa~B$'$ conserva el aislamiento donde el
cumplimiento y la recuperabilidad lo exigen, y mantiene junta la transaccionalidad que
debe permanecer junta. La matriz ponderada del ATAM le da 3,91 sobre 5,00, frente a 3,44 de
la granularidad fina y 2,82 del monolito modular.

La Figura~\ref{fig:c4-contexto} (Anexo~\ref{anx:figuras}) sitúa el sistema frente a sus actores y sistemas externos, y
la Figura~\ref{fig:c4-contenedores} materializa la alternativa elegida: cinco servicios de
dominio y la puerta de enlace sobre cinco bases de datos, con una bandeja de salida en cada
base para entregar los eventos. Las dos salen del código fuente: la primera es una vista UML
del modelo de arquitectura, y la segunda sigue el modelo C4. La vista de contenedores añade los procesos en segundo plano de cada servicio y numera
los pasos de tres flujos: la subida de la lista y los kits, el pago y los avisos.

Con RabbitMQ, un publicador por servicio productor lleva los eventos de su bandeja a una cola
por consumidor, con una cola de reintentos y otra de descartados, y la entrega es al menos una
vez, con consumidores idempotentes (ADR 0037). La bandeja sigue siendo la fuente de verdad. Sin
RabbitMQ, cada consumidor lee la bandeja en la base del productor, sin tope de intentos ni cola
de descartados. El manifiesto declara esas lecturas como respaldo, y una función de aptitud
impide que el camino con RabbitMQ lea bases ajenas.



\begin{figure}[htbp]
\centering
% Vista de contenedores (C4 nivel 2) del código fuente del sistema en su versión del 24 de octubre de 2026,
% con los flujos de la subida y los kits (A), el pago (B) y los avisos (C). Estilos locales: no depende del
% preámbulo; solo usa las bibliotecas arrows.meta y shapes.geometric.
\begin{tikzpicture}[x=1cm, y=1cm, font=\scriptsize,
  per/.style={draw, rounded corners=7pt, align=center, fill=yellow!10, inner sep=3pt, minimum height=7mm},
  cont/.style={draw, rounded corners=2pt, align=center, fill=blue!8, inner sep=2pt, minimum height=8mm},
  proc/.style={draw, rounded corners=2pt, align=center, fill=green!10, inner sep=2pt, minimum height=7mm},
  ext/.style={draw, rounded corners=2pt, align=center, fill=orange!14, inner sep=2pt, minimum height=7mm},
  infra/.style={draw, rounded corners=2pt, align=center, fill=gray!12, inner sep=2pt, minimum height=7mm},
  bd/.style={draw, cylinder, shape border rotate=90, aspect=0.22, align=center, fill=white,
    minimum width=1.25cm, minimum height=8mm, inner sep=1pt},
  grupo/.style={draw=gray!70, rounded corners=3pt, fill=gray!3},
  http/.style={-{Stealth[length=1.6mm]}, semithick},
  evt/.style={-{Stealth[length=1.6mm]}, densely dashed, semithick},
  cola/.style={-{Stealth[length=1.6mm]}, semithick, draw=violet!75!black},
  hueco/.style={draw=white, line width=2.8pt},
  paso/.style={circle, draw, fill=white, inner sep=0.5pt, font=\tiny\bfseries, minimum size=3.6mm},
  rot/.style={font=\tiny, fill=white, inner sep=1pt},
  col/.style={anchor=north west, text width=4.9cm, align=left, inner sep=0pt, font=\scriptsize}]
% Personas y aplicaciones
\node[per, text width=2.6cm] (acudiente) at (2.45,0) {\textbf{Acudiente}};
\node[per, text width=3.8cm] (equipo) at (8.6,0) {\textbf{Equipo de Tukit}\\operador y administrador};
\node[cont, text width=3.4cm] (web) at (2.45,-1.45) {\textbf{Aplicación web progresiva}\\Next.js, puerto 3000};
\node[ext, text width=1.85cm] (google_places) at (5.5,-1.45) {\textbf{Google Maps}\\sugerencias de dirección};
\node[cont, text width=3.4cm] (admin) at (8.6,-1.45) {\textbf{Panel de administración}\\Next.js, puerto 3001};
\node[ext, text width=2.5cm] (wompi) at (13.4,-1.45) {\textbf{Wompi}\\pasarela de pago};
\node[cont, minimum width=14.3cm, text width=14cm] (gateway) at (7.35,-2.85) {\textbf{Puerta de enlace}
  (Python, puerto 8010): reenvía cada petición al servicio que atiende el prefijo de su ruta};
% Servicios que llaman a otros: operación, inteligencia y comercio
\draw[grupo] (0.2,-3.65) rectangle (4.55,-6.25);
\draw[grupo] (4.85,-3.65) rectangle (10.15,-6.25);
\draw[grupo] (10.45,-3.65) rectangle (14.95,-6.25);
\node[cont, text width=2.3cm] (operacion) at (1.6,-4.3) {\textbf{Operación}\\puerto 8015};
\node[bd] (base_operacion) at (3.7,-4.3) {base de\\operación};
\node[proc, text width=1.7cm] (instantanea) at (1.2,-5.55) {instantánea de métricas (a mano)};
\node[proc, text width=1.7cm] (repartidor) at (3.5,-5.55) {repartidor de avisos};
\node[cont, text width=2.1cm] (inteligencia) at (6.4,-4.3) {\textbf{Inteligencia}\\puerto 8013};
\node[bd] (base_inteligencia) at (9.2,-4.3) {base y cola\\de trabajos};
\node[proc, text width=1.45cm] (sondeador_listas) at (5.75,-5.55) {sondeador\\de listas};
\node[proc, text width=1.45cm] (trabajador) at (7.5,-5.55) {trabajador\\durable};
\node[proc, text width=1.45cm] (publicador_inteligencia) at (9.25,-5.55) {publicador\\de eventos};
\node[cont, text width=2.4cm] (comercio) at (11.85,-4.3) {\textbf{Comercio}\\puerto 8014};
\node[bd] (base_comercio) at (13.95,-4.3) {base de\\comercio};
\node[proc, text width=1.6cm] (barredor) at (11.5,-5.55) {barredor de reservas};
\node[proc, text width=1.45cm] (publicador_comercio) at (13.75,-5.55) {publicador\\de eventos};
% Bus de llamadas a identidad y catálogo, que no llaman a otros servicios
\draw[semithick] (2.95,-6.7) -- (12.55,-6.7);
\foreach \x in {2.95,8.55,12.55}{\draw[semithick] (\x,-6.25) -- (\x,-6.7);}
\draw[semithick] (10.3,-3.25) -- (10.3,-6.7);
\node[rot, align=center] at (11.42,-6.7) {llamadas HTTP a\\identidad y catálogo};
% RabbitMQ: intercambio de eventos entre las dos filas de servicios
\node[draw=violet!75!black, fill=violet!10, rounded corners=2pt, minimum width=10.35cm, minimum height=6.5mm,
  inner sep=0pt] (rabbitmq) at (9.125,-7.25) {};
\node[align=center, inner sep=0pt, font=\scriptsize] at (5.95,-7.25) {\textbf{RabbitMQ}\\[-1.5pt]intercambio tukit.eventos};
% Servicios proveedores
\draw[grupo] (2.9,-7.8) rectangle (6.75,-10.3);
\node[cont, text width=1.8cm] (identidad) at (3.95,-8.45) {\textbf{Identidad}\\puerto 8011};
\node[bd] (base_identidad) at (6.0,-8.45) {base de\\identidad};
\node[proc, text width=2.4cm] (sondeador_colegios) at (4.75,-9.65) {sondeador de colegios};
\draw[grupo] (9.05,-7.8) rectangle (13.1,-10.3);
\node[cont, text width=2.0cm] (catalogo) at (10.2,-8.45) {\textbf{Catálogo}\\puerto 8012};
\node[bd] (base_catalogo) at (12.3,-8.45) {base de\\catálogo};
\node[proc, text width=2.0cm, fill=green!4, densely dashed] (herramientas) at (10.2,-9.65) {herramientas de\\curaduría (a mano)};
\node[proc, text width=1.3cm] (publicador_catalogo) at (12.15,-9.65) {publicador\\de eventos};
% Sistemas externos
\node[ext, text width=1.9cm] (gemini) at (7.85,-8.45) {\textbf{Gemini}\\lectura y vectores};
\node[ext, text width=1.7cm] (tiendas) at (14.4,-9) {\textbf{Tiendas en línea}\\productos y precios};
\node[ext, text width=1.9cm] (correo) at (1.25,-8.45) {\textbf{Correo saliente}\\SMTP};
\node[ext, text width=1.9cm] (whatsapp) at (1.25,-9.55) {\textbf{WhatsApp}\\enlace wa.me};
% Fila de infraestructura
\node[infra, text width=4.6cm] (postgres) at (4.9,-11.05) {\textbf{PostgreSQL 16 con pgvector}\\un servidor con las
  cinco bases;\\cada una con su bandeja de salida};
\node[infra, text width=4.2cm, densely dotted, anchor=north] (redis) at (9.6,-10.45) {\textbf{Redis}, opcional\\caché de vectores y de lecturas\\y contador del límite de tasa\\de la puerta de enlace};
\node[infra, text width=3.3cm, densely dotted] (observabilidad) at (13.55,-11.05) {\textbf{Observabilidad opcional}\\OpenTelemetry y Jaeger};
% Personas, aplicaciones y pasarela
\draw[http] (acudiente) -- (web);
\draw[http] (equipo) -- (admin);
\draw[http] (web) -- (google_places);
\draw[http] (web.south) -- (web.south |- gateway.north);
\draw[http] (admin.south) -- (admin.south |- gateway.north);
\draw[http] ([xshift=2mm]web.north west) -- (0.88,0.65) -| node[rot, pos=0.25, above] {abre el widget de pago de Wompi} (wompi.north);
\draw[http] (wompi.south) -- (wompi.south |- gateway.north) node[rot, midway, right] {aviso firmado};
\draw[http] ([xshift=-2mm]comercio.north east) -- ++(0,0.45) -- (15.25,-3.45) |- (wompi.east);
\draw[http] (gateway.south -| operacion) -- (operacion.north);
\draw[http] (gateway.south -| inteligencia) -- (inteligencia.north);
\draw[http] (gateway.south -| comercio) -- (comercio.north);
% Llamadas entre servicios que llaman
\draw[http] ([xshift=-2mm]operacion.north east) -- ++(0,0.35) -| ([xshift=3mm]inteligencia.north west);
\draw[http] ([xshift=-3mm]inteligencia.north east) -- ++(0,0.4) -| ([xshift=3mm]comercio.north west);
% Llamadas a identidad y catálogo desde el bus; la del catálogo cruza RabbitMQ
\draw[http] (3.6,-6.7) -- (3.6,0 |- identidad.north);
\draw[hueco] (10.3,-6.9) -- (10.3,-7.6);
\draw[http] (10.3,-6.7) -- (10.3,0 |- catalogo.north);
% Bandejas de salida que lee cada publicador
\draw[evt] (base_inteligencia.south) -- (base_inteligencia.south |- publicador_inteligencia.north);
\draw[evt] (base_comercio.south) -- (base_comercio.south |- publicador_comercio.north);
\draw[evt] (base_catalogo.south) -- (base_catalogo.south |- publicador_catalogo.north);
% Eventos por RabbitMQ: de los publicadores al intercambio y del intercambio a la cola de cada consumidor
\foreach \x in {4.1,5.75,7.1,9.25}{\draw[hueco] (\x,-6.58) -- (\x,-6.82);}
\draw[cola] (publicador_inteligencia.south) -- (9.25,0 |- rabbitmq.north);
\draw[cola] (publicador_comercio.south) -- (13.75,0 |- rabbitmq.north);
\draw[cola] (11.5,0 |- publicador_catalogo.north) -- (11.5,0 |- rabbitmq.south);
\draw[cola] (4.1,0 |- rabbitmq.north) -- (4.1,0 |- repartidor.south);
\draw[cola] (5.75,0 |- rabbitmq.north) -- (sondeador_listas.south);
\draw[cola] (7.1,0 |- rabbitmq.north) -- (7.1,0 |- trabajador.south);
\draw[cola] (5.2,0 |- rabbitmq.south) -- (5.2,0 |- sondeador_colegios.north);
% Llamadas a sistemas externos
\draw[hueco] (7.85,-6.58) -- (7.85,-6.82);
\draw[hueco] (7.85,-6.9) -- (7.85,-7.6);
\draw[http] (7.85,0 |- trabajador.south) -- (gemini.north);
\draw[http] (catalogo.west) -- (gemini.east);
\draw[http] (13.1,-9) -- (tiendas.west);
\draw[http] (2.75,0 |- repartidor.south) |- (correo.east);
\draw[semithick] (2.75,0 |- repartidor.south) |- (whatsapp.east);
% Pasos de los flujos
\node[paso] at (2.15,-2.15) {A1};
\node[paso] at (7.85,-7.73) {A2};
\node[paso] at (8.3,-2.15) {A3};
\node[paso] at (5.75,-6.45) {A4};
\node[paso] at (2.75,-2.15) {B1};
\node[paso] at (11.6,0.65) {B2};
\node[paso] at (13.05,-2.15) {B3};
\node[paso] at (barredor.north west) {B4};
\node[paso] at (4.1,-6.45) {C1};
\node[paso] at (2.95,-6.45) {C2};
\node[paso] at (2.75,-9) {C3};
% Pasos de los flujos y notación, debajo de la figura
\node[col] (colA) at (0.2,-12.1) {\textbf{A. Subida de la lista y kits}\\[1pt]
  \textbf{A1} El acudiente sube la lista; inteligencia guarda el trabajo y lo anuncia en su bandeja de salida.\\[1pt]
  \textbf{A2} El trabajador toma el trabajo de RabbitMQ o de su tabla, lee las páginas con Gemini, pide candidatos al catálogo y deja los kits por revisar.\\[1pt]
  \textbf{A3} El equipo revisa los kits en el panel y los publica; inteligencia deja el evento de kits publicados en su bandeja de salida.\\[1pt]
  \textbf{A4} Si la lista no estaba verificada, el sondeador de listas arma el kit al recibir el evento de lista verificada.};
\node[col] (colB) at (5.35,-12.1) {\textbf{B. Pago}\\[1pt]
  \textbf{B1} El acudiente elige un kit y escribe la dirección con sugerencias de Google Maps; inteligencia crea la orden en comercio, que reserva el inventario al precio vigente.\\[1pt]
  \textbf{B2} La aplicación web abre el widget de Wompi y el acudiente paga.\\[1pt]
  \textbf{B3} Wompi envía el aviso firmado o la página de resultado pide sincronizar el pago; comercio deja el evento de pedido pagado en su bandeja.\\[1pt]
  \textbf{B4} El barredor libera la reserva del pedido que no se paga en 24 horas y lo cancela.};
\node[col] (colC) at (10.5,-12.1) {\textbf{C. Avisos}\\[1pt]
  \textbf{C1} El repartidor de avisos recibe por RabbitMQ los eventos de kits publicados, pago, incidencia, despacho, reparto y entrega, y no repite el aviso de un evento ya enviado.\\[1pt]
  \textbf{C2} Pide a identidad el contacto del acudiente, sin guardarlo.\\[1pt]
  \textbf{C3} Si el acudiente tiene celular, arma un enlace de WhatsApp con el aviso; si no, se lo envía por correo.\\[1pt]
  El despacho lo registra el equipo en el panel, con el nombre y el celular del repartidor propio; las transportadoras con guía son una evolución prevista.};
\node[anchor=north west, text width=15.1cm, align=left, inner sep=0pt, font=\scriptsize]
  at ([yshift=-2mm]colA.south west |- current bounding box.south) {Línea continua: llamada HTTP. Línea discontinua:
  el publicador lee la bandeja de salida de su base. Línea violeta: evento que viaja por RabbitMQ hasta la cola de su
  consumidor (avisos, colegios, listas y trabajos), con reintentos y descartados; sin RabbitMQ, los publicadores no
  arrancan y cada consumidor lee la bandeja del productor. Línea punteada: componente opcional. Línea sin punta: enlace que el sistema
  arma y no envía. Recuadro gris: el servicio con sus procesos y su base.};
\end{tikzpicture}

\caption{Vista de contenedores (C4 nivel 2) del sistema construido, con los procesos en
segundo plano de cada servicio y los pasos de la subida de la lista y los kits (A), el pago (B)
y los avisos (C).}
\label{fig:c4-contenedores}
\end{figure}

Redis es opcional y solo guarda lo que se puede recalcular o perder (ADR 0038): el vector de cada
renglón durante 24 horas, la lectura aceptada de cada lista durante 30 días y los contadores del
límite de tasa de la puerta de enlace. Si no responde en 250~ms, cada servicio calcula sin caché
y la puerta de enlace cuenta el límite en su memoria.

% =====================================================================
\subsection{Modelo de dominio y contextos delimitados}
\label{sec:dominio}

La descomposición sigue el enfoque de diseño guiado por el dominio \citep{evans_domain-driven_2004}.
Cada servicio corresponde a un contexto delimitado con un lenguaje ubicuo propio, y las
traducciones entre contextos ocurren explícitamente en los adaptadores.

Las migraciones del sistema construido crean 74 tablas en cinco bases de datos, una por
servicio (Tabla~\ref{tab:reparto}); la bandeja de salida de eventos se repite en las cinco,
de modo que hay 70 nombres distintos. El recuento cambió durante la implementación, entre
67 y 74 tablas según la versión.

% Tabla tab:reparto movida al Anexo A.3 (tablas del cuerpo)


\subsubsection{La matrícula como ancla de tenencia}

El modelo de datos de los microservicios anteriores no sabía a qué colegio pertenecían los
pedidos ni los usuarios: el identificador de institución aparecía cero veces en los
servicios de comercio y de cuentas. Donde sí existía, lo enviaba el cliente en el cuerpo de
la petición sin validación alguna, lo cual es peor que su ausencia porque aparenta un
control que no existe.

La solución aparentemente obvia, un identificador de colegio en el usuario, es incorrecta,
porque un acudiente puede tener hijos en colegios distintos (CO-04). La tenencia reside, en
cambio, en la relación $\langle$estudiante, sede, grado, año lectivo$\rangle$, que se modela
como la entidad matrícula, con una relación de muchos a muchos con los acudientes
(ADR 0012). Su identificador viaja a los otros cuatro servicios como referencia lógica y es
el eje de las políticas de visibilidad. Se ubica en la \emph{línea} del pedido y no en su
cabecera, porque en la cabecera dos hermanos obligarían a dos pedidos y a dos pagos, lo que
contradice la propuesta de valor del producto. La comprobación de que cada línea apunta a
una matrícula visible para el acudiente se delimitó como parte del endurecimiento productivo
(Sección~\ref{sec:trabajo-futuro}).

\subsubsection{Asimetría entre lectura y escritura}

Las políticas de visibilidad son deliberadamente asimétricas. Para \emph{leer} un recurso
basta que toque alguna matrícula visible para el solicitante; para \emph{escribirlo} se
exige que \emph{todas} las matrículas que el recurso referencia estén contenidas en su
ámbito. Un pedido con dos hermanos, uno visible y otro no, puede consultarse pero no
modificarse. La simetría ingenua produciría, en ese escenario, una fuga de información o un
sistema inutilizable.

% =====================================================================
\subsection{Autorización con ámbito y seguridad a nivel de campo}
\label{sec:autorizacion}

El modelo de autorización de los microservicios anteriores representaba el rol como una
cadena de texto en la tabla de usuarios, y el control de acceso se reducía a verificar que
esa cadena perteneciera a un conjunto permitido. Dos usuarios con el mismo rol sobre
colegios distintos superaban, así, exactamente el mismo control. La única defensa consistía
en que cada punto de entrada recordara comparar el colegio del token con el del recurso
solicitado, y la auditoría de los 152 puntos de entrada de esos servicios encontró 22 en los
que esa comparación faltaba o era parcial (Sección~\ref{sec:defectos}).

El diagnóstico requiere precisión. Uno de esos servicios, el de instituciones, ya definía
una relación de pertenencia con la tripla \emph{(usuario, institución, rol)}, introducida
para corregir una referencia directa insegura real. Expresar autoridad con ámbito no era,
por tanto, imposible: el defecto era de alcance y de completitud. El mecanismo cubría un
servicio de siete; admitía un único nivel de ámbito, la institución, cuando el recurso
protegido pertenece a una sede y un grado; permitía un solo rol por institución, y carecía
de vigencia temporal. Esta última carencia producía una contradicción ilustrativa: la
documentación del modelo justificaba conservar el vínculo para auditar quién tuvo acceso y
cuándo, pero la única forma de revocarlo era eliminarlo. El diseño generaliza esa solución,
ya validada contra un ataque real, con la jerarquía de ámbitos, la vigencia temporal y la
aplicación uniforme a todos los servicios.

Con el control de acceso basado en roles con ámbito, la pregunta que el sistema responde
deja de ser \emph{qué es este usuario} y pasa a ser \emph{qué puede hacer este usuario sobre
este recurso concreto}:

\begin{equation}
\mathrm{puede}(u, p, a) \iff \exists\, r \in \mathcal{R}_{\text{vig}}(u, a) : (r, p) \in \mathcal{P}
\end{equation}

\noindent donde $\mathcal{P}$ es el conjunto de pares de rol y permiso, y
$\mathcal{R}_{\text{vig}}(u, a)$ el de roles vigentes del usuario $u$ sobre el ámbito $a$ o
sobre cualquier ámbito que lo contenga, según la jerarquía
$\text{global} \succ \text{colegio} \succ \text{sede} \succ \text{matrícula}$. La contención
opera en un solo sentido: un rol otorgado sobre un colegio alcanza sus sedes, pero ser
acudiente de una matrícula no confiere autoridad sobre la sede.

\subsubsection{La función falla cerrada}

Las políticas de seguridad a nivel de fila de los microservicios anteriores iniciaban su
condición base verificando si el contexto de usuario era nulo, de modo que toda conexión sin
contexto establecido veía la tabla completa. El control fallaba \emph{abierto}, y su modo de
fallo más probable, omitir el contexto, era precisamente el que concedía acceso total. La
función aquí especificada no tiene ninguna rama que conceda por ausencia de información.

En el sistema construido, identidad guarda las asignaciones de rol con su vigencia y los
pares de rol y permiso, y evalúa los permisos con ámbito (ADR 0012); las rutas protegidas
de inteligencia fallan cerradas (SG-06). La aplicación uniforme (AD-03 y SG-02) se delimitó como
endurecimiento productivo (Sección~\ref{sec:trabajo-futuro}): solo una ruta de producción consulta la función de permiso, y varias rutas que
devuelven datos personales no exigen sesión ni rol. Tampoco hay políticas de seguridad a
nivel de fila, de modo que la evaluación local del ADR 0015 sigue siendo diseño
(Sección~\ref{sec:defectos}).

\subsubsection{Dos ejes independientes, no uno}

La seguridad a nivel de campo usa el mismo mecanismo, con un componente de campo en el
permiso: leer un estudiante y leer su nombre son permisos distintos. El rediseño inicial
afirmó que esos permisos se derivaban \emph{automáticamente} del metadato de clasificación de
datos, y el análisis del registro de tratamiento mostró que no es así. La clasificación mide
sensibilidad y el permiso expresa quién puede ver qué: son dos ejes independientes. De la
clasificación solo se deriva el valor por omisión, de modo que todo campo clasificado como
dato de menor o como personal sensible nace denegado y todo campo público nace permitido.
Colapsar ambos ejes conduce a implementar la mitad del mecanismo creyendo que está completo.

% =====================================================================
\subsection{Protección de datos personales desde el diseño}
\label{sec:datos}

El sistema trata datos de menores de edad, datos personales de los acudientes, como el
contacto y la dirección de entrega, y datos de las instituciones educativas. La Ley 1581 de
2012 \citep{congreso_ley1581_2012} rige los dos primeros, y los terceros cuando identifican
a personas naturales: exige autorización previa y una finalidad legítima informada al
titular, y reconoce el derecho de supresión.

Se adoptó un esquema de cinco niveles de clasificación (público, interno, personal,
personal sensible y dato de menor), aplicado campo por campo en el registro de tratamiento
(Anexo~M), que cubre las cinco bases del sistema construido. Su primera versión, sobre los
microservicios anteriores, revisó 398 columnas de 50 tablas y encontró 126 campos con datos
personales, 20 de ellos clasificados como dato de menor.

\subsubsection{Políticas que el esquema anterior no podía sostener}

El ejercicio reveló un patrón recurrente: políticas correctamente enunciadas en la
documentación que el esquema hacía inejecutables. La Tabla~\ref{tab:politicas} recoge los
casos y su corrección.

% Tabla tab:politicas movida al Anexo A.3 (tablas del cuerpo)


El patrón tiene valor metodológico más allá de este proyecto: una política de tratamiento
de datos es verificable solo si el modelo lógico expone las columnas que la hacen
computable. Documentar la política sin instrumentar el esquema produce cumplimiento
aparente. El riesgo inverso también existe: el sistema construido tiene esas columnas, pero
la supresión propagada existe en el modelo y no como flujo de extremo a extremo
(Sección~\ref{sec:defectos}). El SRS v6.1 deja la supresión propagada, el panel de privacidad
y la retención automática fuera del alcance de este trabajo, y atiende los derechos del
titular mediante un caso de soporte.

\subsubsection{Minimización por columna: la dirección de entrega}

Los microservicios anteriores no tenían ninguna dirección de entrega, en un negocio cuya
propuesta de valor consiste en entregar el kit en el domicilio. Tampoco figuraba en el
esquema objetivo, de modo que la omisión se arrastraba del código al diseño: un defecto por
omisión sobrevive más que uno por comisión, porque nada lo señala.

El dato es sensible de forma poco habitual: más que información de contacto, es la
ubicación física del hogar de una familia con menores, cruzada con el colegio al que
asisten. El ADR 0014 distingue dos cosas que se parecen sin ser la misma: la libreta de
direcciones reutilizables del acudiente, en identidad, editable y sujeta a supresión, y la
copia congelada en el envío, en comercio, inmutable una vez confirmado el pedido. El
razonamiento es análogo al de congelar el precio en la línea del pedido: un hecho histórico
y un dato editable son entidades distintas aunque se escriban igual. Si el envío
referenciara la libreta, corregir un número de apartamento reescribiría el destino de una
entrega pasada.

La partición de columnas responde al principio de minimización. La ciudad y la localidad
quedan en claro porque de ellas depende la asignación de la zona de reparto, mientras que
la calle, el número y el teléfono se cifran en la aplicación: planificar una ruta requiere
conocer la zona, no el apartamento. En el sistema construido esa partición rige para la
libreta, que guarda cifrados la dirección y el contacto. La copia del envío conserva el
destinatario y la dirección en claro, como prueba ante un reclamo de la Ley 1480 de 2011, de
modo que la minimización por columna no cubre todavía el dato congelado.

% =====================================================================
\subsection{Organización del código y funciones de aptitud}
\label{sec:monorepo}

\subsubsection{Un repositorio no es un desplegable}

La implementación se organiza en un monorrepositorio con seis desplegables independientes
(ADR 0013). La decisión exige justificación porque contradice la asociación habitual entre
microservicio y repositorio, asociación que carece de fundamento normativo. La norma
ISO/IEC/IEEE 42010 \citep{iso_42010_2022} compone una descripción arquitectónica de vistas
que responden a preocupaciones distintas, y el modelo 4+1 \citep{kruchten_41_1995} separa
la \emph{vista de desarrollo}, que organiza módulos y bibliotecas, de la \emph{vista
física}, que asigna el software a nodos de ejecución. Dónde vive el código y cómo se
despliega son decisiones de vistas diferentes; en la terminología del modelo C4, un
contenedor no equivale a un repositorio.

La motivación concreta es medible. En los microservicios anteriores, trasladar una función
de redacción de datos a la biblioteca compartida exigió modificar siete archivos de
configuración de dependencias y publicar una etiqueta de versión antes de que cualquier
servicio pudiera instalarla. En el monorrepositorio es un único cambio.

\subsubsection{Fronteras verificadas, no declaradas}

La objeción legítima al repositorio único es que facilita la erosión de las fronteras. La
respuesta de este diseño es instrumentarlas mediante funciones de aptitud
\citep{ford_building_2017}: pruebas automáticas y continuas sobre una característica
arquitectónica. Un manifiesto declara, para cada servicio, su lenguaje, su base de datos, las tablas de las que es
dueño y las referencias que cruzan hacia otros dominios, y doce pruebas lo verifican en cada
cambio: que ningún servicio importe código de otro ni declare una tabla ajena, que toda referencia
cruzada esté declarada, que ninguna tabla tenga dos dueños y que el grafo de dependencias sea
acíclico. Dos pruebas más, añadidas con las colas, y una con Redis, extienden la verificación a
esos canales. La lista completa de las doce pruebas está en el Anexo~\ref{anx:ampliaciones}.
El motivo de instrumentarlas es que el defecto que encontró el rediseño, la tabla de productos
declarada como maestra en dos servicios a la vez, no se produjo por falta de criterio sino porque
nada lo verificaba de forma continua.

Una función de aptitud que nunca ha fallado no está demostrada, pues podría no estar
midiendo nada. Por ello el conjunto incluye un segundo nivel, con veintiocho pruebas, que
suministra a los detectores infracciones plantadas y verifica que las detecten, además de
casos que no deben marcarse. Entre esas infracciones están una tabla ajena declarada como
el doble maestro y el ciclo de dependencias que existió. Ante una infracción plantada, la
integración falla e indica el archivo, la línea, el dominio invadido y la remediación
sugerida.

% =====================================================================
\subsection{Asignación de lenguajes: una decisión medida}
\label{sec:lenguajes}

La hipótesis de partida sostenía que los componentes intensivos en cómputo debían reescribirse en
un lenguaje compilado. Para contrastarla se perfiló una lista escolar de veinte ítems, la mediana
del corpus real, contra un catálogo de 3.187 productos. Las cifras son reportadas por la
documentación del código, no recalculadas. El 95,3\,\% de la latencia lo producía una sola línea,
que reconstruía un grafo de sustituibilidad en cada ítem; construirlo cuesta 1.673~ms y consultarlo,
0,137~ms. Corregido ese defecto, la reducción que aporta cambiar de lenguaje queda, según la ley de
Amdahl \citep{amdahl_validity_1967}, entre el 20,0\,\% y el 21,7\,\%, y el 98,5\,\% de lo que
resta corresponde al optimizador NSGA-II sobre pymoo \citep{blank_pymoo_2020}, sin equivalente
maduro en un lenguaje compilado. Por eso los seis desplegables están en Python, y el manifiesto
declara el lenguaje de cada servicio para que incorporar otro sea una adición. El cálculo
completo y el desglose de latencia (Tabla~\ref{tab:latencia}) están en el
Anexo~\ref{anx:ampliaciones}.

% =====================================================================
\subsection{Decisiones registradas}
\label{sec:adr}

Las decisiones arquitectónicas se documentan como registros en el formato propuesto por
Nygard \citep{nygard_documenting_2011} (título, contexto, decisión, estado y consecuencias),
al que el proyecto añade un apartado de alternativas consideradas. La entrega reúne
treinta y ocho registros, numerados de 0001 a 0038, en el registro de decisiones de
arquitectura (Tabla~\ref{tab:anexos}); la Tabla~\ref{tab:adr} destaca los que estructuran el
diseño de este capítulo.

El ADR 0015 ilustra una propiedad del método. Nació de una contradicción entre decisiones previas
que solo se hizo visible al redactar el criterio de aceptación de un requisito: la política de
seguridad no podía evaluarse en ninguno de los tres servicios donde hacía falta, porque la
matrícula vive en la base de identidad. Un registro obliga a escribir el fundamento donde puede
contrastarse con otro.

% Tabla tab:adr movida al Anexo A.3 (tablas del cuerpo)


% =====================================================================
\subsection{Síntesis}

El diseño presentado se distingue del original en tres aspectos, que constituyen el aporte
de ingeniería de este objetivo específico. El primero es la justificación de la
descomposición, que pasó de la escalabilidad, refutada por el cálculo de carga, a la
recuperabilidad asimétrica, la frontera de cumplimiento y el perfil de recurso. Con ella,
los catorce servicios documentados se redujeron a seis desplegables con código, cuyas
fronteras se trazaron de nuevo. El segundo es la matrícula, una entidad que el modelo
anterior no tenía y que permite expresar tenencia, autoridad con ámbito y revocación
parcial del consentimiento. El tercero es que las fronteras entre servicios dejaron de ser
afirmaciones y pasaron a ser aserciones verificables de forma continua. La autorización
uniforme y la depuración del nombre del menor en el documento cargado se delimitaron como
endurecimiento productivo (Sección~\ref{sec:trabajo-futuro}).

    \setcounter{subsection}{0}
\section{DESARROLLO DE LA SOLUCIÓN}
\label{cap:desarrolloSolucion}

Este capítulo documenta la preparación de los datos y el modelado de CRISP-DM en las Fases~3
y~4, ligadas a los objetivos específicos tercero y cuarto: cómo se construyó el mecanismo que
convierte una lista escolar fotografiada en un kit de útiles concreto. El capítulo de diseño
(Sección~\ref{sec:impulsores} y siguientes) estableció dónde viven los componentes y por qué las
fronteras entre servicios se trazan así. Este explica cómo se modeló el comportamiento dentro de
esas fronteras, qué alternativas se descartaron y qué se reemplazó por una opción mejor durante la
construcción.

% =====================================================================
\subsection{El pipeline de resolución}
\label{sec:pipeline}

El problema se descompone en cuatro etapas secuenciales, y cada una consume la salida de la
anterior. La percepción extrae los renglones de una imagen o un PDF sin estructura. La
normalización recupera del catálogo los productos candidatos de cada renglón. La resolución decide
cuál corresponde o marca el renglón como ambiguo, y la optimización arma con los productos
resueltos los kits que pide el acudiente, que equilibran costo, fidelidad de marca, disponibilidad y personaje. La
percepción, la resolución y la optimización se alojan en el servicio de inteligencia, cuya carga
intensiva en CPU y memoria motivó su separación en el diseño. La recuperación de candidatos la
ejecuta el catálogo sobre su propia base.

La orquestación de las cuatro etapas se decidió en el ADR 0017. El pipeline se modela como trabajos
durables con estado persistido, reanudables por etapa y con una espera explícita para la revisión
humana. El registro consideró LangGraph, con persistencia de estado sobre PostgreSQL, y lo difirió:
los trabajos durables ya dan reanudación, reintentos idempotentes y auditoría con dependencias que
el equipo controla, y adoptarlo después no cambiaría el modelo de datos. Descartó también un motor
externo como Temporal o Airflow, desproporcionado para un pipeline que cabe en dos tablas, y dejó
LangGraph como camino de crecimiento si el pipeline llega a tener ramas.

El ADR 0017 preveía una biblioteca de máquinas de estados y una clave de reanudación
formada por el trabajo, la etapa y el intento. El sistema construido usa una tabla de transiciones
propia, con seis estados del trabajo, y registra cada etapa una sola vez por trabajo. El sexto,
cancelado, es terminal y se alcanza cuando el acudiente cancela su solicitud. Según el SAD,
la solución es más simple y basta para reintentar: un reintento vuelve a encolar el trabajo y
recalcula sus ítems y sus kits sin que el acudiente suba otra vez el documento ni se pague de nuevo
la percepción. Su costo es que no conserva el historial de intentos de cada etapa.

% =====================================================================
\subsection{De los prototipos al sistema construido}
\label{sec:lenguajes-dev}

La solución llegó a su forma actual por iteraciones sucesivas, cada una en su propio proyecto de
código. La prueba de concepto, del 29 de junio al 6 de septiembre de 2026, reunió por primera vez
el pipeline completo: lectura con Tesseract y PyMuPDF, raspado de tiendas, búsqueda en un índice
FAISS y un optimizador NSGA-II de tres objetivos. El 30 de julio la búsqueda pasó a PostgreSQL, con
pgvector y trigramas fusionados por rangos. Entre el 31 de julio y el 6 de agosto, un banco de
alternativas comparó almacenes vectoriales y la recuperación con y sin un modelo de lenguaje, y dejó
planteada la tarea como una correspondencia con el catálogo, sin generación de texto.

Del 29 de agosto al 21 de septiembre, los prototipos de tesis sometieron cada etapa a experimentos
con muestra y métrica: percepción, acuerdo de anotación, embeddings, bases de datos y calibración
del umbral. El 14 de septiembre, un prototipo de calibración reunió esos resultados. En paralelo, del 2 al
18 de septiembre, el producto se organizó como una flota de microservicios con librerías comunes. El
servicio de búsqueda, del 7 al 17 de septiembre, llevó la recuperación a un servicio propio, y los
frontales, del 6 al 19 de septiembre, añadieron la revisión de renglones y los perfiles de kit.

El sistema construido, desde el 22 de septiembre y hasta el 24 de octubre,
reunió ese trabajo en un monorrepositorio con
despliegue independiente por servicio (ADR 0013). Tiene seis desplegables con código, cinco
servicios de dominio y la puerta de enlace, sobre cinco bases de datos; el séptimo previsto, la
ingesta, se reemplazó por el catálogo curado desde la demanda de las listas
(Sección~\ref{sec:consolidacion}). Las cuatro etapas se implementaron en Python, incluida la
optimización con pymoo \citep{blank_pymoo_2020}, porque el cálculo de Amdahl
\citep{amdahl_validity_1967} del capítulo de diseño acotó en un 21,7\,\% como máximo la reducción
de latencia de reescribirlas en un lenguaje compilado (Sección~\ref{sec:lenguajes}). La migración
empezó por identidad, que concentraba la frontera de cumplimiento más urgente, y por catálogo, que
tenía el defecto de doble maestro (Sección~\ref{sec:monorepo}). Durante ella, las funciones de
aptitud detectaban si el código portado reintroducía una referencia directa a otro dominio
\citep{ford_building_2017}. La medición del umbral de resolución del 4 de octubre pertenece a una
línea de desarrollo que no está integrada en el sistema construido, y sus cifras se citan como
reportadas. El afinamiento de los kits, del 11 al 15 de octubre, se integró el 22 de octubre.
Entre el 19 y el 23 de octubre se añadieron la recuperación ante las caídas del catálogo y de
los procesos de fondo, las sugerencias de dirección, el alta de colegios no aliados y las colas
por RabbitMQ (ADR 0031 a 0037). El 24 de octubre se sumaron la caché y el límite de tasa con
Redis (ADR 0038).

La Tabla~\ref{tab:rag-fases} resume, por fase de CRISP-DM, el resultado de esas iteraciones para
las 84 estrategias de recuperación: las implementadas en el sistema construido, las que quedaron en
prototipos, las evaluadas y descartadas y las que no se consideraron, no se evaluaron o no aplican
al dominio. El anexo CRISP-DM, listado en la Tabla~\ref{tab:anexos}, describe cada estrategia
dentro de su fase, con su evidencia, sus resultados y su justificación. Dos rasgos dominan.
Veinte estrategias existen solo en prototipos, de modo que sus cifras no describen el sistema
construido. Y la última etapa no genera texto libre: el kit es una selección de productos del
catálogo validada por esquema, aunque la base no impone esa pertenencia con una clave foránea,
porque la línea del kit y el catálogo están en bases distintas.

\begin{table}[htbp]
\centering
\caption{Estrategias de recuperación por fase de CRISP-DM y estado en el sistema construido}
\label{tab:rag-fases}
\footnotesize
\begin{tabular}{@{}lrrrrp{6.2cm}@{}}
\toprule
\textbf{Fase} & \textbf{Impl.} & \textbf{Prot.} & \textbf{Desc.} & \textbf{Otras} & \textbf{Lo que usa el sistema construido} \\
\midrule
Comprensión del negocio & 1 & 0 & 1 & 5 & Sugerencias ancladas al catálogo, sin generación de texto \\
Comprensión de los datos & 2 & 1 & 0 & 0 & Catálogo curado desde la demanda; 175 listas de 17 colegios \\
Preparación de los datos & 13 & 6 & 1 & 4 & Metadatos y niveles del catálogo; percepción con Gemini validada por esquema; limpieza, corrección, sinónimos y atributos del renglón \\
Modelado & 11 & 9 & 5 & 10 & gemini-embedding-001; fusión aditiva del puntaje léxico y el coseno; filtros de metadatos; hasta tres opciones, una por nivel \\
Evaluación & 1 & 3 & 0 & 0 & $\tau=0{,}75$ con abstención y revisión de la lista completa \\
Despliegue & 6 & 1 & 1 & 3 & pgvector en la base del catálogo; caché de vectores en Redis; índice de conceptos en memoria; trabajos durables con espera humana \\
\midrule
Total & 34 & 20 & 8 & 22 & \\
\bottomrule
\end{tabular}
\par\smallskip
{\footnotesize Impl.: implementada; Prot.: solo en prototipos; Desc.: evaluada y descartada; Otras: considerada sin evaluar, no considerada o no aplica al dominio.}
\end{table}


% =====================================================================
\subsection{Percepción documental}
\label{sec:percepcion-dev}

La primera etapa recibe una fotografía o un PDF sin estructura garantizada y debe producir los
renglones con su cantidad y su descripción. El 1 de septiembre de 2026, los prototipos de tesis
compararon cinco enfoques: un flujo clásico de reconocimiento óptico de caracteres (PyMuPDF,
Tesseract y un analizador de expresiones regulares), Google Vision, GOT-OCR2.0, Docling y el modelo
multimodal Gemini. Se evaluaron con F1 sobre un conjunto de oro propio, en la línea de los bancos de
OCR para modelos multimodales \citep{liu_ocrbench_2024} y de extracción clave-valor
\citep{shen_ocr_2026}. La razón para preferir un modelo multimodal es estructural. El
reconocimiento con reglas depende de que el documento tenga una disposición regular de renglones, y
buena parte del corpus del objetivo primero son fotografías de listas manuscritas o impresas,
tomadas con el teléfono del acudiente y sin luz controlada. Un modelo que generaliza sobre la
disposición visual, en vez de suponer una, era la única familia compatible con esa variabilidad.
Las cifras están en la Sección~\ref{sec:percepcion}; la elección no tiene un ADR propio.

En el sistema construido, cada página se convierte en imagen y se envía a Gemini con un prompt en
español y el esquema de la respuesta, con salida JSON y temperatura 0. Si la respuesta no es JSON o
no cumple el esquema, el trabajo falla con su motivo: no hay reintento escondido ni renglón
inventado. Solo la falla del transporte tiene un reintento, declarado y único.

El nombre del menor es un dato
protegido por la Ley 1581 de 2012 \citep{congreso_ley1581_2012}, y la propuesta inicial era
anonimizarlo antes de guardar la lista y antes de invocar el modelo externo; los prototipos lo
intentaron con expresiones regulares y con un reconocedor de entidades. El sistema construido lo
protege de otra forma, por minimización: no lo recoge. Según PV-03 del SRS, no pide el nombre legal del menor, el
apodo con que el acudiente lo identifica es opcional y se cifra en reposo, y la percepción extrae
solo los ítems y el grado. Es mejor porque un dato que no se recoge no puede llegar a los registros
de operación, a los conjuntos de evaluación ni al modelo, y no depende de que un reconocedor lo
detecte. Queda un caso: si el nombre figura en el documento cargado, viaja con la página hacia el
proveedor, y si el modelo lo lee como renglón, queda 30 días en la lectura que guarda Redis. Su
depuración, que pide PE-06, se delimitó como parte del endurecimiento productivo previo al
despliegue con familias (Secciones~\ref{sec:defectos} y~\ref{sec:trabajo-futuro}), porque la
evaluación del trabajo usó listas con estudiantes ficticios.

% =====================================================================
\subsection{Normalización y resolución}
\label{sec:normalizacion}

Cada renglón extraído se busca en el catálogo con recuperación híbrida: una señal léxica, que
captura coincidencias exactas de marca y referencia, y una señal densa, que captura equivalencias de
significado entre descripciones redactadas de forma distinta
\citep{reimers_sentence-bert_2019,karpukhin_dense_2020}. Del marco de generación aumentada por
recuperación \citep{lewis_retrieval-augmented_2020}, esta etapa toma solo el componente de
recuperación. Antes de buscar, el texto se limpia, se
corrige la ortografía con el léxico del catálogo, se reduce al núcleo en singular y se separan los
atributos impuestos por la lista, como el color, el tamaño o la marca; los sinónimos viven en cada
concepto del catálogo.

La
prueba de concepto guardó los vectores en un índice FAISS \citep{johnson_billion-scale_2021} hasta
el 30 de julio. El sistema construido guarda el vector de 1024 dimensiones de cada producto en el
PostgreSQL del catálogo, con pgvector (Sección~\ref{sec:dominio}), junto con el modelo y el texto
del que se calculó. Así el vector se respalda con el producto, sin un segundo almacén
que sincronizar \citep{kleppmann_designing_2017}, y una huella del texto indica cuándo recalcularlo.
Los archivos versionados del catálogo no lo llevan: un proceso por lotes lo calcula contra la base.
La búsqueda se hace en memoria, con el coseno calculado en el proceso: los índices HNSW y de
trigramas existen, pero no se consultan. El SAD advierte que harán falta antes de llegar a los
50.000 productos del escenario ESC-E4.

La comparación con MTEB
\citep{muennighoff_mteb_2023} eligió e5-large, de 560 millones de parámetros, que habría que alojar
y servir. El sistema construido usa gemini-embedding-001, del mismo proveedor que la percepción. La
ventaja es operativa (no hay modelo que alojar ni GPU) y no una mejora de calidad demostrada, porque
los dos no se compararon sobre el corpus. El vector del renglón se calcula en cada solicitud y el de
cada concepto es el promedio de sus productos. Si el servicio de embeddings falla, la búsqueda sigue
solo con la señal léxica y el resultado lo advierte.

Los prototipos
fusionaban las dos señales por rangos (RRF) y reordenaban con un codificador cruzado. Sobre los
2.710 pares del conjunto de calibración, el área bajo la curva fue de 0,728 para la señal léxica,
0,805 para la densa, 0,739 para RRF y 0,880 para el reordenador. El sistema construido resuelve en
una sola fase: un puntaje léxico anclado en el núcleo del renglón más 0,15 veces el coseno, con un
rescate por coseno de 0,86 o más cuando el puntaje léxico no supera 0,55. A cambio, el
administrador de Tukit revisa completa cada lista nueva antes de publicarla (PE-04), en cuatro horas
o menos en temporada (ML-09). El SDD da tres razones. No depende de un reordenador ni de una GPU;
corrige el error que hacía tomar «chaquiras» por otro producto; y en la corrida del servicio de
búsqueda anterior, 1.407 de 2.176 ítems (64,7\,\%) necesitaban de todos modos a una persona. La
revisión se hace una vez por colegio y grado y se reutiliza (NR-07). Las constantes de la fusión (0,15; 0,86; 0,55) proceden de ajustes sobre casos de error
concretos, como el de las «chaquiras», y no de un barrido; la Tabla~\ref{tab:constantes} registra el
origen y el estado de cada una.

La estructura del catálogo hace el resto. Sus 149 conceptos, 145 de ellos activos, agrupan los
productos equivalentes; en el levantamiento de listas, 25 conceptos cubren el 80,97\,\% de los
renglones con concepto. Los atributos y los niveles distinguen
las variantes, y el sistema se abstiene cuando el producto contradice un atributo impuesto. Cada
atributo se compara a su manera (exacto, como mínimo, con tolerancia o por equivalencias), y el
que el producto no declara no cumple (ADR 0034). Un grafo
de sustitución entre productos se evaluó y se descartó: subía la cobertura del 41,3\,\% al
65,3\,\%, pero bajaba la precisión del 86,7\,\% al 83,5\,\%. La tabla de sustitución dirigida entre
conceptos curados no tiene datos y no se consulta (Sección~\ref{sec:consolidacion}).

\subsubsection{Umbral de resolución}

El umbral pasó por tres valores. El servicio de búsqueda anterior fijó $\tau=0{,}85$ y un margen
$\delta=0{,}05$ con siete ejemplos sintéticos. El prototipo de calibración lo situó en $\tau=0{,}630$
sobre el puntaje del reordenador, con 542 ítems, y así quedó en el ADR 0009. El sistema construido
usa $\tau=0{,}75$ sobre su puntaje propio, sin margen, como valor inicial de un ciclo de calibración
con revisión humana (Sección~\ref{sec:ciclos-calibracion}). Su efecto está contenido porque el umbral
solo señala renglones dentro de una lista que de todos modos se revisa completa, de modo que un
$\tau$ mal ajustado cuesta tiempo de revisión y no un producto equivocado en el carrito. El valor
vigente y su evidencia se registran en la Tabla~\ref{tab:constantes}\pendienteEvidencia{valor final de tau sobre el puntaje del sistema construido}.

\subsubsection{Algoritmo de resolución de un renglón}
\label{sec:algoritmo-resolucion}

La resolución es el algoritmo propio de OE3, y se presenta con la estructura del curso de
algoritmos del departamento: análisis del problema, diseño de entradas y salidas, estrategia,
seudocódigo y complejidad con comparación frente al método exhaustivo.

\paragraph{Análisis del problema.} Dado un renglón de la lista, por ejemplo «pompones de colores
x 20», hay que decidir a cuál de los conceptos canónicos del catálogo corresponde, es decir, a qué
clase de producto equivalente se refiere, o declarar que no corresponde a ninguno. Dos errores
motivaron el diseño. El primero es el de las «chaquiras de colores», que una regla de contención
literal asignaba al concepto «colores» (lápices de colores) porque la palabra aparecía en ambos. El
segundo es el de los sinónimos que ninguna lista contiene, como «estuche» por «cartuchera». La
solución tenía que corregir el primero sin perder el segundo y sin que la señal semántica pudiera
empeorar un acierto léxico.

\paragraph{Diseño: entradas y salidas.} La Tabla~\ref{tab:es-resolucion} define los tipos.

\begin{table}[H]
\centering\footnotesize
\caption{Entradas y salidas del algoritmo de resolución de un renglón.}
\label{tab:es-resolucion}
\begin{tabularx}{\textwidth}{@{}p{2.6cm}p{3.6cm}X@{}}
\toprule
Elemento & Tipo & Descripción \\
\midrule
Entrada $r$ & cadena de texto & Renglón de la lista, con cantidad y empaque \\
Entrada $v_r$ & vector real de 1.024 componentes, opcional & Embedding del renglón; si falta, el algoritmo usa solo la señal léxica \\
Entrada $\mathcal{C}$ & conjunto de conceptos, cada uno con nombre, lista de sinónimos y vector $v_c$ & 149 conceptos, 145 activos \\
Parámetros & tres números reales en $[0,1]$ & Aporte del coseno $\gamma=0{,}15$, umbral de rescate $\theta=0{,}86$ y piso del rescate $\phi=0{,}55$ \\
Salida & terna (concepto o nulo, nombre, puntaje en $[0,1]$) & Mejor concepto y su puntaje; un puntaje menor que $\tau$ remite el renglón a revisión \\
\bottomrule
\end{tabularx}
\end{table}

\paragraph{Estrategia algorítmica.} El método es una búsqueda del máximo puntaje sobre los
conceptos, con un puntaje en tres ramas que se evalúan en orden. Primero se extrae el núcleo del
renglón, que es la primera palabra útil después de quitar cantidades y empaques, porque en español
el núcleo del sintagma va primero. Si el núcleo está en la cabeza del nombre o de un sinónimo del
concepto, el puntaje es $0{,}6+0{,}4\,J+\gamma\cos$, donde $J$ es el parecido de Jaccard entre las
palabras significativas y $\cos$ es el coseno entre $v_r$ y $v_c$, recortado en 0. Si el núcleo
aparece en otra posición, el puntaje es $J+\gamma\cos$. Si no aparece y el coseno alcanza $\theta$, el
rescate asigna $\phi+0{,}4(\cos-\theta)/(1-\theta)$, que parte del piso $\phi$, apenas sobre el corte de revisión de
0,5. En cualquier otro caso el puntaje es 0. Como $\gamma\cos\ge 0$, el puntaje híbrido nunca es
menor que el léxico, que es la garantía que la tesis afirma, y el umbral alto de rescate impide que
un parecido de tema, como el de «chaquiras» con «colores», se cuele.

\begin{algorithm}[H]
\caption{Resolución de un renglón contra el catálogo de conceptos}
\label{alg:resolucion}
\KwEntrada{renglón $r$, vector $v_r$ (opcional), conceptos $\mathcal{C}$, parámetros $\gamma,\theta,\phi$}
\KwSalida{concepto $c^*$, puntaje $p^*$}
$n \leftarrow$ núcleo$(r)$; $(c^*, p^*) \leftarrow (\text{nulo}, 0)$\;
\ForEach{concepto $c \in \mathcal{C}$}{
  $\cos \leftarrow \max(0, \text{coseno}(v_r, v_c))$ si hay $v_r$ y $v_c$; en otro caso $0$\;
  \ForEach{candidato $k \in \{\text{nombre}(c)\} \cup \text{sinónimos}(c)$}{
    \uIf{$n$ está en la cabeza de $k$}{ $p \leftarrow \min(1,\ 0{,}6 + 0{,}4\,J(r,k) + \gamma\cos)$ }
    \uElseIf{$n$ está en alguna palabra de $k$}{ $p \leftarrow \min(1,\ J(r,k) + \gamma\cos)$ }
    \uElseIf{$\cos \ge \theta$}{ $p \leftarrow \phi + 0{,}4\,(\cos-\theta)/(1-\theta)$ }
    \Else{ $p \leftarrow 0$ }
    \lIf{$p > p^*$}{ $(c^*, p^*) \leftarrow (c, p)$ }
  }
}
\Return{$(c^*, p^*)$}\;
\end{algorithm}

\paragraph{Complejidad y comparación con el método exhaustivo.} Sea $|\mathcal{C}|$ el número de
conceptos, $s$ el número medio de candidatos por concepto (nombre y sinónimos), $w$ el número de
palabras de un candidato y $d=1.024$ la dimensión del vector. El coseno se calcula una vez por
concepto, con costo $O(|\mathcal{C}|\,d)$; la comparación léxica cuesta $O(w^2)$ por candidato con
conjuntos pequeños. El total por renglón es $O(|\mathcal{C}|\,(d+s\,w^2))$, lineal en el
tamaño del catálogo de conceptos. Con 149 conceptos y $d=1.024$, el coseno requiere unas 152.000
multiplicaciones por renglón (149 por 1.024). El método exhaustivo de referencia compara el
renglón con cada producto del catálogo en lugar de cada concepto; con 1.030 productos con precio
(Sección~\ref{sec:recuperacion}), su costo es $O(N\,d)$ con $N=1.030$, unas 1,05 millones de
multiplicaciones, es decir, 6,9 veces más que la resolución por conceptos. Esa relación es el cociente
$N/|\mathcal{C}|$ entre productos y conceptos, y crece con el catálogo, mientras que el número de
conceptos crece más despacio porque agrupa productos equivalentes. Ambos métodos son lineales, de
modo que la ventaja no es asintótica sino de constante. La razón principal para agrupar por conceptos
es la abstención: el exhaustivo por producto no distingue entre un producto ausente y un
concepto que el catálogo no cubre.

% =====================================================================
\subsection{Optimización del kit}
\label{sec:optimizacion-dev}

Con los productos resueltos, el kit se arma sobre las líneas comprables con cuatro objetivos: costo,
fidelidad de marca, disponibilidad y personaje. Ninguna ponderación fija representa a todos los
acudientes, así que el sistema construye el frente de Pareto en vez de optimizar una función de
utilidad única \citep{zitzler_multiobjective_1999}.

La prueba de concepto usaba NSGA-II
\citep{deb_fast_2002} con una población de 60, 100 generaciones, un tope de 50 ítems y 30 segundos,
y fundía la calidad con el personaje. En el banco de alternativas, sobre nueve instancias
sintéticas, quedó a un 0,07\,\% del óptimo, frente al 11,92\,\% de un algoritmo voraz. El sistema
construido conserva el algoritmo, con pymoo \citep{blank_pymoo_2020}, y cambia tres cosas. Calcula
el frente exacto cuando hay 4.096 combinaciones o menos; separa el personaje como objetivo propio; y
hace crecer la población y las generaciones con la lista, con una semilla fija. El servicio no
arranca sin pymoo.

Del frente salen siempre el Ahorrador y el Premium. El Fan sale solo si el acudiente pidió un
personaje, y «A tu Medida», el de mejor marca que cabe en el tope o, sin tope, a mitad de camino
entre los otros dos, solo si eligió el punto medio de gasto o declaró un tope (ADR 0021). Ofrecer
pocas opciones bien diferenciadas facilita decidir
\citep{thaler_nudge_2008,iyengar_when_2000,schwartz_paradox_2003} (NR-02), y un kit que nadie
pidió es ruido para el acudiente y trabajo para el revisor. La regla de recomendación marca el Fan
cuando hay personaje y «A tu Medida» cuando se armó; si no, el Premium cuando el acudiente pone la
calidad primero y acepta un sobreprecio del 10\,\%, y en los demás casos el Ahorrador. La regla sale
de la encuesta del Dossier OE1, en la que el 70\,\% de los acudientes no pagaría más de un 5\,\% de
sobreprecio. El tope es el presupuesto por hijo, sin ampliarlo por el sobreprecio aceptado. Si el
kit marcado no cabe, se marca el de mejor marca que sí cabe y, solo si ninguno cabe, el más barato,
con la razón y su costo (ADR 0020). El presupuesto acota la recomendación, no el frente: los demás
kits siguen a la vista.

Los prototipos inferían las
preferencias con reglas y, de forma opcional, con un modelo de lenguaje local. El sistema construido
las pide como opciones cerradas por hijo: cuánto gastar, el presupuesto, el sobreprecio aceptado y
el personaje (PS-01, PS-02 y PS-04). El SDD lo justifica en tres razones: la preferencia queda
explícita y auditable, no depende de un modelo y el acudiente puede cambiarla. Los pesos por omisión
son 0,35, 0,30, 0,25 y 0,10, y salen de la pregunta P21 de la encuesta normalizada (42 de 60
encuestados pagan como máximo un 5\,\% más y el personaje quedó último en la prioridad). Los de OP-04 del
SRS v6.1, 0,285, 0,273, 0,255 y 0,187, se contrastan con los anteriores en la calibración de la
utilidad multiatributo (Sección~\ref{sec:optimizacion-cal})\pendienteEvidencia{pesos MAUT finales}. El
hipervolumen del frente, indicador de calidad de un conjunto de soluciones no dominadas
\citep{zitzler_multiobjective_1999}, se calcula sobre listas reales en la
Sección~\ref{sec:optimizacion}\pendienteEvidencia{hipervolumen sobre listas reales}.

% =====================================================================
\subsection{Constantes numéricas del sistema}
\label{sec:constantes}

En esta sección se registra cada constante numérica que condiciona el comportamiento del sistema,
para que ninguna quede sin justificación. Para cada una, la Tabla~\ref{tab:constantes} indica su
valor, el módulo donde vive, su origen y su estado. El origen es de uno de tres tipos: una
referencia de la literatura, un requisito del SRS con su razón, o un experimento de los
prototipos o de la calibración. El estado es «vigente» cuando el valor tiene una justificación
registrada y «en ciclo de calibración» cuando el valor actual es el de partida de un ciclo de
revisión humana (Sección~\ref{sec:ciclos-calibracion}) cuyo resultado final se reporta en la
Sección~\ref{sec:optimizacion-cal}. Los valores proceden del código del repositorio y de la tabla de
parámetros del documento de algoritmos del proyecto.

{\footnotesize
\setlength{\tabcolsep}{3pt}
\renewcommand{\arraystretch}{1.1}
\begin{longtable}{@{}p{3.1cm}p{2.3cm}p{2.5cm}p{4.7cm}p{1.6cm}@{}}
\caption{Constantes numéricas del sistema, con su origen y su estado.}\label{tab:constantes}\\
\toprule
Constante & Valor & Módulo & Origen & Estado \\
\midrule
\endfirsthead
\toprule
Constante & Valor & Módulo & Origen & Estado \\
\midrule
\endhead
\bottomrule
\endfoot
\multicolumn{5}{@{}l}{\emph{Percepción}}\\
Tope de páginas & 40 & percepción & Experimento: el PDF más largo del corpus tiene 39 páginas & Vigente \\
Temperatura del modelo & 0 & percepción & Decodificación voraz, para reducir la variabilidad entre lecturas & Vigente \\
Tiempo límite por llamada & 40 s & percepción & Una página típica tarda unos 5 s; 40 s indica una llamada colgada & Vigente \\
Reintentos y espera & 1 reintento, 3 s & percepción & Solo ante fallas de transporte, para no ocultar respuestas mal formadas & Vigente \\
\midrule
\multicolumn{5}{@{}l}{\emph{Normalización y resolución}}\\
Dimensión del vector & 1.024 & catálogo & Fija la columna de la base; se pide explícitamente al modelo de \emph{embeddings} & Vigente \\
Lote de \emph{embeddings} & 100 textos & catálogo & Una petición por lote, para evitar cuerpos de petición excesivos & Vigente \\
Umbral de reconocimiento & 0,5 & catálogo & Corte por debajo del cual el renglón no corresponde a ningún concepto; fijado al corregir el caso «chaquiras» & Vigente \\
Aporte del coseno $\gamma$ & 0,15 & catálogo & Ajuste sobre casos de error; pequeño para que el coseno desempate y no decida & En ciclo de calibración \\
Umbral de rescate $\theta$ & 0,86 & catálogo & Ajuste sobre casos de error; separa un sinónimo de un parecido de tema & En ciclo de calibración \\
Piso del rescate $\phi$ & 0,55 & catálogo & Apenas sobre el corte de 0,5, para que un rescate entre a revisión & En ciclo de calibración \\
Umbral de autoaceptación $\tau$ & 0,75 & inteligencia & Requisito PE-04 del SRS; valor inicial anterior a la calibración (0,85 y 0,630 en versiones previas); ver Tabla~\ref{tab:cal-optim} & En ciclo de calibración \\
Tolerancia de medidas & 0,25 y rango 0,85 a 1,35 & catálogo & Ajuste con las auditorías de los kits (Tabla~\ref{tab:kits}) & Vigente \\
Rango de precio en tiendas & \$1.000 a \$400.000 & catálogo & Descarta accesorios sueltos y lotes institucionales & Vigente \\
\midrule
\multicolumn{5}{@{}l}{\emph{Optimización}}\\
Límite de enumeración & 4.096 ($2^{12}$) & inteligencia & Por debajo de ese tamaño del espacio el frente exacto se calcula en milisegundos; por encima se usa NSGA-II & Vigente \\
Población & $\min(200,\max(40,12n))$, par & inteligencia & Escala con las $n$ líneas decidibles; el par evita desperdiciar un individuo & Vigente \\
Generaciones & $\min(120,\max(30,20\log_{10}|\mathcal{E}|))$ & inteligencia & Escala con el logaritmo del espacio $|\mathcal{E}|$ de búsqueda & Vigente \\
Cruce SBX & $p_c=0{,}9$, $\eta_c=15$ & inteligencia & Valores estándar de \citet{deb_simulated_1995} & Vigente \\
Mutación polinomial & $p=1$, $\eta_m=20$ & inteligencia & Valores estándar de \citet{deb_combined_1996} & Vigente \\
Semilla & 0 & inteligencia & Reproducibilidad del frente entre ejecuciones & Vigente \\
Pesos por omisión (costo, marca, calidad, personaje) & 0,35; 0,30; 0,25; 0,10 & inteligencia & Encuesta, pregunta P21 (n=60), normalizada; ver Tabla~\ref{tab:cal-optim} & En ciclo de calibración \\
Sobreprecio para Premium & 10\,\% & inteligencia & Encuesta: 42 de 60 pagan como máximo un 5\,\% más; el 10\,\% marca a quien acepta pagar más & Vigente \\
Opciones por renglón & 3 & catálogo & Pocas opciones diferenciadas facilitan la decisión \citep{iyengar_when_2000} & Vigente \\
\midrule
\multicolumn{5}{@{}l}{\emph{Requisitos de desempeño y operación}}\\
Tiempo de generación del kit & 8 s & SRS (DP-03) & Requisito de desempeño; se contrasta en la Sección~\ref{sec:optimizacion} & Vigente \\
Ventana de revisión & 4 h en temporada & SRS (ML-09) & Requisito: el administrador revisa cada lista nueva antes de publicarla & Vigente \\
Disyuntor & 5 fallas, 30 s & biblioteca común & Patrón de resiliencia; 30 s es el escenario ESC-R2 del ATAM & Vigente \\
Caché de lecturas & 30 días & inteligencia & Evita pagar de nuevo la percepción de un mismo documento & Vigente \\
Caché de vectores & 24 h & catálogo & Evita recalcular el vector de un renglón repetido & Vigente \\
\end{longtable}
}

% =====================================================================
\subsection{Aplicación web y pago}
\label{sec:web}

El acudiente usa una aplicación web progresiva. Siguiendo las pautas de interacción con sistemas de
IA \citep{amershi_guidelines_2019}, la interfaz muestra el estado de cada renglón en nueve clases y
dice «Lo está revisando una persona» cuando la lista espera la revisión. El acudiente elige uno de
sus kits y puede cambiar el producto de una línea por el sugerido o por una de sus alternativas:
el cambio se guarda en su solicitud, se aplica al leer el kit y se cobra (ADR 0019). La dirección
de entrega se completa con las sugerencias de Google Places y, sin ese servicio, se escribe a
mano (ADR 0035).

El pago se abre en el widget de Wompi, una ventana sobre la página del kit o de «Mis pedidos»
(ADR 0024). Con tarjeta o Nequi el acudiente no sale de Tukit; con PSE, el banco toma la pestaña y
Wompi lo devuelve a la página de resultado. Comercio firma los datos del widget y guarda solo la
referencia, el identificador, el estado y el medio de la transacción (ADR 0010). La confirmación
llega por el aviso de la pasarela y por la consulta de la transacción cuando el acudiente regresa,
y las dos pasan por el mismo procedimiento idempotente (ADR 0025). El acudiente elige el medio dentro del widget, de modo que Tukit no da a PSE la prioridad que
preveía el ADR 0010. El flujo de pago se verificó en el entorno de pruebas (\emph{sandbox}) de la
pasarela con tarjeta, Nequi y PSE\pendienteEvidencia{pago de extremo a extremo en el sandbox de Wompi}. El cobro con
dinero real depende del despliegue productivo y se delimitó como trabajo futuro
(Sección~\ref{sec:trabajo-futuro}).

% Contenido que llega de otras sesiones (calibracion y MAUT; NL2C).
\subsection{Calibración de la optimización multiobjetivo}
\label{sec:optimizacion-cal}

En esta sección se describen la calibración del optimizador y el método de utilidad multiatributo que ordena los kits del frente de Pareto.

\subsection{Traducción de preferencias en lenguaje natural a restricciones}
\label{sec:nl2c}

En esta sección se describe cómo una preferencia escrita en lenguaje natural se traduce a restricciones estructuradas para el optimizador.


    \setcounter{subsection}{0}
\section{RESULTADOS}

Este capítulo presenta los resultados de las Fases~3, 4 y~5 del proyecto, ligadas a los objetivos
tercero, cuarto y quinto. Los de los componentes de datos y de recuperación siguen la fase de
evaluación de CRISP-DM, y los del código, el plan del informe de pruebas
(Sección~\ref{sec:calidad}). Cada cifra va con su muestra, su condición y su grado de respaldo: recalculada desde una
salida conservada, respaldada por una salida existente o solo reportada. Cuando dos fuentes
discrepan o falta el programa que produjo una cifra, el texto lo dice en lugar de dar un valor.

% =====================================================================
\subsection{Percepción documental}
\label{sec:percepcion}

El primer criterio del árbol de utilidad exige extraer de una lista escolar, entregada como imagen
o PDF sin estructura, los renglones con su cantidad y su descripción. El 1 de septiembre de 2026,
los prototipos de tesis compararon cinco motores sobre CGS-10, un conjunto de oro de diez documentos
y 1.293 ítems. La métrica es F1 en sus variantes micro, agregada sobre todos los renglones, y macro,
promediada por documento (Tabla~\ref{tab:percepcion}).

\begin{table}[htbp]
\centering
\caption{Percepción documental en los prototipos de tesis (CGS-10, diez documentos)}
\label{tab:percepcion}
\small
\begin{tabular}{@{}lrr@{}}
\toprule
\textbf{Enfoque} & \textbf{F1 micro} & \textbf{F1 macro} \\
\midrule
Gemini (multimodal) & 0,948 & 0,944 \\
Google Vision & 0,673 & 0,672 \\
GOT-OCR2.0 & 0,597 & 0,569 \\
Docling & 0,364 & 0,424 \\
OCR con reglas heurísticas & 0,298 & 0,419 \\
\bottomrule
\end{tabular}
\end{table}

Los valores se recalcularon desde la salida conservada de la comparación, con su regla de
emparejamiento (similitud mínima de 0,75), sin volver a ejecutar los motores. El acta del OE3, que
está pendiente de firma y no forma parte de la entrega, consigna 0,925, un valor que no proviene de
esa salida. La Figura~\ref{fig:bench-F3-2} añade la latencia por documento que registra la misma
salida; puede incluir esperas por el límite de tasa, de modo que no sirve para verificar DP-04 del
SRS. La latencia de la percepción en el sistema construido sigue pendiente: la verificación usó una
percepción simulada, y la medición con el modelo real corresponde al segundo nivel de las pruebas
(ESC-D1 del ATAM).

La comparación tiene tres amenazas a su validez. La principal es el oro: CGS-10 se armó corrigiendo,
en cada documento, la extracción de uno de los motores evaluados, sin un segundo anotador. Coincide
con la salida de Gemini en el 86,7\,\% de los ítems literales y en el 93,1\,\% sin tildes ni
puntuación, así que el 0,948 no es independiente de esa salida. La segunda es la muestra: diez
documentos y una sola corrida por motor. La tercera es que el sistema construido usa otro modelo de
Gemini, con temperatura 0 y sin semilla. La ventaja del modelo multimodal frente a las reglas es
amplia, pero su magnitud está sesgada a favor del motor elegido.

\input{\graficasdir/F3-2.tex}

% =====================================================================
\subsection{Recuperación y normalización}
\label{sec:recuperacion}

El 2 de septiembre, los prototipos de tesis compararon cuatro estrategias sobre PostgreSQL con
pgvector: léxica pura, léxica expandida, vectorial densa e híbrida con fusión recíproca de rangos
(RRF, $k=60$). El sistema construido adoptó en su lugar la fusión aditiva de una sola fase
descrita en la Sección~\ref{sec:normalizacion}.

\subsubsection{Latencia}

La léxica pura tuvo un p50 de 6,66\,ms y una media de 40,60\,ms; la léxica expandida, de 6,59 y
52,76\,ms; la vectorial densa, de 46,51 y 106,61\,ms, y la híbrida con RRF, de 53,17 y 155,51\,ms.
La medición registra el tiempo y el número de resultados, no la relevancia, y no documenta un
calentamiento; la media está entre dos y ocho veces el p50 por las colas largas. En la misma fecha, con 50 consultas
en CPU, pgvector con índice HNSW tuvo una media de 5,06\,ms (p95 de 6,22\,ms); la búsqueda exacta en
pgvector, 363,68\,ms, y el índice HNSW de FAISS, 0,098\,ms. El QA acepta los 5,06\,ms de pgvector
por sus garantías transaccionales. La latencia de la recuperación del sistema construido sigue
pendiente; la prueba de carga de la verificación midió la consulta del catálogo, no la recuperación
de candidatos (Sección~\ref{sec:calidad}).

\subsubsection{Calidad de la recuperación}

Con la biblioteca de MTEB \citep{muennighoff_mteb_2023} se compararon dos modelos de embedding en
la tarea SpanishPassageRetrievalS2P, de pasajes de salud y de dominio distinto al escolar.
multilingual-e5-large obtuvo un MRR@10 de 0,6163, frente a 0,4735 de
paraphrase-multilingual-mpnet-base-v2; las salidas no registran la versión de la biblioteca. Los
550,14\,ms y 69,88\,ms que algunos documentos técnicos presentan como latencia por consulta resultan
de dividir el tiempo total de la tarea entre un número aproximado de consultas (1.088,20\,s entre
1.978,03 en el modelo grande), así que no lo son. El resultado tampoco describe el sistema
construido, que calcula los embeddings con gemini-embedding-001 desde el 23 de septiembre de 2026.

Sobre Recall@10 hay una discrepancia sin resolver. Un informe de los prototipos reporta 1,0000. La
única salida conservada, la de las 50 consultas, da 0,10 tanto para FAISS como para pgvector con
HNSW frente a la búsqueda exacta, y el programa que la produjo no se conserva. En el banco de
alternativas, la búsqueda exacta sobre 10.000 vectores sintéticos mantuvo un Recall@10 de 1,0. Este
capítulo no afirma un valor de Recall@10 para los índices aproximados.

\input{\graficasdir/F5-2.tex}

\subsubsection{Acuerdo en la clasificación del corpus}

El 7 de septiembre, los prototipos de tesis midieron un kappa de Cohen de 0,7331 sobre 200 pares de
12 documentos, que se reprodujo desde las entradas conservadas; su intervalo de confianza del
95\,\% por remuestreo es [0,670; 0,794]. Un segundo valor, 0,7049 con 25 documentos, no se puede
recalcular. El kappa no mide el acuerdo entre dos personas: compara la categoría que asigna una
heurística con la que asigna un modelo de lenguaje, en 45 de las 48 categorías del corpus. Por eso
no cumple la meta del SPMP, un kappa de 0,7 o más con un segundo anotador humano sobre al menos 30
documentos, y el límite inferior del intervalo queda además por debajo de 0,7.

\subsubsection{Umbral de resolución automática}

Los tres valores del umbral (Sección~\ref{sec:normalizacion}) operan sobre escalas distintas y no
son intercambiables. El de 0,630 se calibró con 542 ítems y un AUC de 0,880, pero 186 de sus
etiquetas (34,3\,\%) se propagaron a descripciones idénticas. La muestra efectiva es de 356
descripciones, así que los intervalos que tratan los 542 ítems como independientes son más estrechos
de lo justificado, y el mismo conjunto sirvió para elegir el umbral y para medirlo. La
Figura~\ref{fig:bench-F5-2} muestra cómo cambian la cobertura y la precisión del primer candidato
al desplazarlo; en 0,630 se resuelven 354 de los 542 ítems (65,3\,\%), y 307 de ellos (86,7\,\%)
correctamente.

La única medición del umbral de 0,75 pertenece a una línea de desarrollo del 4 de octubre, anterior
a la versión evaluada y no integrada en el sistema construido. Un conjunto de oro de 2.588
renglones, tomados de 4.719 ítems y 98 conceptos, midió la regla de una versión anterior. Las
versiones evaluada y construida conservan el umbral, pero cambiaron el análisis del renglón y el
catálogo, de modo que la medición no las describe. Aquella regla resolvió de forma directa 2.252 renglones, 2.234 de ellos
correctos (precisión directa de 0,992), dejó 298 en la zona gris (11,51\,\%) y 38 fuera de la
resolución. Sobre todas las líneas, el acierto fue de 0,9784; si en la zona gris decide el
acudiente, de 0,9737, y si decide el sistema, de 0,8632. La medición solo ejercita el canal léxico,
porque los 609 productos tenían el embedding vacío, y 1.138 de sus renglones son coincidencias
exactas triviales. El umbral de 0,75 sigue sin calibrar.

\subsubsection{Extremo a extremo}

La cifra de 83\,\% sobre 60 sublistas y 2.176 ítems, del servicio de búsqueda anterior, mide la
fracción de ítems con al menos un candidato recuperado (83,04\,\%), no la resolución. En esa
corrida, 769 ítems (35,3\,\%) se resolvieron de forma automática, 795 (36,5\,\%) pasaron a revisión
y 612 (28,1\,\%) quedaron sin correspondencia. Se ejecutó sin conexión, en una base de prueba, sobre
un espejo de 838 productos con los 48 nombres más cortos de cada categoría.

De las primeras versiones del sistema construido hay una medición de distribución, no de precisión. Entre el 21 y el 23 de
septiembre, un ensayo de huecos aplicó el pipeline a 14 PDF de listas, con 199 páginas y 4.414
renglones. Se reconocieron 2.202; 926 terminaron en abstención (20,98\,\%), 626 sin coincidencia,
66 como hueco del catálogo y 594 en una clase «fuera» que los documentos no definen. El ensayo dejó
1.397 faltantes únicos para la curaduría del catálogo.

\subsubsection{Auditoría de los kits}

Las mediciones anteriores evalúan la resolución de cada renglón. Una línea de afinamiento de los
kits evaluó además el resultado que recibe la familia. Una auditoría previa, del 7 de octubre,
revisó los kits de un grado, Séptimo: al incorporar los atributos del catálogo, las líneas del
Ahorrador que cumplen pasaron de 31 a 49 y las que no cumplen, de 30 a 3, y de ahí salieron los
atributos y el filtro por renglón. El 11 y el 12 de octubre de 2026 se auditaron cuatro veces, renglón por renglón, los kits
Ahorrador y Premium de las trece listas de 2026 del Colegio Agustiniano Ciudad Salitre. Un renglón
cumple si los dos perfiles cobran lo que pide la lista, y son comprables todos los renglones salvo
los que Tukit no vende, como los uniformes o los textos con título. Entre una auditoría y la
siguiente se ajustaron la comparación de atributos y la curaduría de las listas y del catálogo, y
antes de la cuarta el catálogo pasó de 679 a 1.030 productos con precio (Tabla~\ref{tab:kits}).

\begin{table}[htbp]
\centering
\caption[Auditorías de los kits del Colegio Agustiniano Ciudad Salitre]{Auditorías de los kits del Colegio Agustiniano Ciudad Salitre, 11 y 12 de octubre de
2026; las dos últimas, posteriores a la versión evaluada (cifras reportadas; costo de los trece kits de cada perfil)}
\label{tab:kits}
\small
\begin{tabular}{@{}lrrrrr@{}}
\toprule
\textbf{Auditoría} & \textbf{Comprables} & \textbf{Huecos} & \textbf{Cobros} &
\multicolumn{2}{c}{\textbf{Costo (millones de pesos)}} \\
\cmidrule(l){5-6}
 & \textbf{que cumplen} & \textbf{declarados} & \textbf{sin cumplir} & \textbf{Ahorrador} &
\textbf{Premium} \\
\midrule
1 (11 oct.) & 263 de 888 (29,6\,\%) & 240 & 589 & 13,9 & 32,0 \\
2 (11 oct.) & 465 de 888 (52,4\,\%) & 385 & 242 & 9,9 & 17,6 \\
3 (11 oct.) & 617 de 886 (69,6\,\%) & 220 & 49 & 9,0 & 14,2 \\
4 (12 oct.) & 758 de 885 (85,6\,\%) & 69 & 58 & 8,1 & 17,3 \\
\bottomrule
\end{tabular}
\end{table}

El cumplimiento subió en cada auditoría. Los cobros de un producto que no corresponde a lo pedido
bajaron de 589 a 49 en las tres primeras y subieron a 58 en la cuarta, sobre todo por 23 renglones
en los que el Ahorrador cobra un cuaderno que no dice ser grande. Los huecos no bajaron de forma
continua: los 240 iniciales incluían 180 renglones que Tukit no vende, y de los 69 finales, 45 son
de electrónica, que las tiendas registradas no ofrecen. El Premium se encareció en la cuarta
auditoría porque cobra 151 renglones que antes eran huecos. Las cifras provienen de los informes de
auditoría y cubren un solo colegio, cuyas listas se cargaron desde datos curados. Dependen de una
clasificación humana de cada renglón. Las auditorías 1 y 2 midieron versiones del 5 y el 8 de
octubre, que forman parte de la historia de la versión evaluada, y la segunda tiene su mismo código
de producto. Las auditorías 3 y 4 midieron versiones intermedias del afinamiento, que el sistema
construido integra con cambios posteriores en la tolerancia de las medidas, el tamaño del cuaderno
y el catálogo, de modo que ninguna auditoría mide el sistema construido. El anexo
CRISP-DM detalla cada iteración y sus límites.

% =====================================================================
\subsection{Optimización}
\label{sec:optimizacion}

Sobre la calidad del optimizador (Sección~\ref{sec:optimizacion-dev}) hay tres evidencias de
alcance distinto. El 31 de julio, el banco de alternativas comparó NSGA-II con un algoritmo voraz y
uno aleatorio en nueve instancias sintéticas: halló el óptimo en ocho y quedó a un 0,07\,\% de él,
frente al 11,92\,\% del voraz y el 159,34\,\% del aleatorio. Solo se conserva el informe, y el
problema tenía tres objetivos, de modo que no valida el optimizador del sistema construido. En este,
las pruebas calculan un hipervolumen \citep{zitzler_multiobjective_1999} en dos dimensiones sobre un
ejemplo de enumeración y comprueban, con 65.536 combinaciones, que el frente de NSGA-II no tenga
soluciones dominadas y que la semilla lo haga reproducible; no hay registro de su ejecución. Un
banco del 4 de octubre, que forma parte de la versión evaluada, compara NSGA-II con un óptimo exacto
por programación dinámica sobre 174 listas del corpus. En el Ahorrador lo iguala en el 100\,\% de
las listas. Con presupuesto, el kit recomendado daba el 40\,\% de la calidad posible; después del
ADR 0020 cabe siempre en el tope e iguala el óptimo. Son cifras reportadas con una sola semilla y no
reproducidas, y la resolución usó solo el canal léxico y una foto fija del catálogo. El hipervolumen
sigue sin calcularse fuera de la prueba automatizada.

Sobre 29 listas reales, otra medición reportada dio 5,3~s de media, 6,3~s de mediana y 9,0~s de
máximo para producir todos los kits de una lista, en una CPU de 16 núcleos sin otras cargas. El
máximo supera los 8~s de DP-03. No es la prueba del escenario ESC-D4 del ATAM, que pide al menos 30
ejecuciones y el percentil 95, y midió el optimizador anterior a las correcciones del 4 y el 7 de
octubre. El tiempo en el peor caso de la versión evaluada sigue pendiente.

% =====================================================================
\subsection{Verificación de la versión evaluada}
\label{sec:calidad}

Los porcentajes de cobertura que citan algunos documentos sobre los microservicios anteriores no
son mediciones. Son los umbrales mínimos que su integración continua exigía para aprobar: 90\,\% en
cuentas y en comercio, 84\,\% en instituciones y 80\,\% en búsqueda. Esas ejecuciones terminaron con
éxito el 18 de septiembre de 2026, de modo que la cobertura alcanzó al menos el umbral, pero no se
conserva el valor medido.

La versión evaluada se verificó con un plan de diez fases, distintas de las cinco del proyecto:
la línea base (0), la matriz de requisitos no funcionales (1), las pruebas unitarias, de
integración, funcionales con conjuntos de oro, adversariales de la IA y de los patrones de
resiliencia (2 a 6), la degradación de componentes (7), la carga (8), la telemetría (9) y el
informe (10). Es un resultado de la Fase~4 de la VFP, que espera pruebas automatizadas del motor de
optimización y de las interfaces, y cubre en parte las pruebas funcionales, no funcionales y de
carga de la Fase~5. El informe de pruebas de la entrega reúne el método, las evidencias y los
resultados. Las fases 0 a 6 corrieron sobre la versión de integración del 8 de octubre de 2026,
cuyo código de producto es el mismo de la versión evaluada salvo la integración continua y dos
valores por omisión de la base de datos. El primer nivel de las fases 7 a 9 corrió sobre la versión
evaluada, en una sola corrida, en una máquina virtual de 4 procesadores y 8~GiB y con la percepción
simulada. La verificación no hizo ninguna llamada real al modelo: el simulador atendió 2.342
percepciones.

Con las pruebas propias del sistema pasan 1.647 pruebas, y la cobertura es de 72,43\,\% de 12.039
sentencias. Con las que añadió la verificación, pendientes de integrar, pasan 3.911, otras 129 son
fallas esperadas que documentan defectos y la cobertura sube a 94,32\,\%. La integración continua
ejecuta las pruebas, pero no mide la cobertura, y su estado en la versión evaluada no se comprobó.
El modelo del catálogo difiere de sus migraciones en dos valores por omisión. De las 60 filas de la
matriz que liga cada meta con su ensayo, se ejecutaron 36: 7 cumplen, 7 cumplen en lo medido a la
espera de los niveles 2 y 3, y 22 no cumplen. Otras 13 están pendientes, 2 esperan un hallazgo ya
anticipado, 8 no se pueden medir con los componentes actuales y 1 es el protocolo con usuarios.
Quedan 78 hallazgos abiertos: 14 de severidad alta, 56 de severidad media y 8 de severidad baja.

El primer nivel de las fases 7 a 9 dejó 18 ensayos, de los que 14 cumplen y 4 no. En la carga, una
rampa hasta 100 usuarios generó 27.894 peticiones sin errores, y el percentil 95 de la consulta del
catálogo fue de 22~ms con un usuario y de 160~ms con cien, por debajo de los 300~ms de la meta. No
hubo ningún rechazo por límite de tasa porque la versión evaluada no tenía limitador, de modo que en
ella esa meta no se cumple; el sistema construido añade uno con Redis, pendiente de medición. En la degradación, siete de diez caídas de un componente cumplieron: los demás servicios
siguieron respondiendo sin errores y se recuperaron en 2,5~s o menos. Con el catálogo colgado, el
disyuntor abrió a los 26,3~s, dentro de los 30~s del escenario ESC-R2, porque cada intento falló
en unos 5~s. Si cada intento agotara su espera, el modelo da 300,8~s desde inteligencia y 150,8~s
desde comercio. Ese peor caso no se midió, y el hallazgo queda acotado, no corregido. Pasados los
30~s, con diez llamadas simultáneas, el disyuntor deja pasar las diez en lugar de una sola de
prueba.

Las tres caídas que no cumplieron dejaron ver tres limitaciones de la versión evaluada. Con el
catálogo detenido, subir una lista respondía con un error 500 y la lista no se guardaba. Después de
caer la base de datos o el intermediario de mensajes, los procesos de fondo morían y no se
reiniciaban. Y un mismo pago producía dos avisos al acudiente, sin cobro doble. El sistema
construido corrige las tres. Con el catálogo caído, la subida responde 202, guarda lo percibido y
la lista se importa cuando el catálogo vuelve (ADR 0031). Los procesos de fondo sobreviven a las
caídas pasajeras de la base o del intermediario, aunque un error de programación todavía los
detiene (ADR 0032), y un evento produce un solo aviso (ADR 0033). La corrección se verificó el 20
de octubre de 2026 en una versión de trabajo, antes de integrarla: la subida responde 202 en
33~ms y se recupera en 2,7~s, los procesos de fondo se recuperan en 7,1~s y 6,9~s, y el pago
produce un solo aviso. Esa verificación no volvió a correr los demás ensayos, y la corrida
completa sobre el sistema construido está pendiente.

Quedan pendientes la corrida completa sobre el sistema construido y la medición de lo que añadió
después de la versión evaluada: las colas por RabbitMQ, la caché y el límite de tasa con Redis, las sugerencias de dirección, el alta de
colegios no aliados y la comparación por atributo. También lo están los niveles 2 y 3 de la
degradación, la carga y la telemetría, con la percepción real, el pico estacional y las caídas
combinadas, y el estudio con usuarios de la Fase~5
(Sección~\ref{sec:oe5}). En la telemetría, las trazas no llevan el identificador de correlación, de
modo que una carga no se puede seguir de extremo a extremo.

% =====================================================================
\subsection{Auditoría de seguridad y privacidad}
\label{sec:defectos}

Dentro del OE2 se hicieron dos auditorías, listadas en la Tabla~\ref{tab:anexos}: la de
autorización y la de tratamiento de datos personales (Anexo~M), esta conforme a la Ley 1581 de 2012
\citep{congreso_ley1581_2012} y al Decreto 1377 de 2013 \citep{presidencia_decreto1377_2013}. Las
dos revisaron el código de los microservicios anteriores. La de autorización revisó 152 puntos de
entrada y encontró 22 sin una comparación completa de ámbito: 8 sin ninguna y 14 con una parcial.
El Anexo~M registra catorce defectos, dos críticos y cinco de severidad alta.

\subsubsection{Defectos de los microservicios anteriores}

El de mayor alcance estaba en las políticas de seguridad a nivel de fila de las tablas con datos de
estudiantes. La condición base concedía visibilidad de la tabla completa cuando el contexto de
usuario no estaba establecido, de modo que el error más probable en un sistema distribuido, que un
servicio olvide fijar ese contexto, otorgaba acceso irrestricto. Una segunda cláusula daba acceso
completo al rol de administrador comparando solo el identificador del usuario, sin ancla de ámbito,
y ese defecto motivó el modelo de autorización con ámbito del capítulo de diseño. El ADR 0015 lo
registra en 16 tablas de cinco servicios y advierte que corregirlo va antes que su propia decisión;
las migraciones posteriores conservaban la apertura, sin pruebas de regresión. Entre los demás, el
manejador genérico de errores de validación devolvía al cliente el valor rechazado, contraseña
incluida, y la supresión de datos se marcaba como completada sin confirmación de los servicios.

\subsubsection{Estado en el sistema construido}

El sistema construido no tiene políticas de seguridad a nivel de fila. El ámbito se controla en la
aplicación, con la cadena de contención de identidad, que falla cerrada
(Sección~\ref{sec:autorizacion}); las cifras de 152 y 22 son las de la auditoría, no una medición
propia. De los catorce defectos del Anexo~M, dos están corregidos: el valor rechazado ya no vuelve
al cliente, y el correo del acudiente ya no queda en un registro que la supresión no alcanza. El
texto de la lista sin depurar se corrigió para la lista oficial, pero sigue en el documento que
carga el acudiente; la depuración del nombre se sustituyó por la minimización
(Sección~\ref{sec:percepcion-dev}), salvo en ese mismo documento. Los tres defectos de la supresión
cambiaron de forma, pero la supresión sigue sin mecanismo: el estado de propagación existe en el
modelo, sin publicador ni consumidor del evento.

Cuatro siguen vigentes: la gestión de cuentas no deja registro de auditoría, el consentimiento por
finalidad no se escribe ni se consulta, la respuesta de la importación de listas guarda el extracto
completo del documento y ninguna tabla de auditoría tiene una política de purga. Otro no tiene
efecto mientras no se escriba el registro de auditoría. Uno es parcial: las preferencias y los pesos
del acudiente están cifrados, pero la razón de la recomendación queda en claro y no se borra. El de
la anonimización no es evaluable, porque no hay ese proceso.

Un defecto de comercio quedó corregido en la versión evaluada. La consulta de los pedidos del
acudiente toma el usuario de la sesión, validada con identidad, e ignora el que llegue en la
consulta; el detalle, el pago y la cancelación de un pedido propio también exigen sesión, y un
pedido ajeno responde 404. La puerta de enlace sigue sin exigir sesión.

En el sistema construido siguen abiertos otros defectos. La consulta de contactos del servicio de identidad devuelve el correo de
cualquier usuario activo sin autenticación, y las rutas de los paneles de catálogo, inteligencia y
operación no exigen sesión ni rol. Cuatro caminos crean matrículas, y solo uno, el registro de un
hijo desde la cuenta del acudiente, escribe el consentimiento del menor. El celular del acudiente se
guarda sin cifrar en la base de identidad.

Las pruebas propias del sistema no detectaron ninguno de estos defectos. Medir cuánto código
ejecutan no dice si buscan las condiciones adversas correctas, y varias comprueban garantías leyendo
el código fuente en lugar de ejecutar el comportamiento. La verificación de la versión evaluada
(Sección~\ref{sec:calidad}) los confirmó y encontró otros. Un kit se puede comprar sin sesión, aun
mientras espera revisión, y varias rutas internas aceptan llamadas sin verificar quién llama. Las
cinco bases de datos comparten un mismo rol de acceso, de modo que las credenciales de cualquier
servicio leen los usuarios de identidad. La página del documento cargado sale sin depurar hacia el
proveedor del modelo (\mbox{PE-06}), y el nombre ficticio de las pruebas quedó guardado. El retracto existe
en comercio, pero ninguna aplicación lo ofrece y la puerta de enlace no lo expone; además, se acepta
pasado el plazo legal y no libera la reserva. Un pedido de 100.000 unidades fue aceptado para pago.
La respuesta de la puerta de enlace expone direcciones internas, el mensaje de un trabajo fallido
contiene la consulta SQL con sus parámetros y, sin la llave de cifrado, el cifrador guarda en claro
sin avisar.

% =====================================================================
\subsection{Validación con usuarios (OE5): planificada y pendiente}
\label{sec:oe5}

El quinto objetivo específico de la VFP pide evaluar el sistema con pruebas funcionales, no
funcionales, de carga y de usabilidad con usuarios reales, y medir la satisfacción y la reducción de
la carga cognitiva frente al proceso tradicional de compra. La verificación de la
Sección~\ref{sec:calidad} cubre en parte las pruebas funcionales, no funcionales y de carga. Las
sesiones con usuarios, con SUS y NPS, son una actividad planificada de la Fase~5 que todavía no se
ha hecho, y sus resultados no se estiman aquí. Los
escenarios de prueba del ATAM fijan los criterios de aceptación, entre ellos un SUS de 68 o más y 50
usuarios concurrentes con una tasa de error inferior al 1\,\%.

El método es el de la VFP. Antes de involucrar participantes exige un protocolo ético aprobado
\citep{lazar_research_2017}, conforme a la Ley 1581 de 2012 y al Decreto 1377 de 2013, porque la
interacción toca datos de menores. Prevé sesiones observadas con tareas representativas, el
cuestionario SUS \citep{brooke_sus_1996,bangor_empirical_2008,lewis_factor_2009} con al menos cinco
participantes e idealmente veinte o más, verbalización concurrente, el índice NPS
\citep{reichheld_one_2003} y, para la carga cognitiva, el conteo de interacciones discretas (clics,
búsquedas y validaciones manuales) frente al flujo tradicional. UI-07 del SRS v6.1, el QA, la WBS y
el ATAM usan NASA-TLX; como no hay una decisión registrada de ese cambio, el OE5 se juzga con el
método de la VFP.

El informe de pruebas propone un protocolo para esas sesiones. Recluta 30 acudientes que hayan
comprado útiles escolares en los dos últimos años y no hayan usado Tukit, para llegar a 26 sesiones
válidas; si no se alcanzan, el informe dará la muestra lograda y la potencia real, sin declarar
cumplimiento. Cada sesión, de unos 45 minutos, compara la compra en dos o tres tiendas en línea con
la obtención del kit en Tukit hasta un pago simulado, en orden contrabalanceado, con un estudiante
ficticio y sin datos de menores. Quedan cinco decisiones abiertas, entre ellas quién facilita y
observa las sesiones y quién tramita el aval de ética, si la universidad lo exige. El protocolo
ético aprobado, los consentimientos y el análisis no forman parte de esta versión.

     \setcounter{subsection}{0}

\section{CONCLUSIONES}

\subsection{Análisis de Impacto del Proyecto}

En esta sección se analiza el impacto del trabajo en dos planos. Primero se describe lo que aporta a la ingeniería de sistemas y luego su efecto en los contextos global, económico, ambiental y social.

\subsubsection{Impacto en la ingeniería de sistemas}

Además del sistema construido, el trabajo deja tres prácticas de ingeniería que se pueden
trasladar a otros proyectos, y cada una tiene un límite. Primero, la arquitectura se eligió
comparando alternativas con el método ATAM \citep{bass_software_2021}: un monolito modular,
microservicios de grano fino y la variante de grano grueso adoptada. La matriz ponderada del
documento ATAM da 3,91 sobre 5 a la variante elegida, frente a 3,44 y 2,82; el acta del
segundo objetivo había registrado 4,27 en una evaluación anterior, con otro conjunto de
alternativas. La decisión se apoyó en impulsores comprobables (recuperabilidad asimétrica,
frontera de cumplimiento y perfil de recurso heterogéneo) y no en la escalabilidad invocada
al principio, que el cálculo de carga descartó por órdenes de magnitud.

Segundo, el sistema se auditó más allá de sus requisitos funcionales. Las dos auditorías del
segundo objetivo revisaron el código de los servicios anteriores. La de autorización encontró
22 de 152 puntos de entrada sin una comparación completa de ámbito, y la de datos personales,
catorce hallazgos, entre ellos políticas de seguridad a nivel de fila que fallaban abiertas
cuando faltaba el contexto de usuario. El sistema construido corrigió una parte: ya no devuelve
los valores rechazados en los errores, no copia el correo del acudiente en los avisos y no
guarda el texto de la lista oficial. El ADR 0015 había decidido copiar las claves de ámbito
en cada tabla protegida y escribir contra ellas las políticas de fila; el sistema construido no
tiene esas políticas y controla el ámbito en la aplicación. Los
anexos de auditoría de autorización y de datos personales registran la situación de cada
hallazgo, y la Sección~\ref{sec:defectos} reúne los defectos que siguen abiertos. Tercero, las
fronteras entre servicios pasaron de convención documental a funciones de aptitud ejecutables
\citep{ford_building_2017}: doce pruebas contrastan el código con el manifiesto y otras
veintiocho comprueban que esas pruebas detectan infracciones plantadas. La integración continua
las ejecuta, aunque 11 de sus 12 ejecuciones registradas hasta el final de septiembre
terminaron en fallo y su estado en la versión evaluada no se comprobó. Las tres prácticas sirven a un proyecto que deba sostener ante terceros que su
diseño hace lo que declara, siempre que el seguimiento de los hallazgos forme parte del
método.

\subsubsection{Impacto en contextos globales, económicos, ambientales y sociales}

El problema que dio origen a Tukit es social antes que tecnológico. La encuesta del primer
objetivo específico (N=65 respuestas recibidas, n=60 analizables) describe una compra
costosa para los acudientes: el 65\,\% le dedica cuatro horas o más y el 68,3\,\% compra en
tres o más lugares. Las listas son, además, heterogéneas (1.451 formas distintas de nombrar
los mismos artículos sobre 3.608 renglones reunidos de 17 colegios), y cada hogar las
interpreta por su cuenta. Si el sistema reduce ese esfuerzo, el efecto sería económico, al
facilitar la comparación de precios y sustitutos que el proceso manual no ofrece, y de tiempo
para las familias en las semanas de mayor presión del calendario escolar. Ese efecto y la carga
cognitiva son lo que mide el OE5 (Sección~\ref{sec:oe5})\pendienteEvidencia{resultados SUS, NPS y carga cognitiva del OE5 para el impacto social}. El modelo de negocio adoptado
\citep{osterwalder_business_2010}, B2B2C, distribuye ese valor sin imponer un costo al colegio:
la institución no paga por participar, y el ingreso proviene de la compra del acudiente. En el
contexto colombiano, donde la capacidad de pago de las instituciones y de sus hogares es
heterogénea, una solución que dependiera del presupuesto del colegio limitaría su adopción a
las instituciones con mayor capacidad financiera. En el plano ambiental, el trabajo no midió un efecto directo. La disponibilidad que optimiza el
cuarto objetivo es la fracción de líneas del kit que Tukit puede surtir, un criterio comercial. Que
reduzca compras redundantes es una hipótesis cuya medición exige operar con familias, y la economía
circular de libros y útiles se plantea como trabajo futuro (Sección~\ref{sec:trabajo-futuro}). El alcance social del hallazgo tiene un límite: la encuesta se
concentró en un solo colegio del sector Salitre, con sobrerrepresentación del estrato 4
(68,3\,\%) y ausencia de los estratos 1 y 6. Por eso, generalizar la insatisfacción declarada a
la población escolar colombiana excede lo que esta muestra permite afirmar.

\subsection{Conclusiones y Trabajo Futuro}

En esta sección se evalúa el cumplimiento de cada objetivo específico y su razón, y luego se
presenta el trabajo futuro. La Tabla~\ref{tab:cumplimiento} resume, por objetivo, la evidencia, la
sección donde se presenta y el estado, y los párrafos siguientes explican por qué se asignó ese
estado.

\begin{table}[H]
\centering\footnotesize
\setlength{\tabcolsep}{3pt}
\caption{Cumplimiento por objetivo específico.}
\label{tab:cumplimiento}
\begin{tabularx}{\textwidth}{@{}p{0.9cm}Xp{2.9cm}p{2.3cm}@{}}
\toprule
OE & Evidencia & Sección & Estado \\
\midrule
OE1 & Corpus de 17 colegios, 3.608 renglones, 1.451 formas de nombrar y 99 conceptos canónicos; encuesta (n=60) & Cap.~III y IV & Cumplido \\
OE2 & Comparación ATAM de tres alternativas, árbol de utilidad de 17 escenarios, vistas 4+1 y 38 ADR & Cap.~V & Cumplido \\
OE3 & F1 micro de 0,948 en lectura (CGS-10); resolución con $\tau$ calibrado; NL2C & Sec.~\ref{sec:percepcion}, \ref{sec:optimizacion-cal} y~\ref{sec:nl2c} & Cumplido, con las amenazas de la Sec.~\ref{sec:amenazas} \\
OE4 & Optimizador de cuatro objetivos, igualdad con el óptimo exacto en 174 listas, MAUT, hipervolumen y aplicación web con pago de pruebas & Sec.~\ref{sec:optimizacion} y~\ref{sec:optimizacion-cal} & Cumplido\pendienteEvidencia{estado final de OE4 tras MAUT, hipervolumen y NL2C} \\
OE5 & Pruebas de unidad, integración, sistema y aceptación con usuarios (SUS, NPS, RTLX) & Sec.~\ref{sec:control-calidad} y~\ref{sec:oe5} & Cumplido\pendienteEvidencia{estado final de OE5 segun C1, C2 y C3 del plan de analisis} \\
\bottomrule
\end{tabularx}
\end{table}

El OE1 se cumplió porque el corpus permitió definir 99 conceptos canónicos, de los cuales 25
cubren cerca del 80\,\% de los renglones observados: el dominio es heterogéneo en la superficie,
pero tiene una estructura tratable. El OE2 se cumplió porque las decisiones de arquitectura se
justificaron con ATAM \citep{bass_software_2021}, vistas 4+1 \citep{kruchten_41_1995} y la norma
ISO/IEC/IEEE 42010 \citep{iso_42010_2022}; el diseño previó siete desplegables y el sistema
construido tiene seis, agrupados por razón de cambio, sobre cinco bases de datos (ADR 0013).

El OE3 se cumplió con un cambio de diseño entre los prototipos y el sistema construido. En
percepción, el modelo multimodal superó al reconocimiento con reglas (F1 micro de 0,948 frente a
0,298), con la salvedad de que el oro se construyó desde su salida. En resolución, el sistema
construido sustituyó la fusión recíproca de rangos y el reordenador por una sola fase que suma el
parecido léxico y el coseno, y compensa con la revisión de cada lista nueva antes de publicarla. El
OE4 se cumplió porque el optimizador enumera el frente exacto hasta 4.096 combinaciones y recurre a
NSGA-II \citep{deb_fast_2002,blank_pymoo_2020} por encima, con cuatro objetivos (costo, fidelidad
de marca, disponibilidad y personaje), y presenta los kits Ahorrador y Premium siempre, Fan si el
acudiente pidió un personaje y A tu Medida si eligió el punto medio de gasto o un tope. El personaje
es un hallazgo del desarrollo: la encuesta midió esa preferencia y los tres criterios de la VFP no la
recogían. El OE5 se evaluó con el método de la VFP y el plan de análisis fijado antes de la primera
sesión, y su estado depende de los criterios C1, C2 y C3 de ese plan.

\subsubsection{Trabajo futuro}
\label{sec:trabajo-futuro}

El trabajo futuro se organiza en seis líneas, alineadas con la sección de trabajo futuro de la
propuesta (VFP v4). Cada una quedó fuera de este trabajo por una razón explícita, que se indica a
continuación.
\begin{enumerate}
\item \textbf{Alianza con tiendas de uniformes e integración de libros.} Alianza con las tres
tiendas de uniformes y flujo con Buscalibre y las editoriales de textos escolares. Quedó fuera
porque este trabajo se centró en un solo negocio, el de útiles y papelería.
\item \textbf{Aplicaciones nativas y modo sin conexión.} El acudiente usa una aplicación web
progresiva; las aplicaciones nativas y el modo sin conexión amplían su alcance, pero no cambian el
algoritmo que se evaluó.
\item \textbf{Endurecimiento productivo.} Autorización uniforme en las rutas, rol de acceso por base
de datos, depuración del nombre del menor en el documento cargado, supresión de datos, medición de
la cobertura en la integración continua y cobro real con la pasarela. Es la condición para operar
con familias reales y se plantea antes del despliegue.
\item \textbf{Optimización de rutas.} Agrupar la entrega entre tiendas, que no afecta la
construcción del kit.
\item \textbf{Economía circular de libros.} Reventa e intercambio de textos entre hogares, que
requiere una base de usuarios activa.
\item \textbf{Calibración con más anotadores y más colegios.} Repetir la calibración del umbral y
de los pesos con un conjunto anotado por dos personas y con listas de más colegios, y comparar
NSGA-II con otros optimizadores de referencia sobre el mismo banco.
\end{enumerate}

    \section{REFERENCIAS}

% Estandar IEEE, coherente en todo el documento. Las entradas viven en
% bib/tukit.bib y se citan con \citep{clave} o \citet{clave}.
% Compilacion: pdflatex -> bibtex -> pdflatex -> pdflatex

\renewcommand{\refname}{}
\vspace{-2em}
\bibliographystyle{IEEEtran}
\bibliography{bib/tukit}

    \section{ANEXOS}

Los anexos acompañan a esta memoria en la carpeta Anexos de la entrega, cada uno con su fuente
editable y su PDF. La Tabla~\ref{tab:anexos} los lista con su contenido y su estado. Los
documentos técnicos de soporte (VFP, SRS, SPMP, WBS, SAD, SDD, documento maestro de diseño,
informes ATAM y QA, informe de pruebas, que está en la carpeta QA, Dossier de OE1 y diagramas
de OE2) están en sus carpetas de la entrega y se citan en el texto cuando se usan.

\begin{table}[htbp]
\centering
\caption{Índice de anexos de la entrega}
\label{tab:anexos}
\small
\renewcommand{\arraystretch}{1.2}
\begin{tabular}{@{}p{4.2cm}p{7.4cm}p{2.4cm}@{}}
\toprule
\textbf{Anexo} & \textbf{Contenido} & \textbf{Estado} \\
\midrule
Registros de decisión (ADR) & Treinta y ocho registros de decisión de arquitectura (0001 a 0038),
  en formato Nygard, reunidos en un solo documento & Entregado \\
\addlinespace
Actas de aprobación & Actas de los hitos OE1 y OE2 & Entregadas; sin firmas \\
\addlinespace
Auditoría de autorización & Autorización de los servicios anteriores y del sistema construido,
  con la situación de cada hallazgo & Entregado \\
\addlinespace
Anexo M, datos personales & Registro de tratamiento de datos personales del sistema
  construido, con la clasificación de los datos de menores y su base legal (Ley 1581 de 2012) &
  Entregado \\
\addlinespace
Diccionario de datos & Diccionario de datos y esquema objetivo & Entregado \\
\addlinespace
Empatía y entrevistas & Guiones de entrevista y mapas de empatía por actor & Entregado \\
\addlinespace
Matriz de evidencia & Matriz de objetivos específicos y evidencia & Entregado \\
\addlinespace
Priorización MoSCoW & Priorización de los requisitos del SRS v6.1, con su estado en el sistema
  construido & Entregado \\
\addlinespace
CRISP-DM de recuperación & Recuperación ordenada por las seis fases de CRISP-DM, con la
  evolución desde los prototipos y el estado, las pruebas y la justificación de cada
  estrategia & Entregado \\
\addlinespace
Diagramas de clases & Clases de los servicios y de la biblioteca común & Entregado \\
\midrule
Previstos en la WBS & Protocolo ético y consentimientos, instrumentos de validación con
  usuarios, kits e hipervolumen del frente de Pareto y demostración del incremento con su
  retrospectiva. La cobertura se reporta en el informe de pruebas & Se incorporan con la entrega final \\
\bottomrule
\end{tabular}
\end{table}

El protocolo ético, los consentimientos y los instrumentos de validación corresponden al OE5
(Sección~\ref{sec:oe5})\pendienteEvidencia{anexo de protocolo etico, consentimiento e instrumentos del OE5}. Los diagramas de contexto y de contenedores
de la alternativa elegida son las Figuras~\ref{fig:c4-contexto} y~\ref{fig:c4-contenedores};
los de las tres alternativas evaluadas están en el informe ATAM y en los diagramas de OE2, y
los de componentes y de despliegue, en el SAD.

\subsection{Tabla de constantes numéricas}
\label{anx:constantes}

La Tabla~\ref{tab:constantes} reúne cada constante numérica del sistema con su valor, su módulo, su origen y su estado, según se explica en la Sección~\ref{sec:constantes}.

{\footnotesize
\setlength{\tabcolsep}{3pt}
\renewcommand{\arraystretch}{1.1}
\begin{longtable}{@{}p{3.1cm}p{2.3cm}p{2.5cm}p{4.7cm}p{1.6cm}@{}}
\caption{Constantes numéricas del sistema, con su origen y su estado.}\label{tab:constantes}\\
\toprule
Constante & Valor & Módulo & Origen & Estado \\
\midrule
\endfirsthead
\toprule
Constante & Valor & Módulo & Origen & Estado \\
\midrule
\endhead
\bottomrule
\endfoot
\multicolumn{5}{@{}l}{\emph{Percepción}}\\
Tope de páginas & 40 & percepción & Experimento: el PDF más largo del corpus tiene 39 páginas & Vigente \\
Temperatura del modelo & 0 & percepción & Decodificación voraz, para reducir la variabilidad entre lecturas & Vigente \\
Tiempo límite por llamada & 40 s & percepción & Una página típica tarda unos 5 s; 40 s indica una llamada colgada & Vigente \\
Reintentos y espera & 1 reintento, 3 s & percepción & Solo ante fallas de transporte, para no ocultar respuestas mal formadas & Vigente \\
\midrule
\multicolumn{5}{@{}l}{\emph{Normalización y resolución}}\\
Dimensión del vector & 1.024 & catálogo & Fija la columna de la base; se pide explícitamente al modelo de \emph{embeddings} & Vigente \\
Lote de \emph{embeddings} & 100 textos & catálogo & Una petición por lote, para evitar cuerpos de petición excesivos & Vigente \\
Umbral de reconocimiento & 0,5 & catálogo & Corte por debajo del cual el renglón no corresponde a ningún concepto; fijado al corregir el caso «chaquiras» & Vigente \\
Aporte del coseno $\gamma$ & 0,15 & catálogo & Ajuste sobre casos de error; pequeño para que el coseno desempate y no decida & En ciclo de calibración \\
Umbral de rescate $\theta$ & 0,86 & catálogo & Ajuste sobre casos de error; separa un sinónimo de un parecido de tema & En ciclo de calibración \\
Piso del rescate $\phi$ & 0,55 & catálogo & Apenas sobre el corte de 0,5, para que un rescate entre a revisión & En ciclo de calibración \\
Umbral de autoaceptación $\tau$ & 0,75 & inteligencia & Requisito PE-04 del SRS; valor inicial anterior a la calibración (0,85 y 0,630 en versiones previas); ver Tabla~\ref{tab:cal-optim} & En ciclo de calibración \\
Tolerancia de medidas & 0,25 y rango 0,85 a 1,35 & catálogo & Ajuste con las auditorías de los kits (Tabla~\ref{tab:kits}) & Vigente \\
Rango de precio en tiendas & \$1.000 a \$400.000 & catálogo & Descarta accesorios sueltos y lotes institucionales & Vigente \\
\midrule
\multicolumn{5}{@{}l}{\emph{Optimización}}\\
Límite de enumeración & 4.096 ($2^{12}$) & inteligencia & Por debajo de ese tamaño del espacio el frente exacto se calcula en milisegundos; por encima se usa NSGA-II & Vigente \\
Población & $\min(200,\max(40,12n))$, par & inteligencia & Escala con las $n$ líneas decidibles; el par evita desperdiciar un individuo & Vigente \\
Generaciones & $\min(120,\max(30,20\log_{10}|\mathcal{E}|))$ & inteligencia & Escala con el logaritmo del espacio $|\mathcal{E}|$ de búsqueda & Vigente \\
Cruce SBX & $p_c=0{,}9$, $\eta_c=15$ & inteligencia & Valores estándar de \citet{deb_simulated_1995} & Vigente \\
Mutación polinomial & $p=1$, $\eta_m=20$ & inteligencia & Valores estándar de \citet{deb_combined_1996} & Vigente \\
Semilla & 0 & inteligencia & Reproducibilidad del frente entre ejecuciones & Vigente \\
Pesos por omisión (costo, marca, calidad, personaje) & 0,35; 0,30; 0,25; 0,10 & inteligencia & Encuesta, pregunta P21 (n=60), normalizada; ver Tabla~\ref{tab:cal-optim} & En ciclo de calibración \\
Sobreprecio para Premium & 10\,\% & inteligencia & Encuesta: 42 de 60 pagan como máximo un 5\,\% más; el 10\,\% marca a quien acepta pagar más & Vigente \\
Opciones por renglón & 3 & catálogo & Pocas opciones diferenciadas facilitan la decisión \citep{iyengar_when_2000} & Vigente \\
\midrule
\multicolumn{5}{@{}l}{\emph{Requisitos de desempeño y operación}}\\
Tiempo de generación del kit & 8 s & SRS (DP-03) & Requisito de desempeño; se contrasta en la Sección~\ref{sec:optimizacion} & Vigente \\
Ventana de revisión & 4 h en temporada & SRS (ML-09) & Requisito: el administrador revisa cada lista nueva antes de publicarla & Vigente \\
Disyuntor & 5 fallas, 30 s & biblioteca común & Patrón de resiliencia; 30 s es el escenario ESC-R2 del ATAM & Vigente \\
Caché de lecturas & 30 días & inteligencia & Evita pagar de nuevo la percepción de un mismo documento & Vigente \\
Caché de vectores & 24 h & catálogo & Evita recalcular el vector de un renglón repetido & Vigente \\
\end{longtable}
}

\subsection{Glosario}
\label{anx:glosario}

\begin{description}[leftmargin=2.5cm,style=nextline]
\item[MAUT] Teoría de utilidad multiatributo: método que resume varios criterios en un solo valor de utilidad mediante pesos que suman uno.
\item[NSGA-II] Algoritmo genético de ordenamiento no dominado que aproxima el frente de Pareto de un problema multiobjetivo \citep{deb_fast_2002}.
\item[Frente de Pareto] Conjunto de soluciones no dominadas, es decir, aquellas en las que ninguna otra mejora un objetivo sin empeorar otro.
\item[Hipervolumen] Volumen del espacio de objetivos que un frente domina respecto de un punto de referencia; mide la calidad de un conjunto de soluciones \citep{zitzler_multiobjective_1999}.
\item[NL2C] Traducción de una preferencia en lenguaje natural a restricciones estructuradas que el optimizador entiende.
\item[$\tau$] Umbral de autoaceptación: puntaje por encima del cual un renglón se da por resuelto sin revisión individual.
\item[$\delta$] Margen mínimo entre el mejor y el segundo candidato que usó una versión anterior del umbral.
\item[$\kappa$] Coeficiente kappa de Cohen: acuerdo entre dos clasificadores corregido por el azar \citep{cohen_coefficient_1960}.
\item[SUS] Escala de usabilidad del sistema, de 0 a 100 \citep{brooke_sus_1996}.
\item[NPS] Índice de recomendación neta, de $-100$ a $100$ \citep{reichheld_one_2003}.
\item[RTLX] Versión sin ponderar del índice de carga de trabajo NASA-TLX, de 0 a 100.
\item[CRISP-DM] Proceso de seis etapas para proyectos de datos \citep{chapman_crisp-dm_2000,wirth_crisp-dm_2000}.
\item[ATAM] Método de evaluación de arquitecturas por compromisos entre atributos de calidad \citep{bass_software_2021}.
\item[ADR] Registro de decisión de arquitectura.
\item[Concepto canónico] Clase de productos equivalentes del catálogo a la que se asigna un renglón de la lista.
\end{description}



\end{document}
